% =====================================================================
% Archivo: CEL_Compilado_Final.tex
% Propósito: Documento Maestro del Sistema CEL (Bloques I a VII) - Generado Automáticamente
% Destinatario: SENER / CNE
% Estilo: cel.cls (Institucional)
% =====================================================================

\documentclass{cel}

% --- Metadatos del Documento ---
\title{Análisis Marco Regulatorio de CEL}
\subtitle{Reingeniería del Sistema de Certificados de Energía Limpia (LSE 2025)}
\author{SENER}
\date{Enero 2026}
\institucion{Secretaría de Energía (SENER)}
\unidad{Comisión Nacional de Energía (CNE)}
\setDocumentoCorto{CEL Enero 2026}
\palabrasclave{CEL, Soberanía, CNE, Justicia Energética, Transparencia}
\version{Final 1.0}

% --- Metadatos PDF/UA ---
\hypersetup{
  pdftitle={Análisis Integral del Sistema CEL - Reingeniería 2025},
  pdfauthor={SENER},
  pdfsubject={Modernización del S-CEL bajo la Nueva Ley del Sector Eléctrico},
  pdfkeywords={CEL, CNE, LSE 2025, Transición Energética, Soberanía},
  pdfversion={1}
}

\begin{document}
\pagenumbering{roman}

% =====================================================================
% PORTADA PRINCIPAL
% =====================================================================
\portadaseccion{}{Análisis Marco Regulatorio de CEL}{Enero 2026}

% =====================================================================
% ÍNDICE
% =====================================================================
\tableofcontents
\newpage


\portadaseccion{BLOQUE I}{Proceso de inscripción al Sistema de Certificados de Energía Limpia (S-CEL)}{}
\clearpage
\pagenumbering{arabic}
\setcounter{page}{1}
\phantomsection
\addcontentsline{toc}{section}{Proceso de inscripción al Sistema de Certificados de Energía Limpia (S-CEL)}
\section*{Proceso de inscripción al Sistema de Certificados de Energía Limpia (S-CEL)}
\nopagebreak

\textbf{Alcance del punto:} Reingeniería del proceso de inscripción al S-CEL bajo la soberanía de datos de la CNE y el cumplimiento integral de impacto social (SENER).

\textbf{Instrumentos jurídicos revisados:} Ley del Sector Eléctrico (LSE 2025), Reglamento de la LSE (RLSE 2025), Ley de Planeación y Transición Energética (LPTE).



\subsection{A. Marco Normativo Actualizado (LSE 2025)}
\nopagebreak

A diferencia del marco anterior (LIE/CRE), el nuevo sistema S-CEL 2.0 se fundamenta en la \textbf{Soberanía de Datos} y la integración vinculante con la planeación de la Secretaría de Energía (SENER).


\vspace{0.3cm}
\Needspace{15\baselineskip}
\begin{tabladoradoLargo}
\renewcommand{\tabularxcolumn}[1]{p{#1}}
\begin{xltabular}{\textwidth}{|X|X|X|X|}
\caption{Matriz de Datos} \\
\toprule
\rowcolor{gobmxDorado} \encabezadotablaDorada{Actor / Fuente} & \encabezadotablaDorada{Instrumento legal} & \encabezadotablaDorada{Artículo / Referencia} & \encabezadotablaDorada{Contexto Normativo Actualizado} \\
\midrule
\endhead
\textbf{Generadoras y Cargas} & LSE 2025 & Art. 11 Fracc. XVI & La CNE otorga los CEL y administra el registro único de certificados. \\ \hline
\textbf{Soberanía de Datos} & LSE 2025 & Art. 149 & Registro de Certificados: Matriculación unitaria e historial de propiedad inviolable (CNE). \\ \hline
\textbf{Inscripción General} & Reglamento LSE & Art. 182 & Regulación del otorgamiento, liquidación y transacciones. Sistema electrónico centralizado. \\ \hline
\textbf{Participantes Obligados} & LSE 2025 & Atribuciones CNE & Inscripción de Usuarios Calificados y Comercializadoras bajo criterios de fomento y eficiencia. \\ \hline
\textbf{Exención Social (GLD)} & Reglamento LSE & Art. 203 Fracc. I & Generación \textless{} 0.7 MW exenta de presentar Manifestación de Impacto Social (MIS). \\ \hline
\textbf{Validación Social (\textgreater{} 0.7 MW)} & Reglamento LSE & Art. 205 & La autorización de la MIS (SENER) es requisito previo obligatorio a la construcción e inscripción. \\ \hline
\textbf{Permisionarios Legados} & Reglamento LSE & Transitorio 19 & Migración expedita obligatoria para figuras LSPEE como requisito para acceder al S-CEL. \\ \hline
\bottomrule
\end{xltabular}
\end{tabladoradoLargo}


\vspace{0.4cm}
\Needspace{25\baselineskip}
\subsection{B. Matriz de Auditoría y Reingeniería (LIE vs. LSE 2025)}
\nopagebreak

\vspace{0.3cm}
\Needspace{15\baselineskip}
\begin{tabladoradoLargo}
\renewcommand{\tabularxcolumn}[1]{p{#1}}
\begin{xltabular}{\textwidth}{|X|X|X|X|}
\caption{Matriz de Datos} \\
\toprule
\rowcolor{gobmxDorado} \encabezadotablaDorada{Elemento Auditado} & \encabezadotablaDorada{Marco Anterior (LIE/CRE)} & \encabezadotablaDorada{Nuevo Marco (LSE 2025)} & \encabezadotablaDorada{Riesgo Resuelto} \\
\midrule
\endhead
\textbf{Autoridad de Registro} & CRE (Órgano Regulador Desconectado) & CNE (Autoridad con Soberanía de Datos y Planeación) & Fragmentación de la información y falta de alineación con la política energética nacional. \\ \hline
\textbf{Validación Social} & Post-facto o inexistente para inscripción & \textbf{Pre-requisito Mandatorio} (Art. 205 Reglamento) & Conflictos territoriales y nulidad de proyectos por falta de licencia social. \\ \hline
\textbf{Tratamiento de Legados} & Mantenimiento de esquemas de \textquotedbl{}Autoabasto\textquotedbl{} & \textbf{Migración Obligatoria} (Transitorio 19) & Dispersión normativa y competencia desleal por figuras jurídicas extintas. \\ \hline
\textbf{Costo de Inscripción} & Pago de derechos gravosos para participantes & \textbf{Gratuidad Estratégica} (Bases de Registro CNE) & Barreras de entrada para la demanda controlable y la comercialización pública. \\ \hline
\textbf{Naturaleza del Acto} & Trámite Administrativo sujeto a discreción & \textbf{Acto de Publicidad Registral Automatizado} & Corrupción y lentitud burocrática en el otorgamiento de títulos. \\ \hline
\bottomrule
\end{xltabular}
\end{tabladoradoLargo}


\subsection{C. Desarrollo Analítico}
\nopagebreak

\subsubsection{1. El Paradigma de la Soberanía de Datos (Art. 149 LSE)}
\nopagebreak

Bajo la LSE 2025, el \textbf{S-CEL 2.0 \textbf{ deja de ser un simple repositorio documental para convertirse en una infraestructura de }Soberanía de Datos}.


\begin{itemize}
\item \textbf{Matriculación Inviolable:} Cada CEL posee un identificador único que vincula la central, la tecnología y el timestamp de generación en tiempo real.
\item \textbf{Fe Pública Digital:} El Registro de la CNE actúa como fedatario público de los atributos limpios, eliminando la necesidad de validaciones de terceros (Unidades Acreditadas) para el proceso de inscripción básico, migrando hacia auditorías de sistemas.

\end{itemize}
\subsubsection{2. Integración del Impacto Social (Art. 201 y 205 RLSE)}
\nopagebreak

Un cambio crítico es la subordinación de la inscripción técnica a la validación social:


\begin{itemize}
\item \textbf{Filtro Social: \textbf{ SENER debe emitir la autorización definitiva de la }Manifestación de Impacto Social (MIS)} antes de que la CNE proceda con la inscripción definitiva de cualquier central mayor a 0.7 MW.
\item \textbf{Justicia Energética:} Este mecanismo asegura que los CELs emitidos provengan de proyectos que respetan los derechos de las comunidades y contribuyen a la justicia energética (Art. 3 Fracc. XXVIII LSE).

\end{itemize}
\subsubsection{3. Procedimiento de Inscripción Optimizado}
\nopagebreak

\paragraph{Ruta A: Generación Exenta (GLD \textless{} 0.7 MW)}
\nopagebreak

\begin{enumerate}
\item \textbf{Simplificación Total:} No requiere MIS (Art. 203 RLSE).
\item \textbf{Inscripción Automática:} Basada en el Contrato de Interconexión y verificación de la Suministradora que la representa (Art. 188 RLSE).

\end{enumerate}
\paragraph{Ruta B: Centrales Permisionadas y Legadas}
\nopagebreak

\begin{enumerate}
\item \textbf{Validación de Migración:} Para legados, presentar Oficio de Migración (Transitorio 19 RLSE).
\item \textbf{Validación MIS:} Consulta síncrona al padrón de SENER.
\item \textbf{Inscripción Provisional:} Figura de \textquotedbl{}Resguardo de Atributos\textquotedbl{} mientras se completa la verificación física de interconexión.


\end{enumerate}
\subsection{D. Propuestas de Ajuste Normativo}
\nopagebreak

Para implementar operativamente estos cambios, se proponen las siguientes Bases de Registro ante la CNE:


\begin{enumerate}
\item \textbf{Billetera S-CEL con \textquotedbl{}Escrow\textquotedbl{}:} Implementar un mecanismo donde los CELs generados en fase de pruebas se retengan administrativamente hasta que la resolución de inscripción definitiva sea firme.
\item \textbf{Interoperabilidad SENER-CNE:} Obligatoriedad de conexión entre el sistema de Gestión Social de SENER y el Registro S-CEL para evitar \textquotedbl{}skip-steps\textquotedbl{} en la validación de impacto social.
\item \textbf{Certificación de Datos vs. Documentos:} Migrar el requisito de \textquotedbl{}copias certificadas\textquotedbl{} de documentos técnicos por \textquotedbl{}firmas electrónicas avanzadas\textquotedbl{} de los responsables de telemetría de CENACE.


\end{enumerate}
\vspace{0.4cm}
\Needspace{25\baselineskip}
\subsection{E. Estrategia de Transitorios y Migración}
\nopagebreak

\vspace{0.3cm}
\Needspace{15\baselineskip}
\begin{tabladoradoLargo}
\renewcommand{\tabularxcolumn}[1]{p{#1}}
\begin{xltabular}{\textwidth}{|X|X|X|X|}
\caption{Matriz de Datos} \\
\toprule
\rowcolor{gobmxDorado} \encabezadotablaDorada{Sujeto} & \encabezadotablaDorada{Acción Requerida} & \encabezadotablaDorada{Plazo Estimado (Días)} & \encabezadotablaDorada{Fundamento Legal} \\
\midrule
\endhead
\textbf{Autoabastos Vigentes} & Solicitar Migración Expedita a Ley LSE 2025. & \textless{} 120 días & Transitorio 19 RLSE \\ \hline
\textbf{Generadoras en Trámite (CRE)} & Ratificar solicitud ante CNE bajo nuevo marco. & Inmediato & Transitorio 24 RLSE \\ \hline
\textbf{Usuarios Calificados Inscritos} & Actualizar datos en el nuevo Registro S-CEL 2.0 (CNE). & \textless{} 180 días & Art. 156 RLSE \\ \hline
\textbf{Proyectos con MIS pendiente} & Concluir trámite en SENER conforme a RLSE. & Según proceso SENER & Art. 201/213 RLSE \\ \hline
\bottomrule
\end{xltabular}
\end{tabladoradoLargo}

\phantomsection
\addcontentsline{toc}{section}{Registro de Participantes y Administración de Cuentas}
\section*{Registro de Participantes y Administración de Cuentas}
\nopagebreak

\subsection{1. Auditoría de la \textquotedbl{}Foto Actual\textquotedbl{} (Marco Anterior: LIE/CRE)}
\nopagebreak

\subsubsection{Instrumentos Citados}
\nopagebreak

\begin{itemize}
\item \textbf{RES/174/2016}: Estableció el Sistema de Gestión de Certificados y Cumplimiento de Obligaciones de Energías Limpias (S-CEL).
\item \textbf{A/067/2017}: Modificó las disposiciones para incluir conceptos como \textquotedbl{}CEL en proceso\textquotedbl{} y \textquotedbl{}Unidades Acreditadas\textquotedbl{} para la certificación técnica.
\item \textbf{Ley de la Industria Eléctrica (LIE) Art. 123}: Definía a los \textquotedbl{}Participantes Obligados\textquotedbl{} bajo un esquema de mercado fragmentado.

\end{itemize}
\subsubsection{Diagnóstico de Fallas Operativas}
\nopagebreak

\begin{enumerate}
\item \textbf{Burocracia Excesiva (Unidades Acreditadas)}: La obligatoriedad de dictámenes técnicos externos (Disp. 11 de RES/174/2016) generaba cuellos de botella y costos transaccionales altos para los generadores limpios, retrasando el otorgamiento de activos.
\item \textbf{Desconexión de Datos}: El S-CEL operaba de forma semiautomática con reportes del CENACE y Distribuidores que a menudo requerían aclaraciones manuales (Disp. 30), lo que generaba incertidumbre en la liquidación de obligaciones.
\item \textbf{Complejidad en Transacciones}: La introducción de estados como \textquotedbl{}CEL en proceso\textquotedbl{} (A/067/2017) buscaba evitar la doble venta pero complicaba la fluidez del mercado secundario al congelar activos por falta de interoperabilidad financiera inmediata.
\item \textbf{Vacío de Soberanía}: Los datos de generación y consumo residían en infraestructuras que no priorizaban la soberanía tecnológica ni la transparencia directa hacia la autoridad reguladora centralizada.


\end{enumerate}
\subsection{2. Reingeniería Legal (Nuevo Marco LSE 2025 / CNE)}
\nopagebreak

La transición al marco de \textbf{Soberanía Energética \textbf{ exige una migración del S-CEL al }S-RCEL (Sistema Registro y Control de Energía Limpia)}, bajo control directo de la CNE.


\vspace{0.4cm}
\Needspace{25\baselineskip}
\subsubsection{Sustitución de Referencias}
\nopagebreak

\vspace{0.3cm}
\Needspace{15\baselineskip}
\begin{tabladoradoLargo}
\renewcommand{\tabularxcolumn}[1]{p{#1}}
\begin{xltabular}{\textwidth}{|X|X|X|X|}
\caption{Matriz de Datos} \\
\toprule
\rowcolor{gobmxDorado} \encabezadotablaDorada{Concepto} & \encabezadotablaDorada{Marco Anterior (LIE/CRE)} & \encabezadotablaDorada{Nuevo Marco (LSE 2025 / RLSE)} & \encabezadotablaDorada{Justificación de Reingeniería} \\
\midrule
\endhead
\textbf{Autoridad de Registro} & CRE (Comisión Reguladora) & \textbf{CNE (Comisión Nacional de Energía)} & Centralización de facultades para garantizar planeación vinculante (Art. 182 RLSE). \\ \hline
\textbf{Obligatoriedad de Sistema} & Registro de CEL (Art. 128 LIE) & \textbf{S-RCEL Obligatorio (Art. 184 RLSE)} & Transacciones solo válidas a través del sistema electrónico habilitado para evitar mercados paralelos opacos. \\ \hline
\textbf{Sujetos Obligados} & Suministradores / Usuarios Calificados & \textbf{Centros de Carga (Art. 186 RLSE)} & Se redefine la obligación proporcional al consumo directo, simplificando la fiscalización. \\ \hline
\textbf{Principios Operativos} & Eficiencia de Mercado & \textbf{Justicia Energética y Simplificación (Art. 6 RLSE)} & Actuación administrativa bajo celeridad, eficacia y transparencia con enfoque social. \\ \hline
\bottomrule
\end{xltabular}
\end{tabladoradoLargo}

\subsubsection{Adiciones Clave: Soberanía e Impacto Social}
\nopagebreak

\begin{itemize}
\item \textbf{Soberanía de Datos}: En cumplimiento con los principios de seguridad energética, el S-RCEL debe operar sobre infraestructura nacional, garantizando la inmutabilidad de los registros de generación limpia.
\item \textbf{Impacto Social (Art. 201+ RLSE) \textbf{: El registro de nuevos generadores está condicionado a la validación de la }Manifestación de Impacto Social}, asegurando que la emisión de CELs responda a proyectos con licencia social.


\end{itemize}
\subsection{3. Propuesta Normativa (Nuevas DACG para el S-RCEL)}
\nopagebreak

Se proponen los siguientes artículos técnicos para integrar en las nuevas Disposiciones Administrativas de Carácter General:


\subsubsection{Interoperabilidad y Automatización}
\nopagebreak

\vspace{0.2cm}
\Needspace{8\baselineskip}
\begin{glowBox}[gobmxDorado]{Nota Informativa}
\textbf{Art. X (Automatización)}: La CNE implementará un motor de conciliación automática con la Empresa Pública del Estado para validar telemetría de generación en tiempo real, eliminando la necesidad de validaciones manuales mensuales.
\end{glowBox}
\textgreater{}

\vspace{0.2cm}
\Needspace{8\baselineskip}
\begin{glowBox}[gobmxDorado]{Nota Informativa}
\textbf{Art. Y (Interoperabilidad)}: El S-RCEL contará con una API segura para la interoperabilidad con sistemas fiscales, permitiendo que la liquidación de CEL sea un proceso de \textquotedbl{}un solo click\textquotedbl{} vinculado al cumplimiento de obligaciones del ejercicio.
\end{glowBox}

\subsubsection{Seguridad y Continuity}
\nopagebreak

\begin{itemize}
\item \textbf{Cifrado y Trazabilidad}: Cada folio de CEL será generado mediante algoritmos de cifrado que garanticen su unicidad y origen desde la Central Eléctrica.
\item \textbf{Auditoría de Algoritmos}: La CNE auditará semestralmente los algoritmos de asignación de CEL para garantizar el cumplimiento de las metas de transición (Art. 181 RLSE).


\end{itemize}
\subsection{4. Matriz de Transición y Variables Críticas}
\nopagebreak

\vspace{0.4cm}
\Needspace{25\baselineskip}
\subsubsection{Tabla Comparativa: Antes vs. Ahora}
\nopagebreak

\vspace{0.3cm}
\Needspace{15\baselineskip}
\begin{tablaParadigmaV2}{Análisis Comparativo}{Modelo Anterior (S-CEL)}{Modelo Nuevo (S-RCEL)}
\tpRow{\textbf{Validación técnica}}{Dictamen de Unidad Acreditada (Costo Privado)}{\textbf{Validación técnica}}{Validación Sistémica CNE/CENACE (Servicio Público)}
\tpRow{\textbf{Frecuencia de Conciliación}}{Mensual (Reporte Ex-post)}{\textbf{Frecuencia de Conciliación}}{Tiempo Real (Telemetría de Red)}
\tpRow{\textbf{Transacciones}}{Mercado Secundario Opaco}{\textbf{Transacciones}}{Plataforma de Compraventa Centralizada (API CNE)}
\tpRow{\textbf{Vigencia de Activos}}{20 años}{\textbf{Vigencia de Activos}}{Permanente hasta liquidación (Transitorio Sexto RLSE)}
\end{tablaParadigmaV2}

\subsubsection{Variables Críticas de Registro}
\nopagebreak

\begin{enumerate}
\item \textbf{Código ID-CNE}: Identificador único de la Central Eléctrica vinculada a la planeación.
\item \textbf{Factor de Justicia Energética}: Variable que pondera el beneficio social del proyecto.
\item \textbf{Hash de Generación}: Código criptográfico que vincula el MWh físico con el activo digital.


\end{enumerate}
\subsection{5. Régimen Transitorio de Migración}
\nopagebreak

\textbf{Transitorio Primero (Validez de Activos) \textbf{: Conforme al }Transitorio Sexto del RLSE 2025}, los CEL emitidos bajo `RES/174/2016` conservan su validez. El S-RCEL habilitará un módulo de migración para que los poseedores actuales carguen sus folios legados.


\textbf{Transitorio Segundo (Plazos)}: La CNE publicará el cronograma de desconexión del S-CEL (CRE) y activación del S-RCEL en un plazo no mayor a 90 días, asegurando la continuidad operativa de las liquidaciones anuales 2025.


\textbf{Transitorio Tercero (Sanciones) \textbf{: Los incumplimientos en el marco anterior serán subsanables mediante el mecanismo de }Reparación del Daño} previsto en el Art. 302 del RLSE, previo a la migración al nuevo sistema.


\phantomsection
\addcontentsline{toc}{section}{Medición, fuentes de información y datos vinculantes}
\section*{Medición, fuentes de información y datos vinculantes}
\nopagebreak

\textbf{Alcance del punto:} Medición, fuentes de información y datos vinculantes que alimentan el S-CEL

\textbf{Instrumentos jurídicos revisados:} LSE 2025, RLSE 2025 (Art. 180-191), DACG S-CEL (RES/174/2016 y A/067/2017 - marco anterior), RES/142/2017 (Generación Distribuida), Disposiciones de Autoconsumo (2025), Bases del Mercado Eléctrico, Lineamientos CEL



\vspace{0.4cm}
\Needspace{25\baselineskip}
\subsection{A. Fuentes de Información del S-CEL}
\nopagebreak

\vspace{0.3cm}
\Needspace{15\baselineskip}
\begin{tabladoradoLargo}
\renewcommand{\tabularxcolumn}[1]{p{#1}}
\begin{xltabular}{\textwidth}{|X|X|X|X|}
\caption{Matriz de Datos} \\
\toprule
\rowcolor{gobmxDorado} \encabezadotablaDorada{Actor/Fuente} & \encabezadotablaDorada{Instrumento legal} & \encabezadotablaDorada{Artículo / Numeral} & \encabezadotablaDorada{Cita textual} \\
\midrule
\endhead
\textbf{Autoridad Reguladora (CNE)} & LSE 2025 / RLSE 2025 & Art. 182, 190 & \textquotedbl{}La CNE debe emitir las disposiciones administrativas de carácter general que regulen el otorgamiento, liquidación, cancelación voluntaria y transacciones relacionadas con los Certificados de Energías Limpias\textquotedbl{} (Art. 182 RLSE 2025). Marco anterior: CRE (RES/174/2016) \\ \hline
\textbf{CENACE} & DACG S-CEL (RES/174/2016) & Disp. 26 & \textquotedbl{}el Cenace, los Transportistas, los Distribuidores, los Generadores y Generadores Exentos que producen energía por el abasto aislado, informarán a la Comisión mediante el Sistema, la energía eléctrica generada en el mes calendario anterior\textquotedbl{} (Marco anterior - vigente hasta transición) \\ \hline
\textbf{Transportistas} & DACG S-CEL (RES/174/2016) & Disp. 26 & \textquotedbl{}informarán a la Comisión mediante el Sistema, la energía eléctrica generada en el mes calendario anterior por cada Central Eléctrica Limpia\textquotedbl{} \\ \hline
\textbf{Distribuidores} & DACG S-CEL (A/067/2017) & Disp. 32.B & \textquotedbl{}para otorgar los CEL a los Suministradores que representen Generación Limpia Distribuida se utilizará los datos de medición reportados por el Distribuidor\textquotedbl{} \\ \hline
\textbf{Generadores (Abasto Aislado)} & DACG S-CEL (RES/174/2016) & Disp. 26 & \textquotedbl{}los Generadores y Generadores Exentos que producen energía por el abasto aislado, informarán a la Comisión\textquotedbl{} \\ \hline
\textbf{Medición Vinculante CENACE} & DACG S-CEL (A/067/2017) & Disp. 32.B & \textquotedbl{}Los datos de medición de generación utilizados en las liquidaciones y reliquidaciones reportados por el Cenace al S-CEL serán la base para el otorgamiento de los CEL\textquotedbl{} \\ \hline
\textbf{Participantes Obligados (Abasto Aislado)} & DACG S-CEL (A/067/2017) & Disp. 30.B & \textquotedbl{}Para el abasto aislado las Centrales Eléctricas deberán reportar al S-CEL la energía generada neta y consumida mensualmente\textquotedbl{} \\ \hline
\bottomrule
\end{xltabular}
\end{tabladoradoLargo}


\vspace{0.4cm}
\Needspace{25\baselineskip}
\subsection{B. Matriz de Validación Jurídica}
\nopagebreak

\vspace{0.3cm}
\Needspace{15\baselineskip}
\begin{tabladoradoLargo}
\renewcommand{\tabularxcolumn}[1]{p{#1}}
\begin{xltabular}{\textwidth}{|X|X|X|X|}
\caption{Matriz de Datos} \\
\toprule
\rowcolor{gobmxDorado} \encabezadotablaDorada{Hallazgo o limitación actual} & \encabezadotablaDorada{Fundamento jurídico (APA + cita textual)} & \encabezadotablaDorada{Riesgo (jurídico u operativo)} & \encabezadotablaDorada{Ajuste propuesto} \\
\midrule
\endhead
\textbf{Fundamento de medición bajo LSE 2025} & CNE. (2025). \textit{Reglamento de la Ley del Sector Eléctrico}. Artículo 190. \textquotedbl{}La CNE debe emitir las disposiciones administrativas de carácter general para establecer los términos para acreditar a las unidades que certifican a las Centrales Eléctricas que generan Energía Limpia, la medición de variables requeridas para determinar el porcentaje de energía libre de combustible generada por las Centrales Eléctricas susceptibles de recibir Certificados de Energía Limpia y determinar el porcentaje de Energías Limpias consumida por los Centros de Carga\textquotedbl{} & Regulatorio: Nuevo marco de medición bajo CNE & Implementar metodología de medición conforme Art. 190 RLSE 2025 \\ \hline
\textbf{Dependencia de reportes manuales para abasto aislado} & Comisión Reguladora de Energía. (2017). \textit{DACG del Sistema de CEL modificadas} . Disposición 30.B. \textquotedbl{}Para el abasto aislado las Centrales Eléctricas deberán reportar al S-CEL la energía generada neta y consumida mensualmente\textquotedbl{} & Operativo: Errores de captura manual y retrasos en información & Implementar sistemas de medición automática para abasto aislado \\ \hline
\textbf{Falta de validación cruzada automática entre fuentes} & Comisión Reguladora de Energía. (2016). \textit{DACG del Sistema de CEL} . Disposición 30. \textquotedbl{}será empleada para verificar la validez y consistencia de la información que se consideró para el otorgamiento de CEL\textquotedbl{} & Técnico: Inconsistencias no detectadas oportunamente entre fuentes & Desarrollar algoritmos de validación cruzada en tiempo real \\ \hline
\textbf{Cobertura parcial de generación distribuida} & DACG S-CEL (A/067/2017), Disp. 32.B & Dependencia exclusiva de reportes del Distribuidor para GLD & Operativo: Subregistro de generación limpia distribuida \\ \hline
\textbf{Falta de integración con sistemas de medición distribuida \textless{}0.5 MW} & RES/142/2017, Capítulo VI: \textquotedbl{}el Distribuidor, elaborará y actualizará una base de datos y reportes\textquotedbl{} & Operativo: Ausencia de datos de centrales pequeñas con potencial significativo agregado & Implementar protocolos de RES/142/2017 para medición distribuida \\ \hline
\textbf{Ausencia de metodologías de contraprestación en medición} & RES/142/2017, Anexo I: \textquotedbl{}Medición neta de energía: metodología de contraprestación que considera los flujos de energía eléctrica recibidos y entregados\textquotedbl{} & Técnico: Medición inadecuada para esquemas Net Metering y Net Billing & Integrar metodologías de medición bidireccional según RES/142/2017 \\ \hline
\textbf{Falta de integración con figura de autoconsumo} & Disposiciones de Autoconsumo (2025): \textquotedbl{}regular la figura de Autoconsumo de Energía Eléctrica\textquotedbl{} & Operativo: Medición no adaptada para modalidades de autoconsumo & Implementar sistemas de medición específicos para autoconsumo según disposiciones 2025 \\ \hline
\textbf{Ausencia de medición para centrales 0.7-20 MW en autoconsumo} & Disposiciones de Autoconsumo (2025): rango específico 0.7-20 MW & Técnico: Vacío en medición para rango intermedio de autoconsumo & Desarrollar protocolos de medición para autoconsumo interconectado 0.7-20 MW \\ \hline
\textbf{Ausencia de integración del abasto aislado no interconectado} & Comisión Reguladora de Energía. (2017). \textit{DACG del Sistema de CEL modificadas} . Disposición 30. \textquotedbl{}Para el cálculo de las Obligaciones en materia de Energías Limpias para los Centros de Carga que se suministren por abasto aislado no interconectado al SEN, se utilizarán los datos de medición que el Participante del S-CEL entregue\textquotedbl{} & Regulatorio: Falta de control sobre datos no verificables & Establecer mecanismos de verificación independiente \\ \hline
\textbf{Falta de criterios específicos para abasto aislado con necesidades propias} & Acuerdo A/049/2017: \textquotedbl{}criterio de interpretación del concepto 'necesidades propias', establecido en el artículo 22 de la Ley de la Industria Eléctrica\textquotedbl{} & Técnico: Medición no diferenciada para abasto aislado con necesidades propias & Implementar criterios específicos de medición para abasto aislado según A/049/2017 \\ \hline
\textbf{Uso de valores estimados sin criterios claros de calidad} & Comisión Reguladora de Energía. (2017). \textit{DACG del Sistema de CEL modificadas} . Disposición 31. \textquotedbl{}En aquellos casos en que no sea posible determinar el número de CEL... la Comisión podrá usar temporalmente valores estimados\textquotedbl{} & Técnico: Imprecisión en otorgamiento de CEL por estimaciones & Definir criterios de calidad y límites para uso de estimaciones \\ \hline
\textbf{Plazos de reporte heterogéneos entre actores} & DACG S-CEL (RES/174/2016), Disp. 26 & \textquotedbl{}en los diez primeros días hábiles de cada mes\textquotedbl{} vs reportes trimestrales & Operativo: Desincronización de información \\ \hline
\bottomrule
\end{xltabular}
\end{tabladoradoLargo}


\subsection{C. Desarrollo Analítico}
\nopagebreak

\subsubsection{Diagnóstico de la Situación Actual}
\nopagebreak

El sistema de medición y fuentes de información del S-CEL presenta las siguientes características:


\paragraph{Fuentes de Datos Actuales}
\nopagebreak

\begin{itemize}
\item Datos de liquidaciones y reliquidaciones del MEM
\item Base para otorgamiento de CEL a centrales interconectadas
\item Información más confiable y automatizada

\end{itemize}
\begin{enumerate}
\item \textbf{Transportistas y Distribuidores} :

\end{enumerate}
\begin{itemize}
\item Reporte mensual de generación por central
\item Información complementaria a CENACE
\item Datos de generación limpia distribuida (solo Distribuidores)

\end{itemize}
\begin{enumerate}
\item \textbf{Generadores de Abasto Aislado} :

\end{enumerate}
\begin{itemize}
\item Reportes manuales mensuales
\item Mayor riesgo de errores e inconsistencias
\item Diferenciación entre interconectado y no interconectado al SEN

\end{itemize}
\begin{enumerate}
\item \textbf{Participantes Obligados (Abasto Aislado)} :

\end{enumerate}
\begin{itemize}
\item Reportes de consumo para cálculo de obligaciones
\item Datos no verificados independientemente

\end{itemize}
\paragraph{Problemáticas Identificadas}
\nopagebreak

\begin{enumerate}
\item \textbf{Heterogeneidad de Fuentes} :

\end{enumerate}
\begin{itemize}
\item Diferentes niveles de automatización y confiabilidad
\item Plazos de reporte no homogéneos
\item Calidad variable de la información

\end{itemize}
\begin{enumerate}
\item \textbf{Limitaciones de Cobertura} :

\end{enumerate}
\begin{itemize}
\item Generación distribuida dependiente únicamente del Distribuidor
\item Abasto aislado no interconectado sin verificación independiente
\item Falta de integración de sistemas de medición inteligente

\end{itemize}
\begin{enumerate}
\item \textbf{Deficiencias en Validación} :

\end{enumerate}
\begin{itemize}
\item Validación cruzada manual y ex-post
\item Uso frecuente de valores estimados
\item Ausencia de alertas automáticas de inconsistencias

\end{itemize}
\begin{enumerate}
\item \textbf{Riesgos Operativos} :

\end{enumerate}
\begin{itemize}
\item Dependencia de reportes manuales propensos a errores
\item Retrasos en detección de inconsistencias
\item Dificultad para auditar información de abasto aislado

\end{itemize}
\subsubsection{Estado Objetivo}
\nopagebreak

Implementar un sistema integral de medición que garantice:


\paragraph{Principios Rectores}
\nopagebreak

\begin{itemize}
\item \textbf{Universalidad} : Cobertura de toda la generación limpia nacional
\item \textbf{Soberanía de Datos} : Control nacional sobre datos de generación y medición de energía limpia (LSE 2025)
\item \textbf{Automatización} : Captura automática de datos desde sistemas de medición
\item \textbf{Confiabilidad} : Validación cruzada en tiempo real entre múltiples fuentes
\item \textbf{Trazabilidad} : Seguimiento completo desde generación hasta registro
\item \textbf{Verificabilidad} : Mecanismos independientes de verificación conforme Art. 190 RLSE 2025

\end{itemize}
\paragraph{Características del Sistema Objetivo}
\nopagebreak

\begin{itemize}
\item \textbf{Medición Universal} : Cobertura de gran escala, distribuida y aislada
\item \textbf{Integración Completa} : Conexión con todos los sistemas de medición
\item \textbf{Validación Automática} : Algoritmos de verificación cruzada en tiempo real
\item \textbf{Calidad de Datos} : Estándares mínimos de precisión y completitud
\item \textbf{Auditoría Continua} : Monitoreo permanente de calidad de información

\end{itemize}
\vspace{0.4cm}
\Needspace{25\baselineskip}
\subsubsection{Tabla Comparativa: Modelo Actual vs Modelo Objetivo}
\nopagebreak

\vspace{0.3cm}
\Needspace{15\baselineskip}
\begin{tablaParadigmaV2}{Análisis Comparativo}{Modelo Actual}{Modelo Objetivo}
\tpRow{\textbf{Cobertura}}{Limitada a gran escala interconectada}{\textbf{Cobertura}}{Universal: gran escala + distribuida + aislada}
\tpRow{\textbf{Fuentes de Datos}}{Principalmente CENACE + reportes manuales}{\textbf{Fuentes de Datos}}{CENACE + Distribuidores + IoT + Telemetría}
\tpRow{\textbf{Validación}}{Manual, ex-post, tolerancia 2\%}{\textbf{Validación}}{Automática, tiempo real, alertas inmediatas}
\tpRow{\textbf{Abasto Aislado}}{Reportes manuales sin verificación}{\textbf{Abasto Aislado}}{Medición automática con verificación independiente}
\tpRow{\textbf{GLD}}{Solo reportes del Distribuidor}{\textbf{GLD}}{Medición inteligente + múltiples fuentes}
\tpRow{\textbf{Estimaciones}}{Sin criterios claros de calidad}{\textbf{Estimaciones}}{Límites estrictos y criterios de calidad definidos}
\tpRow{\textbf{Plazos}}{Heterogéneos (mensual/trimestral)}{\textbf{Plazos}}{Homogéneos, tiempo real cuando sea posible}
\tpRow{\textbf{Trazabilidad}}{Limitada por fuente}{\textbf{Trazabilidad}}{Completa, desde medidor hasta CEL}
\end{tablaParadigmaV2}

\subsubsection{Arquitectura del Sistema (Descripción Conceptual)}
\nopagebreak

\paragraph{Componentes Principales}
\nopagebreak

\begin{enumerate}
\item \textbf{Hub Central de Datos}
\end{enumerate}
\begin{itemize}
\item Recepción y consolidación de todas las fuentes
\item Normalización y estandarización de formatos
\item Almacenamiento con trazabilidad completa
\end{itemize}
\begin{enumerate}
\item \textbf{Red de Sensores IoT}
\end{enumerate}
\begin{itemize}
\item Medición inteligente para generación distribuida
\item Telemetría para abasto aislado
\item Comunicación en tiempo real vía redes celulares/satelitales
\end{itemize}
\begin{enumerate}
\item \textbf{Motor de Validación Cruzada}
\end{enumerate}
\begin{itemize}
\item Algoritmos de detección de inconsistencias
\item Validación automática entre múltiples fuentes
\item Alertas inmediatas de anomalías
\end{itemize}
\begin{enumerate}
\item \textbf{Sistema de Calidad de Datos}
\end{enumerate}
\begin{itemize}
\item Métricas de completitud, precisión y oportunidad
\item Clasificación automática de calidad por fuente
\item Reportes de calidad para mejora continua
\end{itemize}
\begin{enumerate}
\item \textbf{Plataforma de Auditoría}
\end{enumerate}
\begin{itemize}
\item Verificación independiente de datos críticos
\item Muestreo estadístico para validación
\item Trazabilidad completa de modificaciones

\end{itemize}
\subsubsection{Reingeniería de Procesos (Pasos Operativos)}
\nopagebreak

\paragraph{Proceso Optimizado de Captura y Validación}
\nopagebreak

\textbf{Captura Automática Universal}


\begin{enumerate}
\item * CENACE: Datos de liquidaciones en tiempo real
\end{enumerate}
\begin{itemize}
\item Distribuidores: Medición inteligente de GLD
\item Abasto Aislado: Telemetría automática
\item Transportistas: Validación cruzada de datos CENACE
\end{itemize}
\begin{enumerate}
\item \textbf{Normalización y Consolidación}
\end{enumerate}
\begin{itemize}
\item Conversión automática a formatos estándar
\item Asignación de timestamps precisos
\item Clasificación por tipo de fuente y tecnología
\item Almacenamiento con metadatos completos
\end{itemize}
\begin{enumerate}
\item \textbf{Validación Cruzada Automática}
\end{enumerate}
\begin{itemize}
\item Comparación entre fuentes múltiples
\item Detección de outliers estadísticos
\item Verificación de consistencia temporal
\item Generación automática de alertas
\end{itemize}
\begin{enumerate}
\item \textbf{Gestión de Calidad}
\end{enumerate}
\begin{itemize}
\item Evaluación automática de completitud
\item Cálculo de métricas de precisión
\item Clasificación de confiabilidad por fuente
\item Reportes de calidad para mejora
\end{itemize}
\begin{enumerate}
\item \textbf{Auditoría y Verificación}
\end{enumerate}
\begin{itemize}
\item Muestreo aleatorio para verificación independiente
\item Auditorías programadas de fuentes críticas
\item Validación de estimaciones contra datos reales
\item Documentación completa de ajustes

\end{itemize}
\subsubsection{Beneficios Esperados}
\nopagebreak

\paragraph{Beneficios Operativos}
\nopagebreak

\begin{itemize}
\item \textbf{Cobertura Universal} : Captura del 100\% de generación limpia nacional
\item \textbf{Precisión} : Reducción del 90\% en errores de medición
\item \textbf{Oportunidad} : Datos disponibles en tiempo real vs. 10 días
\item \textbf{Confiabilidad} : Validación cruzada automática entre fuentes

\end{itemize}
\paragraph{Beneficios Regulatorios}
\nopagebreak

\begin{itemize}
\item \textbf{Transparencia} : Trazabilidad completa de datos
\item \textbf{Equidad} : Tratamiento homogéneo de todas las fuentes
\item \textbf{Cumplimiento} : Adherencia automática a estándares de calidad
\item \textbf{Auditoría} : Capacidades avanzadas de verificación

\end{itemize}
\paragraph{Beneficios Ambientales}
\nopagebreak

\begin{itemize}
\item \textbf{Precisión Ambiental} : Medición exacta de generación limpia
\item \textbf{Incentivos Correctos} : Señales precisas para inversión
\item \textbf{Monitoreo} : Seguimiento real del impacto ambiental
\item \textbf{Planificación} : Datos confiables para política energética


\end{itemize}
\subsection{D. Propuestas de Ajuste Normativo}
\nopagebreak

\subsubsection{Instrumento A: Disposiciones CNE conforme Art. 182 y 190 RLSE 2025}
\nopagebreak

\textbf{Fundamento jurídico:} Artículo 182 RLSE 2025 faculta a la CNE para emitir disposiciones administrativas de carácter general que regulen el otorgamiento, liquidación, cancelación voluntaria y transacciones de CEL. Artículo 190 RLSE 2025 faculta a la CNE para establecer términos de acreditación y medición.


\textbf{Disposición propuesta - Medición Universal:} \textquotedbl{}El S-CEL integrará datos de medición de toda la generación limpia nacional, incluyendo gran escala interconectada, generación distribuida y abasto aislado, mediante sistemas automatizados de captura y validación, garantizando la soberanía de datos conforme a la LSE 2025.\textquotedbl{}


\textbf{Disposición propuesta - Validación Cruzada Automática:} \textquotedbl{}El sistema implementará algoritmos de validación cruzada en tiempo real entre todas las fuentes de información, generando alertas automáticas cuando las inconsistencias superen el 1\% o los umbrales estadísticos establecidos por la CNE.\textquotedbl{}


\textbf{Disposición propuesta - Uso Limitado de Estimaciones:} \textquotedbl{}Los valores estimados solo podrán utilizarse cuando la información real no esté disponible por causas de fuerza mayor, por un período máximo de 30 días, y deberán cumplir con criterios de calidad mínimos del 95\% de precisión establecidos por la CNE.\textquotedbl{}


\textbf{Disposición propuesta - Medición Vinculante Integral:} \textquotedbl{}Los datos de medición vinculante incluirán información de CENACE, sistemas de medición inteligente para GLD, telemetría para abasto aislado, y validación cruzada automática entre todas las fuentes, conforme a la metodología establecida en el Art. 190 RLSE 2025.\textquotedbl{}


\textbf{Transición del marco anterior:} Las disposiciones vigentes (RES/174/2016 Disp. 26, A/067/2017 Disp. 32.B) permanecerán aplicables hasta que la CNE emita las nuevas disposiciones conforme al RLSE 2025.


\subsubsection{Instrumento B: Manual de Medición y Calidad de Datos}
\nopagebreak

\textbf{Capítulo 1: Estándares de Medición}


\begin{itemize}
\item Especificaciones técnicas para equipos de medición
\item Protocolos de comunicación y transmisión de datos
\item Requisitos de precisión por tipo de tecnología
\item Procedimientos de calibración y mantenimiento

\end{itemize}
\textbf{Capítulo 2: Validación y Calidad}


\begin{itemize}
\item Algoritmos de validación cruzada
\item Métricas de calidad de datos
\item Procedimientos de detección de anomalías
\item Protocolos de corrección de errores

\end{itemize}
\textbf{Capítulo 3: Integración de Sistemas}


\begin{itemize}
\item APIs para integración con sistemas externos
\item Formatos estándar de intercambio de datos
\item Protocolos de seguridad para transmisión
\item Procedimientos de respaldo y recuperación

\end{itemize}
\textbf{Capítulo 4: Auditoría y Verificación}


\begin{itemize}
\item Procedimientos de auditoría independiente
\item Muestreo estadístico para verificación
\item Protocolos de investigación de inconsistencias
\item Documentación de ajustes y correcciones


\end{itemize}
\subsection{E. Observaciones Técnicas Relevantes}
\nopagebreak

\subsubsection{Riesgos Identificados}
\nopagebreak

\begin{enumerate}
\item \textbf{Riesgo Tecnológico} : Dependencia de sistemas de comunicación para áreas remotas
\item \textbf{Riesgo de Integración} : Complejidad de integrar múltiples sistemas heterogéneos
\item \textbf{Riesgo de Calidad} : Posible degradación inicial durante período de transición
\item \textbf{Riesgo Económico} : Costos significativos de implementación de infraestructura
\item \textbf{Riesgo Operativo} : Resistencia al cambio por parte de participantes actuales

\end{enumerate}
\subsubsection{Supuestos Operativos}
\nopagebreak

\begin{enumerate}
\item \textbf{Infraestructura de Comunicaciones} : Disponibilidad de redes celulares/satelitales
\item \textbf{Capacidad Técnica} : Personal especializado para operación de sistemas avanzados
\item \textbf{Cooperación de Participantes} : Colaboración en implementación de nuevos sistemas
\item \textbf{Recursos Presupuestales} : Financiamiento para modernización tecnológica
\item \textbf{Estabilidad Regulatoria} : Marco normativo estable durante implementación

\end{enumerate}
\subsubsection{Dependencias con Otros Puntos del Bloque}
\nopagebreak

\begin{itemize}
\item \textbf{Punto 1 (Inscripción)} : Los participantes inscritos deben proporcionar datos de medición
\item \textbf{Punto 2 (Administración)} : El sistema de administración debe procesar los datos de medición

\end{itemize}
\subsubsection{Dependencias con Otros Bloques}
\nopagebreak

\begin{itemize}
\item \textbf{Bloque II (Otorgamiento)} : Los datos de medición son base para otorgamiento de CEL
\item \textbf{Bloque III (Disponibilidad)} : La medición precisa es esencial para calcular disponibilidad real
\item \textbf{Bloque VII (Cumplimiento)} : Los datos de medición son base para verificación de obligaciones

\end{itemize}
\subsubsection{Consideraciones de Implementación}
\nopagebreak

\begin{enumerate}
\item \textbf{Implementación Gradual} : Piloto en regiones específicas antes de despliegue nacional
\item \textbf{Capacitación Intensiva} : Programa de capacitación para operadores y participantes
\item \textbf{Período de Transición} : Operación paralela de sistemas durante 6 meses
\item \textbf{Monitoreo Continuo} : Seguimiento de métricas de calidad durante implementación
\item \textbf{Ajustes Iterativos} : Mejora continua basada en experiencia operativa


\end{enumerate}
\subsection{F. Disposiciones Transitorias: Migración al Marco LSE 2025}
\nopagebreak

\subsubsection{Vigencia del Marco Anterior}
\nopagebreak

Las disposiciones del marco anterior (RES/174/2016, A/067/2017) permanecen vigentes hasta que la CNE emita las nuevas disposiciones administrativas de carácter general conforme a los artículos 182 y 190 del RLSE 2025.


\vspace{0.4cm}
\Needspace{25\baselineskip}
\subsubsection{Cronograma de Transición Propuesto}
\nopagebreak

\vspace{0.3cm}
\Needspace{15\baselineskip}
\begin{tabladoradoLargo}
\renewcommand{\tabularxcolumn}[1]{p{#1}}
\begin{xltabular}{\textwidth}{|X|X|X|X|}
\caption{Matriz de Datos} \\
\toprule
\rowcolor{gobmxDorado} \encabezadotablaDorada{Fase} & \encabezadotablaDorada{Actividad} & \encabezadotablaDorada{Plazo} & \encabezadotablaDorada{Fundamento} \\
\midrule
\endhead
\textbf{Fase 1} & Emisión de disposiciones CNE para medición y acreditación (Art. 190 RLSE 2025) & 6 meses desde entrada en vigor RLSE 2025 & Art. 190 RLSE 2025 \\ \hline
\textbf{Fase 2} & Actualización del S-CEL para integrar nuevas metodologías de medición & 12 meses & Art. 182 RLSE 2025 \\ \hline
\textbf{Fase 3} & Migración de datos y validación cruzada con sistemas CNE & 18 meses & Art. 182, 190 RLSE 2025 \\ \hline
\textbf{Fase 4} & Operación completa bajo marco LSE 2025 & 24 meses & LSE 2025 completa \\ \hline
\bottomrule
\end{xltabular}
\end{tabladoradoLargo}

\subsubsection{Compatibilidad de Datos}
\nopagebreak

Los datos de medición capturados bajo el marco anterior (RES/174/2016, A/067/2017) serán compatibles con el nuevo sistema CNE, garantizando la continuidad operativa y la trazabilidad histórica de CEL.


\subsubsection{Responsabilidades durante la Transición}
\nopagebreak

\begin{itemize}
\item \textbf{CNE:} Emitir disposiciones administrativas conforme Art. 182 y 190 RLSE 2025
\item \textbf{CENACE:} Continuar reportando datos de medición conforme a disposiciones vigentes
\item \textbf{Participantes del S-CEL:} Cumplir con requisitos de medición de ambos marcos durante el período transitorio



\end{itemize}
\portadaseccion{BLOQUE II}{Modalidades Operativas de Otorgamiento del CEL (Marco LSE 2025)}{}
\clearpage
\phantomsection
\addcontentsline{toc}{section}{Modalidades Operativas de Otorgamiento del CEL (Marco LSE 2025)}
\section*{Modalidades Operativas de Otorgamiento del CEL (Marco LSE 2025)}
\nopagebreak


\subsection{1. Auditoría de la \textquotedbl{}Foto Actual\textquotedbl{} (Análisis del Marco Anterior)}
\nopagebreak

El sistema de otorgamiento heredado del marco de la LIE (2014) y las resoluciones de la extinta CRE presentaba fallas estructurales que justifican la migración inmediata al nuevo modelo de soberanía energética:


\subsubsection{Instrumentos Citados e Implícitos:}
\nopagebreak

\begin{itemize}
\item \textbf{RES/174/2016 y A/067/2017 (DACG S-CEL):} Establecían un sistema basado en la \textquotedbl{}buena fe\textquotedbl{} del reporte del participante, con validaciones mensuales ex-post.
\item \textbf{RES/1838/2016 y NOM-017-CRE-2019:} Metodologías de Energía Libre de Combustible (ELC) que requerían dictámenes técnicos de Unidades Acreditadas externas, encareciendo el proceso para el generador.
\item \textbf{Ley de Transición Energética (LTE):} Enfocada en metas de mercado, omitiendo la soberanía de datos y el impacto social.

\end{itemize}
\subsubsection{Diagnóstico de Falla Operativa:}
\nopagebreak

\begin{enumerate}
\item \textbf{Burocracia y Fragmentación:} El flujo de información entre CENACE, Transportistas y Distribuidores no era síncrono. Esto generaba retrasos de hasta 60 días en el otorgamiento de un CEL desde que la energía era generada.
\item \textbf{Costos de Transacción Elevados:} La dependencia de \textquotedbl{}Unidades Acreditadas\textquotedbl{} privadas para validar la eficiencia creaba un mercado secundario de certificaciones que no agregaba valor al Sistema Eléctrico Nacional (SEN).
\item \textbf{Desconexión de Datos:} La falta de una API unificada obligaba a procesos de conciliación manual, aumentando el riesgo de errores y duplicidad de certificados.


\end{enumerate}
\subsection{2. Reingeniería Legal: El Nuevo Marco LSE 2025}
\nopagebreak

Bajo la \textbf{Ley del Sector Eléctrico (LSE) 2025 \textbf{ y el }Reglamento de la Ley del Sector Eléctrico (RLSE) 2025 \textbf{, el otorgamiento del CEL se transforma de un trámite administrativo en un }acto de fe pública automatizado}.


\subsubsection{Sustitución de Referencias y Atribuciones:}
\nopagebreak

\begin{itemize}
\item \textbf{De CRE a CNE: \textbf{ La }Comisión Nacional de Energía (CNE) \textbf{ asume la facultad exclusiva de otorgar los CEL conforme al }Artículo 11, Fracción XVI de la LSE 2025}.
\item \textbf{Procedimiento y Requisitos: \textbf{ Los requisitos mínimos para la solicitud de permisos y el otorgamiento se rigen por el }Artículo 23 del Reglamento (RLSE 2025)}, el cual establece los requisitos documentales y técnicos que deben cumplir las solicitudes, incluyendo la obligatoriedad de presentar el número de acuse de recepción o autorización definitiva de la Manifestación de Impacto Social del Sector Energético.
\item \textbf{Soberanía de Datos: \textbf{ Se establece la obligatoriedad de }Sistemas de Medición con sincronía de tiempo} (Art. 2, Fracción XXVII RLSE), permitiendo la validación directa por parte de la CNE sin reportes manuales.
\item \textbf{Justicia Energética: \textbf{ El otorgamiento se vincula a los objetivos de }Justicia Energética} (Art. 6, Fracción II LSE) y respeto a los derechos humanos y sociales (Art. 23 RLSE).


\end{itemize}
\vspace{0.4cm}
\Needspace{25\baselineskip}
\subsection{3. Tabla Comparativa: Antes vs. Ahora}
\nopagebreak

\vspace{0.3cm}
\Needspace{15\baselineskip}
\begin{tablaParadigmaV2}{Análisis Comparativo}{Marco Anterior (LIE/CRE)}{Nuevo Marco (LSE/CNE 2025)}
\tpRow{\textbf{Autoridad Emisora}}{CRE (Unidad de Electricidad)}{\textbf{Autoridad Emisora}}{CNE (Art. 11-XVI LSE)}
\tpRow{\textbf{Validación Técnica}}{Unidades Acreditadas Privadas}{\textbf{Validación Técnica}}{Interoperabilidad CNE-CENACE (Automatizada)}
\tpRow{\textbf{Periodicidad}}{Mensual / Reporte Declarativo}{\textbf{Periodicidad}}{Síncrona / Liquidación Inicial MEM (Real-Time)}
\tpRow{\textbf{Requisito Social}}{Omitido / Evaluación pasiva}{\textbf{Requisito Social}}{Código MISSE (Obligatorio - Art. 23 RLSE)}
\tpRow{\textbf{Enfoque}}{Eficiencia Económica de Mercado}{\textbf{Enfoque}}{Soberanía y Justicia Energética (Art. 6-II LSE)}
\end{tablaParadigmaV2}


\subsection{4. Detalle de Variables Críticas (Datos Exactos)}
\nopagebreak

Para que el nuevo sistema de la CNE otorgue un CEL, el expediente digital debe contener:


\begin{enumerate}
\item \textbf{Datos de Medición (MWh): \textbf{ Inyección Neta validada vía API de CENACE con estampa de tiempo síncrona según el }Art. 2-XXVII del RLSE 2025}.
\item \textbf{Código MISSE: \textbf{ Identificador único de la }Manifestación de Impacto Social del Sector Energético \textbf{ conforme al }Artículo 23 del RLSE}.
\item \textbf{Validación de Permiso: \textbf{ Estatus vigente bajo las figuras de }Generadora \textbf{ o }Autoconsumo} (Art. 17 y 30 LSE).
\item \textbf{Eficiencia (solo Cogeneración/Fósiles): \textbf{ Certificación automática vía CNE conforme a lo dispuesto en el }Art. 11-XIX de la LSE 2025}.


\end{enumerate}
\subsection{5. Régimen Transitorio: Migración Operativa}
\nopagebreak

Para garantizar la continuidad sin detener la operación del SEN, se aplican los siguientes criterios conforme a los \textbf{Artículos Transitorios del Reglamento de la Ley del Sector Eléctrico publicado en el DOF el 03/10/2025}:


\begin{itemize}
\item \textbf{T-Sexto:} Los CEL emitidos bajo el marco anterior mantienen su validez jurídica para la liquidación de obligaciones pasadas. La migración al nuevo sistema electrónico será comunicada oficialmente por la CNE.
\item \textbf{T-Séptimo:} Los participantes con procesos administrativos abiertos pueden regularizarse mediante el pago de multas y la solicitud de migración a las nuevas figuras de la LSE.
\item \textbf{T-Decimonoveno: \textbf{ Los titulares de permisos de autoabastecimiento y cogeneración bajo la LIE tienen }120 días naturales} para iniciar su trámite de migración a las figuras de la LSE 2025.
\item \textbf{Soberanía del Registro: \textbf{ El Registro de CEL de la CNE es la }única fuente de verdad}. Los registros del S-CEL anterior serán migrados masivamente previo proceso de auditoría de datos.

\end{itemize}
\phantomsection
\addcontentsline{toc}{section}{Dictámenes Técnicos, Verificación CNE y Régimen de Transición}
\section*{Dictámenes Técnicos, Verificación CNE y Régimen de Transición}
\nopagebreak

\textbf{Alcance del punto:} Reingeniería del proceso de certificación de atributos limpios, transición de Unidades Acreditadas a esquemas de validación de datos CNE/CENACE y reglas para nucleoeléctricas.

\textbf{Instrumentos jurídicos base:} Ley del Sector Eléctrico 2025 (LSE), Reglamento de la LSE 2025, RES/2910/2017 (Marco anterior referencia).


\vspace{0.4cm}
\Needspace{25\baselineskip}
\subsection{A. Fuentes de información del S-CEL (Nuevo Marco LSE 2025)}
\nopagebreak

\vspace{0.3cm}
\Needspace{15\baselineskip}
\begin{tabladoradoLargo}
\renewcommand{\tabularxcolumn}[1]{p{#1}}
\begin{xltabular}{\textwidth}{|X|X|X|X|}
\caption{Matriz de Datos} \\
\toprule
\rowcolor{gobmxDorado} \encabezadotablaDorada{Actor / Fuente} & \encabezadotablaDorada{Instrumento legal} & \encabezadotablaDorada{Artículo / Numeral} & \encabezadotablaDorada{Cita textual / Aplicación} \\
\midrule
\endhead
\textbf{CNE (Autoridad)} & Reglamento LSE 2025 & Art. 182 & \textquotedbl{}Los Certificados de Energías Limpias son otorgados por la CNE en función de los criterios que para tal efecto emita la Secretaría... y emitir las disposiciones administrativas... que regulen el otorgamiento\textquotedbl{} \\ \hline
\textbf{Sistemas de Medición} & Reglamento LSE 2025 & Art. 2, Fracc. XXVII & \textquotedbl{}Conjunto de elementos que permiten la adquisición, transmisión... de datos de un medidor... [incluyendo] Un sistema de sincronía de tiempo\textquotedbl{} \\ \hline
\textbf{Generación Inyectada} & Reglamento LSE 2025 & Art. 2, Fracc. XIV & \textquotedbl{}Generación de Electricidad Inyectada Total: Suma de la generación de energía eléctrica neta inyectada al Sistema Eléctrico Nacional... conforme a la metodología que emita la Secretaría\textquotedbl{} \\ \hline
\textbf{Validez Histórica} & Reglamento LSE 2025 & Transitorio Sexto & \textquotedbl{}Los Certificados de Energías Limpias emitidos con anterioridad... conservan su validez... hasta su liquidación o cancelación voluntaria\textquotedbl{} \\ \hline
\textbf{Dictamen Técnico (Antecedente)} & RES/2910/2017 & Numeral 3.3 & \textit{Referencia Histórica:} \textquotedbl{}Exposición por escrito... que avala el cumplimiento... derivado de una visita en campo\textquotedbl{} (Se sustituye por validación digital) \\ \hline
\textbf{Transición de Medición} & Reglamento LSE 2025 & Transitorio Décimo Cuarto & \textquotedbl{}Las Centrales Eléctricas... deben informar a la CNE el estatus de sus Sistemas de Medición... La CNE... [establecerá] los periodos y mecanismos para transitar las actualizaciones\textquotedbl{} \\ \hline
\bottomrule
\end{xltabular}
\end{tabladoradoLargo}

\vspace{0.4cm}
\Needspace{25\baselineskip}
\subsection{B. Matriz de validación jurídica (Reingeniería LSE 2025)}
\nopagebreak

\vspace{0.3cm}
\Needspace{15\baselineskip}
\begin{tabladoradoLargo}
\renewcommand{\tabularxcolumn}[1]{p{#1}}
\begin{xltabular}{\textwidth}{|X|X|X|X|}
\caption{Matriz de Datos} \\
\toprule
\rowcolor{gobmxDorado} \encabezadotablaDorada{Hallazgo o limitación (Marco Anterior)} & \encabezadotablaDorada{Fundamento Anterior (RES/2910/2017)} & \encabezadotablaDorada{Solución Nuevo Marco (LSE 2025)} & \encabezadotablaDorada{Ajuste Operativo Inmediato} \\
\midrule
\endhead
\textbf{Vacío Nuclear:} Ausencia de reglas de medición para CEL en nucleoeléctricas. & Transitorio Cuarto RES/2910/2017 (Uso de analogía de cogeneración) & \textbf{Art. 182 Reglamento LSE: \textbf{ Facultad directa de CNE para otorgar CELs bajo criterios SENER. }Transitorio 14:} CNE define transición de medición. & \textbf{Emitir DACG Específica:} Establecer \textquotedbl{}Piso Técnico Nuclear\textquotedbl{} basado en promedio histórico validado por CENACE (eliminando dictamen externo). \\ \hline
\textbf{Conflicto de Interés:} Unidades Acreditadas certificando a empresas de su propio grupo económico. & RES/2910/2017 (Sin prohibición explícita) & \textbf{Art. 6 Reglamento LSE:} Principios de transparencia, imparcialidad, honestidad y buena fe en la actuación administrativa. & \textbf{Veto Normativo:} Prohibir en las nuevas Disposiciones la certificación cruzada intragrupo bajo pena de revocación. \\ \hline
\textbf{Desconexión de Datos:} Dictámenes estáticos (\textquotedbl{}foto\textquotedbl{}) vs operación real. & RES/2910/2017 Numeral 21 (Equipos instalados y calibrados) & \textbf{Art. 2 XXVII Reglamento LSE:} \textquotedbl{}Sistema de Medición\textquotedbl{} integral (telemetría + sincronía). & \textbf{Validación Digital:} Sustituir la visita física por la integración de datos API CENACE para tecnologías puras (Solar/Eólica). \\ \hline
\textbf{Eficiencia Económica:} No considerada para priorizar despacho/CEL. & RES/003/2011 & \textbf{Art. 2 Fracc. X Reglamento LSE:} Definición de \textquotedbl{}Eficiencia\textquotedbl{} incluye minimización de costo total a largo plazo. & \textbf{Criterio de Prioridad:} La validación de CELs priorizará centrales que demuestren eficiencia operativa y menores costos sistémicos. \\ \hline
\bottomrule
\end{xltabular}
\end{tabladoradoLargo}

\subsection{C. Desarrollo analítico}
\nopagebreak

\subsubsection{Diagnóstico: Fallas del Modelo Anterior (RES/2910/2017)}
\nopagebreak

\begin{enumerate}
\item \textbf{Redundancia de Costos:} Se priorizaba la burocracia sobre la física, obligando a centrales con tecnologías obvias y trazables (como fotovoltaica o eólica, donde no hay combustible que auditar) a pagar cuantiosas sumas a terceros privados para certificar lo evidente.
\item \textbf{Riesgo de Datos (Dictamen \textquotedbl{}Foto\textquotedbl{}):} El modelo de Unidad Acreditada se basaba en certificar una \textquotedbl{}foto\textquotedbl{} estática de la planta cada 3 años. Esto dejaba ciega a la autoridad sobre cambios operativos diarios, desviaciones de eficiencia o modificaciones no reportadas durante el periodo de vigencia.
\item \textbf{Privatización de la Fe Pública:} Se delegó la validación del atributo \textquotedbl{}limpio\textquotedbl{} a entes privados con incentivos comerciales cruzados, restando soberanía a la CNE y al CENACE sobre la veracidad de los datos que sustentan el mercado.

\end{enumerate}
\subsubsection{Estado Objetivo (Modelo LSE 2025)}
\nopagebreak

El nuevo marco LSE 2025 plantea un cambio de paradigma: migrar de un modelo de \textbf{Certificación Externalizada \textbf{ (basado en papeles y visitas esporádicas) a un modelo de }Soberanía de Datos} (basado en flujo de información continuo). En este esquema, la CNE recupera la facultad de validación directa, apoyándose en la infraestructura de medición del CENACE para automatizar la emisión de CELs, eliminando intermediarios innecesarios y costos de transacción.


\vspace{0.4cm}
\Needspace{25\baselineskip}
\subsubsection{Tabla comparativa: Modelo Anterior vs. Nuevo Modelo LSE}
\nopagebreak

\vspace{0.3cm}
\Needspace{15\baselineskip}
\begin{tablaParadigmaV2}{Análisis Comparativo}{Modelo Anterior (Unidades Acreditadas)}{Nuevo Modelo (Validación CNE/CENACE - LSE 2025)}
\tpRow{\textbf{Centrales Nucleoeléctricas}}{Sin regla clara (analogía cogeneración). Dependiente de dictamen externo.}{\textbf{Centrales Nucleoeléctricas}}{\textbf{Medición Directa.} Validación de inyección neta por CENACE + Historial CNSNS.}
\tpRow{\textbf{Renovación}}{Trámite manual documental y visita presencial cada 1 o 3 años.}{\textbf{Renovación}}{\textbf{Automática (Tácita).} Si la telemetría es continua y válida, la certificación se mantiene.}
\tpRow{\textbf{Costo para el Generador}}{Alto (Honorarios recurrentes de Unidad Acreditada).}{\textbf{Costo para el Generador}}{\textbf{Bajo/Nulo.} Proceso administrativo sustentado en datos operativos existentes.}
\tpRow{\textbf{Vigencia}}{Periodos fijos rígidos.}{\textbf{Vigencia}}{\textbf{Indefinida (Condicionada).} Sujeta a la transmisión ininterrumpida de datos de medición.}
\tpRow{\textbf{Rol de Terceros}}{Indispensable y obligatorio para todos.}{\textbf{Rol de Terceros}}{\textbf{Acotado.} Solo para tecnologías complejas (cogeneración, bioenergía) que requieren auditoría física de combustible.}
\end{tablaParadigmaV2}

\subsubsection{Arquitectura del Sistema de Validación}
\nopagebreak

Para implementar el Artículo 182 del Reglamento LSE, se propone la siguiente arquitectura lógica:


\begin{enumerate}
\item \textbf{Módulo de Ingesta (CENACE):} Recibe en tiempo real las variables críticas del Sistema de Medición (MWh consumidos, MWh entregados, Timestamp, Estatus Operativo).
\item \textbf{Motor de Reglas (CNE):}
\end{enumerate}
\begin{itemize}
\item \textit{Regla 1 (Tecnología Pura):} Si la fuente es Solar/Eólica/Nuclear/Hidro y la telemetría es válida  $\rightarrow$  \textbf{Certificación Automática}.
\item \textit{Regla 2 (Tecnología Mixta):} Si la fuente es Cogeneración/Bioenergía  $\rightarrow$  Requiere input adicional de flujo de combustible (Medible) o inspección semestral de CNE (si no hay telemetría de combustible).
\end{itemize}
\begin{enumerate}
\item \textbf{Registro S-CEL:} Emisión de certificados basada estrictamente en la \textquotedbl{}Generación de Electricidad Inyectada Total\textquotedbl{} (Art. 2 Fracc. XIV RLSE) validada por el motor de reglas, desacoplando la emisión del proceso comercial.

\end{enumerate}
\subsection{D. Propuestas de ajuste normativo (Texto para DACG CNE)}
\nopagebreak

Se proponen las siguientes redacciones para ser integradas en las nuevas Disposiciones Administrativas de Carácter General (DACG) que emitirá la CNE, fundamentadas en el Art. 182 y 190 del Reglamento LSE 2025.


\subsubsection{Instrumento A: Modificación a las Disposiciones del S-CEL (CNE)}
\nopagebreak

\textbf{Disposición 11.3 (Nuevo - Vía Rápida Nuclear):}


\vspace{0.2cm}
\Needspace{8\baselineskip}
\begin{glowBox}[gobmxDorado]{Nota Informativa}
\textquotedbl{}Para el otorgamiento de Certificados de Energías Limpias a centrales nucleoeléctricas propiedad del Estado o con participación estatal mayoritaria, la CNE validará directamente la Generación de Electricidad Inyectada Total conforme a los registros operativos del CENACE y los reportes de seguridad de la CNSNS. Se exime a estas centrales de la presentación de Dictamen Técnico de terceros, sustituyéndose este requisito por el reporte oficial de medición validado por el operador del sistema.\textquotedbl{}
\end{glowBox}

\textbf{Disposición 34.A (Nuevo - Independencia y Conflicto de Interés):}


\vspace{0.2cm}
\Needspace{8\baselineskip}
\begin{glowBox}[gobmxDorado]{Nota Informativa}
\textquotedbl{}En los casos excepcionales donde la metodología de cálculo requiera la verificación de terceros (específicamente para tecnologías con uso de combustibles fósiles o bioenergía compleja), las unidades de inspección o entidades de verificación \textbf{no podrán tener vinculación societaria, comercial o de grupo económico} con la central generadora a certificar. La violación a esta disposición, conforme a los principios de honestidad e imparcialidad del Art. 6 del Reglamento, resultará en la revocación inmediata de la acreditación y la nulidad de los dictámenes emitidos.\textquotedbl{}
\end{glowBox}

\textbf{Disposición 28.3 (Digitalización y Vigencia):}


\vspace{0.2cm}
\Needspace{8\baselineskip}
\begin{glowBox}[gobmxDorado]{Nota Informativa}
\textquotedbl{}La vigencia de la certificación de Central Eléctrica Limpia será indefinida, quedando sujeta a una \textbf{verificación digital continua} a través de los Sistemas de Medición validados por el CENACE. Cualquier interrupción en la transmisión de datos de medición por un periodo mayor a 24 horas suspenderá automáticamente la elegibilidad de la central para recibir CELs por la energía generada en dicho intervalo, restableciéndose una vez normalizada la telemetría, conforme a los plazos de transición del Transitorio Décimo Cuarto del Reglamento de la LSE.\textquotedbl{}
\end{glowBox}

\subsection{E. Estrategia de Transición (Transitorios)}
\nopagebreak

Para asegurar una migración ordenada del marco de la Ley de la Industria Eléctrica (LIE) al nuevo marco de la Ley del Sector Eléctrico (LSE) 2025, se establecen las siguientes directrices transitorias:


\begin{enumerate}
\item \textbf{Validez de CELs Históricos (Seguridad Jurídica):}
\end{enumerate}
De conformidad con el \textbf{Transitorio Sexto} del Reglamento LSE, se ratifica la plena validez de todos los Certificados de Energías Limpias emitidos con anterioridad a la entrada en vigor del nuevo reglamento, manteniéndose vigentes bajo las condiciones originales de su emisión hasta su liquidación o cancelación voluntaria.

\begin{enumerate}
\item \textbf{Periodo de Gracia para Actualización de Medición (120 días):}
\end{enumerate}
En atención al \textbf{Transitorio Décimo Cuarto}, las Centrales Eléctricas que a la fecha de publicación de las nuevas DACG no cuenten con sistemas de telemetría compatibles con los nuevos estándares de validación de la CNE/CENACE, dispondrán de un plazo de 120 días naturales para regularizar sus sistemas. Durante este periodo de transición, la CNE aceptará como válidos los dictámenes técnicos vigentes emitidos por Unidades Acreditadas bajo el esquema anterior, garantizando que no se detenga la emisión de certificados mientras se moderniza la infraestructura.

\begin{enumerate}
\item \textbf{Migración Administrativa de Permisos:}
\end{enumerate}
Para las centrales que realicen la migración de su Permiso de Generación (de LIE a LSE), el anexo técnico de \textquotedbl{}Atributos Limpios\textquotedbl{} del nuevo permiso funcionará como la certificación base para el S-CEL. Esto eliminará la necesidad de un trámite separado de \textquotedbl{}Dictamen Técnico\textquotedbl{} inicial para el 80\% de los casos (tecnologías puras como Solar, Eólica, Hidroeléctrica y Nuclear), simplificando la carga administrativa y centrando el control en la operación real.


\phantomsection
\addcontentsline{toc}{section}{Conteo, fraccionamiento y acumulación de Certificados de Energía Limpia (S-CEL 2.0)}
\section*{Conteo, fraccionamiento y acumulación de Certificados de Energía Limpia (S-CEL 2.0)}
\nopagebreak

\textbf{Alcance:} Definición de reglas algorítmicas y normativas para la medición, acumulación de fracciones y emisión de CEL bajo la autoridad de la Comisión Nacional de Energía (CNE).

\textbf{Instrumentos jurídicos base:} Ley del Sector Eléctrico 2025 (LSE), Ley de Planeación y Transición Energética (LPTE).


\subsection{A. Auditoría de Fuentes: Transición del Marco Anterior al Nuevo}
\nopagebreak

Esta tabla contrasta las fuentes operativas originales (CRE/RES/174/2016) con los nuevos mandatos de la CNE para justificar la reingeniería del proceso.


\vspace{0.3cm}
\Needspace{15\baselineskip}
\begin{tabladoradoLargo}
\renewcommand{\tabularxcolumn}[1]{p{#1}}
\begin{xltabular}{\textwidth}{|X|X|X|X|}
\caption{Matriz de Datos} \\
\toprule
\rowcolor{gobmxDorado} \encabezadotablaDorada{Concepto} & \encabezadotablaDorada{Referencia Anterior (RES/174/2016 - CRE)} & \encabezadotablaDorada{Limitación Operativa Detectada} & \encabezadotablaDorada{Nuevo Fundamento (LSE 2025 / LPTE - CNE)} \\
\midrule
\endhead
\textbf{Autoridad Emisora} & CRE (Ley de la Industria Eléctrica 2014) & Fragmentación de datos entre CRE y CENACE. Burocracia en validación manual. & \textbf{CNE (Comisión Nacional de Energía).} \textquotedbl{}La CNE está facultada para: ... Otorgar los Certificados de Energías Limpias\textquotedbl{} y \textquotedbl{}Verificar el cumplimiento de los requisitos\textquotedbl{}. \\ \hline
\textbf{Medición y Datos} & Disposición 32 (Acumulación simple). Dependencia de reporte de terceros. & Imprecisión en decimales provocaba pérdidas de valor en generadores pequeños (Micro-generación). & \textbf{Soberanía del Dato.} \textquotedbl{}La CNE puede establecer requerimientos de estimación, medición, y reporte relacionado con la generación de Energías Limpias\textquotedbl{} (Art. 147 Fracc. V LSE 2025). \\ \hline
\textbf{Definición de Limpieza} & Lineamientos CEL (Tecnología + Combustible). & Ambigüedad en co-procesamiento y centrales híbridas. & \textbf{Taxonomía Estricta.} La LSE define explícitamente las Energías Limpias (Art. 3, XXI, incisos a-m), incluyendo criterios de eficiencia de la CNE para cogeneración e hidrógeno. \\ \hline
\textbf{Base de Otorgamiento} & Mercado impulsado por oferta/demanda privada. & Volatilidad y falta de alineación con metas nacionales. & \textbf{Planeación Vinculante.} El otorgamiento debe alinearse a las \textquotedbl{}metas de transición energética, descarbonización\textquotedbl{} establecidas por la Secretaría (SENER) (Art. 147 LSE). \\ \hline
\textbf{Migración} & N/A & Riesgo de vacío legal durante el cambio de sistema. & \textbf{Continuidad Regulada.} Los permisos anteriores siguen surtiendo efectos hasta su vigencia, pero se debe facilitar la migración expedita a la nueva Ley (Transitorios Noveno y Sexto LSE). \\ \hline
\bottomrule
\end{xltabular}
\end{tabladoradoLargo}

\subsection{B. Matriz de Validación Jurídica (LSE 2025)}
\nopagebreak

Se proponen ajustes operativos sustentados en las nuevas facultades de la CNE para eliminar las \textquotedbl{}zonas grises\textquotedbl{} del conteo y garantizar la Justicia Energética.


\vspace{0.3cm}
\Needspace{15\baselineskip}
\begin{tabladoradoLargo}
\renewcommand{\tabularxcolumn}[1]{p{#1}}
\begin{xltabular}{\textwidth}{|X|X|X|X|}
\caption{Matriz de Datos} \\
\toprule
\rowcolor{gobmxDorado} \encabezadotablaDorada{Hallazgo / Limitación Actual} & \encabezadotablaDorada{Fundamento Jurídico LSE 2025 (Cita Textual)} & \encabezadotablaDorada{Riesgo Mitigado} & \encabezadotablaDorada{Ajuste Técnico Propuesto} \\
\midrule
\endhead
\textbf{Redondeo impreciso} (pérdida de fracciones \textless{} 0.1 MWh) & \textbf{Art. 10 Fracc. XIX LSE:} \textquotedbl{}Fomentar el otorgamiento de créditos... para el financiamiento de Centrales Eléctricas de Generación Limpia Distribuida\textquotedbl{}. & Subestimación de CELs afecta la viabilidad financiera de la Generación Distribuida (Justicia Energética). & \textbf{Regla de 4 Decimales:} Acumulación exacta hasta 0.9999 MWh antes de emitir. \\ \hline
\textbf{Centrales Híbridas} (Mezcla de tecnologías sin desglose) & \textbf{Art. 11 Fracc. XIX LSE:} La CNE debe \textquotedbl{}Emitir los criterios de eficiencia utilizados en la definición de Energías Limpias\textquotedbl{}. & \textquotedbl{}Greenwashing\textquotedbl{} o conteo doble en centrales con respaldo fósil no medido correctamente. & \textbf{Acumulación Segregada:} \textquotedbl{}Wallets\textquotedbl{} de acumulación separadas por tipo de tecnología (ej. Solar vs. Cogeneración). \\ \hline
\textbf{Inconsistencias de Datos} (Diferencias \textgreater{} 2\%) & \textbf{Art. 11 Fracc. XI y XII LSE:} \textquotedbl{}Vigilar la operación del Mercado... e instruir las correcciones que deban realizarse a los parámetros registrados\textquotedbl{}. & Retrasos en la liquidación mensual y disputas administrativas largas. & \textbf{Validación Automática:} Interconexión directa API CNE-CENACE. Plazo de corrección: 15 días máx. \\ \hline
\textbf{Paros y Mantenimientos} (Vacío legal en conteo durante pausas) & \textbf{Art. 3 Fracc. XLIX LSE:} Garantizar \textquotedbl{}Calidad, Confiabilidad, Continuidad y seguridad\textquotedbl{} del Sistema Eléctrico. & Otorgamiento indebido de CELs basado en estimaciones erróneas durante paros. & \textbf{Flag de Estado Operativo:} Suspensión automática de conteo estimado durante estados de \textquotedbl{}Mantenimiento\textquotedbl{}. \\ \hline
\bottomrule
\end{xltabular}
\end{tabladoradoLargo}

\subsection{C. Desarrollo Analítico: Modelo S-CEL 2.0}
\nopagebreak

\subsubsection{Diagnóstico de la situación}
\nopagebreak

El modelo anterior operaba bajo una lógica de \textquotedbl{}mínimos necesarios\textquotedbl{}. El nuevo modelo bajo la LSE 2025 opera bajo una lógica de \textbf{trazabilidad y soberanía}, donde cada fracción de MWh cuenta para las Metas de Transición Energética que la SENER está obligada a reportar en el PLATEASE.


\subsubsection{Arquitectura del Sistema Objetivo}
\nopagebreak

El sistema se rediseña para cumplir con el \textbf{Artículo 147 de la LSE}, que faculta a la CNE para establecer requerimientos de medición:


\begin{enumerate}
\item \textbf{Motor de Ingesta (CNE-Link):} Recibe datos brutos del CENACE (Art. 11 VI LSE).
\end{enumerate}
\begin{itemize}
\item \textit{Input:} `Energía\_Inyectada\_MWh` (precisión `float` 4 decimales).
\item \textit{Input:} `Fuente\_Energia\_ID` (Catálogo LSE Art. 3 XXI).
\end{itemize}
\begin{enumerate}
\item \textbf{Lógica de Fraccionamiento (Art. 11 XX LSE):}
\end{enumerate}
\begin{itemize}
\item Si `Tecnología` == \textquotedbl{}Limpia Pura\textquotedbl{} (Solar/Eólica): Factor 1.0.
\item Si `Tecnología` == \textquotedbl{}Cogeneración/Bio\textquotedbl{}: Factor calculado según eficiencia térmica (Art. 3 XXI k, l, m LSE).
\end{itemize}
\begin{enumerate}
\item \textbf{Bóveda de Acumulación (Ledger):}
\end{enumerate}
\begin{itemize}
\item Mantiene el saldo de fracciones mes a mes (\textquotedbl{}Rollover\textquotedbl{}).
\item \textbf{Regla de corte:} Al alcanzar \$\textbackslash\{\}ge 1.0000\$, se emite 1 CEL y se resta 1.0000 del saldo. El remanente permanece en la cuenta del usuario.

\end{itemize}
\vspace{0.4cm}
\Needspace{25\baselineskip}
\subsubsection{Tabla Comparativa: Modelo Anterior vs Modelo LSE 2025}
\nopagebreak

\vspace{0.3cm}
\Needspace{15\baselineskip}
\begin{tablaParadigmaV2}{Análisis Comparativo}{Modelo RES/174/2016}{Modelo LSE 2025 (Objetivo)}
\tpRow{\textbf{Precisión}}{1 o 2 decimales (pérdida de residuales).}{\textbf{Precisión}}{\textbf{4 decimales} (Art. 147 V LSE) para captura total de valor.}
\tpRow{\textbf{Híbridas}}{Conteo global por central.}{\textbf{Híbridas}}{\textbf{Desglose por tecnología} (Art. 3 XXI LSE).}
\tpRow{\textbf{Validación}}{Manual / Reactiva.}{\textbf{Validación}}{\textbf{Automática / Instructiva} (CNE instruye corrección al CENACE).}
\tpRow{\textbf{Transparencia}}{Datos privados del participante.}{\textbf{Transparencia}}{Datos alineados al \textbf{Sistema Nacional de Información Energética} (Art. 8 XI LPTE).}
\end{tablaParadigmaV2}

\subsection{D. Propuestas de Ajuste Normativo (Nuevas DACG)}
\nopagebreak

Estas propuestas sustituyen las disposiciones de la RES/174/2016 y se integran como \textbf{Disposiciones Administrativas de Carácter General (DACG)} emitidas por la CNE.


\subsubsection{Instrumento A: DACG para la Emisión y Validación de CELs (CNE-DACG-001-2026)}
\nopagebreak

\textbf{Disposición 32.A (Nuevo Estándar de Precisión):}


\vspace{0.2cm}
\Needspace{8\baselineskip}
\begin{glowBox}[gobmxDorado]{Nota Informativa}
\textquotedbl{}De conformidad con la facultad de la CNE establecida en el artículo 147 fracción V de la Ley del Sector Eléctrico, la medición para la emisión de CEL se realizará con una precisión de cuatro cifras decimales. Las fracciones de MWh se acumularán mensualmente en el registro del Participante del Mercado. La emisión del CEL se perfeccionará automáticamente en el momento en que el saldo acumulado de fracciones sea igual o superior a la unidad (1.0000 MWh).\textquotedbl{}
\end{glowBox}

\textbf{Disposición 32.B (Trazabilidad por Tecnología):}


\vspace{0.2cm}
\Needspace{8\baselineskip}
\begin{glowBox}[gobmxDorado]{Nota Informativa}
\textquotedbl{}Para dar cumplimiento a la clasificación de Energías Limpias del artículo 3 fracción XXI de la Ley, la acumulación de fracciones deberá segregarse por tecnología de generación. En centrales híbridas o de cogeneración, se aplicarán los criterios de eficiencia emitidos por la CNE (Art. 11 XIX LSE) antes de la acumulación en el saldo de fracciones.\textquotedbl{}
\end{glowBox}

\textbf{Disposición 30.A (Corrección de Inconsistencias):}


\vspace{0.2cm}
\Needspace{8\baselineskip}
\begin{glowBox}[gobmxDorado]{Nota Informativa}
\textquotedbl{}En ejercicio de la facultad de vigilancia del artículo 11 fracción XI de la Ley, la CNE notificará inconsistencias entre la medición reportada y la registrada por el CENACE. El Generador tendrá un plazo perentorio de 15 días hábiles para subsanar la diferencia. En caso de silencio, prevalecerá el dato del CENACE corregido por la CNE según el artículo 11 fracción XII de la Ley.\textquotedbl{}
\end{glowBox}

\subsection{E. Estrategia de Transición (Transitorios y Migración)}
\nopagebreak

Para evitar la parálisis operativa durante la migración del sistema CRE al sistema CNE, se activan los mecanismos de los artículos transitorios de la LSE 2025.


\begin{enumerate}
\item \textbf{Reconocimiento de Saldos Históricos:}

\end{enumerate}
\begin{itemize}
\item De acuerdo con el \textbf{Transitorio Noveno} de la LSE, los permisos y actos administrativos anteriores continúan surtiendo efectos. La CNE reconocerá los saldos de fracciones acumulados bajo la vigencia de la CRE como \textquotedbl{}Saldo Inicial\textquotedbl{} en el nuevo sistema.
\end{itemize}
\begin{enumerate}
\item \textbf{Ventanilla Única de Aclaración:}

\end{enumerate}
\begin{itemize}
\item Se implementará la \textbf{Ventanilla Única} ordenada en el Transitorio Noveno y Sexto de la LSE para resolver discrepancias de conteo del periodo 2024-2025 de manera expedita, facilitando la migración de Generadores y Generadores Exentos al nuevo esquema.
\end{itemize}
\begin{enumerate}
\item \textbf{Periodo de Coexistencia (180 días):}

\end{enumerate}
\begin{itemize}
\item Durante los primeros 180 días (plazo para actualización de metodologías según \textbf{Transitorio Cuarto}), el sistema S-CEL aceptará reportes en formatos anteriores, ejecutando una conversión interna al estándar de 4 decimales para garantizar la continuidad del despacho y la liquidación.

\end{itemize}
\phantomsection
\addcontentsline{toc}{section}{Trazabilidad jurídica y técnica del Certificado de Energía Limpia (S-CEL 2.0)}
\section*{Trazabilidad jurídica y técnica del Certificado de Energía Limpia (S-CEL 2.0)}
\nopagebreak

\textbf{Alcance:} Protocolo técnico-jurídico para el ciclo de vida del CEL bajo el principio de Planeación Vinculante y Soberanía Energética.

\textbf{Marco Jurídico Vigente:} Ley del Sector Eléctrico (LSE 2025), Reglamento de la LSE 2025 (RLSE), Ley de Planeación y Transición Energética (LPTE).


\subsection{A. Auditoría del Marco Anterior (\textquotedbl{}Foto Actual\textquotedbl{})}
\nopagebreak

\textbf{Diagnóstico de Fallo Operativo:}

El modelo anterior, gestionado por la CRE bajo las Disposiciones Administrativas RES/174/2016, operaba bajo una lógica de mercado especulativo desconectado de la realidad operativa del SEN.


\begin{enumerate}
\item \textbf{Centralización Vulnerable:} La plataforma S-CEL (Sistema de Gestión de Certificados y Cumplimiento de Obligaciones de Energías Limpias) residía en servidores desconectados de la medición fiscal del CENACE, dependiendo de reportes de \textquotedbl{}buena fe\textquotedbl{} de los generadores.
\item \textbf{Falta de Validación Técnica:} No existía un mecanismo robusto para verificar que el MWh reportado correspondiera efectivamente a energía limpia inyectada y no a simulaciones o arbitraje regulatorio.
\item \textbf{Opacidad Legal:} La derogación fáctica de la autoridad de la CRE ha dejado a los participantes en un limbo jurídico sobre quién valida sus atributos ambientales hoy.

\end{enumerate}
\textbf{Estatus:} \textit{Obsoleto y No Vinculante}. Se requiere migración inmediata al nuevo estándar de Soberanía.


\subsection{B. Reingeniería Jurídica (Marco LSE 2025)}
\nopagebreak

El control del sistema pasa de un regulador autónomo a la \textbf{Comisión Nacional de Energía (CNE) \textbf{, como órgano desconcentrado encargado de la ejecución de la política energética dictada por }SENER}.


\vspace{0.4cm}
\Needspace{25\baselineskip}
\subsubsection{B.1. Fundamentos de Autoridad (LSE 2025)}
\nopagebreak

\vspace{0.3cm}
\Needspace{15\baselineskip}
\begin{tabladoradoLargo}
\renewcommand{\tabularxcolumn}[1]{p{#1}}
\begin{xltabular}{\textwidth}{|X|X|X|X|}
\caption{Matriz de Datos} \\
\toprule
\rowcolor{gobmxDorado} \encabezadotablaDorada{Facultad} & \encabezadotablaDorada{Autoridad} & \encabezadotablaDorada{Fundamento Legal (LSE 2025)} & \encabezadotablaDorada{Mandato} \\
\midrule
\endhead
\textbf{Emisión de CELs} & \textbf{CNE} & \textbf{Art. 147 Fracc. III} & \textit{\textquotedbl{}La CNE debe otorgar los Certificados de Energías Limpias que correspondan, emitir la regulación para validar su titularidad y verificar el cumplimiento de dichas obligaciones.\textquotedbl{}} \\ \hline
\textbf{Registro (S-CEL 2.0)} & \textbf{CNE} & \textbf{Art. 149} & \textit{\textquotedbl{}La CNE debe crear y mantener el Registro de Certificados, el cual al menos debe tener el matriculado de cada certificado, así como la información correspondiente a su fecha de emisión, vigencia e historial de personas propietarias.\textquotedbl{}} \\ \hline
\textbf{Política y Criterios} & \textbf{SENER} & \textbf{Art. 147 Fracc. I y II} & La Secretaría establece los requisitos de Transición Energética y los criterios para reconocer a las Generadoras Limpias. \\ \hline
\textbf{Validación de Datos} & \textbf{CENACE} & \textbf{Art. 11 LSE / Art. 124 LSE} & CENACE provee los datos de medición fiscal necesarios para que la CNE valide la generación real antes de la emisión. \\ \hline
\bottomrule
\end{xltabular}
\end{tabladoradoLargo}

\subsubsection{B.2. Reglamento de la LSE (RLSE 2025)}
\nopagebreak

El Reglamento detalla la operatividad del sistema en sus artículos 180 al 194.


\begin{itemize}
\item \textbf{Art. 182 (RLSE):} Establece que la emisión es un acto electrónico de fe pública, no un mero trámite administrativo.
\item \textbf{Art. 184 (RLSE):} Obliga a que todas las transacciones de CELs (compra/venta) se realicen dentro de la plataforma oficial para evitar el doble conteo.

\end{itemize}
\subsection{C. Trazabilidad Técnica (Estándar S-CEL 2.0)}
\nopagebreak

Para dar cumplimiento al \textbf{Art. 149 de la LSE \textbf{ (Matriculado e Historial), se define el nuevo estándar técnico del Certificado. El S-CEL 2.0 no es solo una base de datos, es un }Sistema de Trazabilidad Integral}.


\subsubsection{C.1. Identificador Único (UID) Soberano}
\nopagebreak

Cada MWh de energía limpia genera un activo digital único e irrepetible.

Estructura del Matriculado: `MX-[TEC]-[NODO]-[MES]-[AÑO]-[HASH]`


\begin{enumerate}
\item \textbf{`MX`:} Jurisdicción Soberana (México).
\item \textbf{`[TEC]`:} Tecnología de Generación (SOL, EOL, HID, NUC, BIO, GEO).
\item \textbf{`[NODO]`:} Zona de Carga/Generación del SEN (Trazabilidad Regional).
\item \textbf{`[MES/AÑO]`:} Temporalidad de la inyección (Vigencia de 30 meses según Art. 146 LSE).
\item \textbf{`[HASH]`:} Firma criptográfica que vincula el certificado a la medición fiscal del CENACE.

\end{enumerate}
\subsubsection{C.2. Capa de Metadatos (Expediente Digital)}
\nopagebreak

Al consultar el UID en el Registro de la CNE, el sistema despliega:


\begin{itemize}
\item \textbf{Fuente:} Central Generadora y Permiso CNE asociado.
\item \textbf{Validación:} Fecha y hora de la conciliación con CENACE.
\item \textbf{Estatus Social:} \textit{Cumplimiento / Incumplimiento} (Vinculado a la Evaluación de Impacto Social de SENER).
\item \textbf{Historial de Propiedad:} Cadena de custodia desde el Generador hasta el Usuario Final o Suministrador (Obligado).

\end{itemize}
\vspace{0.4cm}
\Needspace{25\baselineskip}
\subsection{D. Tablas Comparativas (Antes vs Ahora)}
\nopagebreak

\vspace{0.3cm}
\Needspace{15\baselineskip}
\begin{tabladoradoLargo}
\renewcommand{\tabularxcolumn}[1]{p{#1}}
\begin{xltabular}{\textwidth}{|X|X|X|X|}
\caption{Matriz de Datos} \\
\toprule
\rowcolor{gobmxDorado} \encabezadotablaDorada{Concepto} & \encabezadotablaDorada{Modelo Anterior (CRE / LIE)} & \encabezadotablaDorada{Nuevo Modelo (CNE / LSE 2025)} & \encabezadotablaDorada{Impacto} \\
\midrule
\endhead
\textbf{Autoridad Emisora} & Comisión Reguladora de Energía (CRE) & \textbf{Comisión Nacional de Energía (CNE)} & Centralización de la política y ejecución en el Estado. \\ \hline
\textbf{Base de Datos} & S-CEL (Plataforma Privada/Externa) & \textbf{Registro de Certificados (Plataforma Soberana)} & Seguridad nacional y soberanía de datos (Art. 149 LSE). \\ \hline
\textbf{Validación} & Declarativa (Buena Fe) & \textbf{Automatizada (Interoperabilidad)} & Eliminación de fraude y \textquotedbl{}CELs de papel\textquotedbl{} sin respaldo energético. \\ \hline
\textbf{Vigencia} & Indefinida (Especulativa) & \textbf{30 Meses (Art. 146 LSE)} & Fomento a la liquidez y desincentivo al acaparamiento. \\ \hline
\bottomrule
\end{xltabular}
\end{tabladoradoLargo}

\subsection{E. Estrategia de Transición (Transitorios)}
\nopagebreak

Para asegurar la continuidad del mercado y el cumplimiento de las Metas Nacionales durante la migración:


\begin{enumerate}
\item \textbf{Reconocimiento de Legados:} Los CELs emitidos por la CRE antes de la entrada en vigor de la LSE 2025 serán reconocidos en el nuevo Registro, marcados con el estatus \textquotedbl{}LEGADO\textquotedbl{}. Mantendrán su validez hasta su liquidación o expiración (Transitorio Sexto RLSE).
\item \textbf{Ventanilla de Migración:} Los Participantes del Mercado deberán darse de alta en el nuevo sistema de la CNE en un plazo no mayor a 90 días hábiles.
\item \textbf{Validación Retroactiva:} La CNE se reserva el derecho de auditar la emisión histórica de CELs para detectar inconsistencias con los datos de CENACE.

\end{enumerate}
\portadaseccion{BLOQUE III}{Disponibilidad Real de Certificados de Energía Limpia (CEL) y Transición al Nuevo Modelo de Soberanía}{}
\clearpage
\phantomsection
\addcontentsline{toc}{section}{Disponibilidad Real de Certificados de Energía Limpia (CEL) y Transición al Nuevo Modelo de Soberanía}
\section*{Disponibilidad Real de Certificados de Energía Limpia (CEL) y Transición al Nuevo Modelo de Soberanía}
\nopagebreak

\textbf{Alcance:} Diagnóstico de disponibilidad, reingeniería del monitoreo bajo la LSE 2025 y estrategia de transición operativa CRE → CNE.

\textbf{Instrumentos jurídicos rectores:} Ley del Sector Eléctrico 2025 (LSE), Reglamento de la Ley del Sector Eléctrico 2025 (RLSE), Ley de Planeación y Transición Energética (LPTE).



\subsection{A. Marco Normativo: Del Mercado a la Planeación Vinculante}
\nopagebreak

Esta sección contrasta las fuentes históricas (aún relevantes para liquidaciones pendientes) con las nuevas facultades de la CNE y SENER bajo el modelo de soberanía energética.


\vspace{0.3cm}
\Needspace{15\baselineskip}
\begin{tabladoradoLargo}
\renewcommand{\tabularxcolumn}[1]{p{#1}}
\begin{xltabular}{\textwidth}{|X|X|X|X|}
\caption{Matriz de Datos} \\
\toprule
\rowcolor{gobmxDorado} \encabezadotablaDorada{Actor} & \encabezadotablaDorada{Instrumento (Marco Anterior - Referencia)} & \encabezadotablaDorada{Instrumento (Nuevo Marco LSE 2025)} & \encabezadotablaDorada{Disposición Clave (Nueva)} \\
\midrule
\endhead
\textbf{SENER} & Aviso Requisito CEL 2019-2022 & \textbf{LSE 2025 / RLSE} & \textbf{Transitorio Décimo Sexto (RLSE):} \textquotedbl{}Por única ocasión, la Secretaría debe publicar los requisitos de Energía Limpia para los años 2025, 2026, 2027 y 2028...\textquotedbl{} \\ \hline
\textbf{CNE (antes CRE)} & RES/174/2016 (Otorgamiento) & \textbf{LSE Art. 11 Fracc. XVI y XVII} & \textquotedbl{}Otorgar los Certificados de Energías Limpias... y operar el registro de los mismos\textquotedbl{}. \\ \hline
\textbf{CNE} & A/013/2019 (Flexibilidad) & \textbf{LSE Art. 147 Fracc. III} & La CNE debe verificar el cumplimiento de las obligaciones de Transición Energética y Descarbonización. \\ \hline
\textbf{Generación Limpia} & RES/1838/2016 (Eficiencia) & \textbf{RLSE Art. 182} & Los CELs se otorgan en función de criterios que emita SENER y regulaciones de CNE para validar titularidad. \\ \hline
\textbf{Metas Nacionales} & Ley de Transición Energética (LTE) & \textbf{LPTE Art. 24} & La Estrategia es el instrumento rector. Las metas de Energías Limpias son porcentajes mínimos vinculantes. \\ \hline
\textbf{Planeación} & PRODESEN (Indicativo) & \textbf{LSE Art. 12} & \textquotedbl{}La planeación del sector eléctrico tiene carácter \textbf{vinculante} y está a cargo de la Secretaría\textquotedbl{}. \\ \hline
\bottomrule
\end{xltabular}
\end{tabladoradoLargo}


\subsection{B. Matriz de Validación Jurídica y Riesgos Operativos}
\nopagebreak

Identificación de brechas entre el modelo especulativo anterior y el modelo de garantía de suministro actual.


\vspace{0.3cm}
\Needspace{15\baselineskip}
\begin{tabladoradoLargo}
\renewcommand{\tabularxcolumn}[1]{p{#1}}
\begin{xltabular}{\textwidth}{|X|X|X|X|}
\caption{Matriz de Datos} \\
\toprule
\rowcolor{gobmxDorado} \encabezadotablaDorada{Hallazgo / Limitación} & \encabezadotablaDorada{Fundamento Jurídico (LSE 2025)} & \encabezadotablaDorada{Riesgo Operativo} & \encabezadotablaDorada{Ajuste Propuesto (Reingeniería LSE)} \\
\midrule
\endhead
\textbf{Opacidad en Disponibilidad Real:} El sistema anterior reportaba \textit{ex-post} (10 días después del cierre de mes). & \textbf{LSE Art. 11 Fracc. XVII:} CNE debe llevar el registro de CELs y vigilar su integridad. & Incapacidad de reacción ante déficits de CELs, comprometiendo las metas de la LPTE. & Implementar \textbf{Tableros de Control en Tiempo Real} enlazados directamente a la medición de CENACE bajo el nuevo Art. 15 del RLSE (Proyectos Piloto). \\ \hline
\textbf{Criterios de Flexibilidad Financiera:} El A/013/2019 basaba la flexibilidad en precios de mercado (UDIs), no en capacidad técnica. & \textbf{LSE Art. 12 Fracc. I:} Procurar la Confiabilidad y Continuidad del servicio público con responsabilidad social. & La volatilidad financiera puede justificar incumplimientos técnicos inaceptables para la soberanía energética. & Sustituir el \textquotedbl{}Precio Implícito\textquotedbl{} por un \textbf{Índice de Suficiencia de Descarbonización} gestionado por CNE. \\ \hline
\textbf{Gestión de Excedentes:} Los CELs no emitidos o erróneos se perdían o transferían sin transparencia (RES/174/2016). & \textbf{RLSE Art. 182:} CNE emite disposiciones para liquidación, cancelación y transacciones. & Acumulación de pasivos ocultos y distorsión de la contabilidad nacional de energías limpias. & Crear la \textbf{Reserva Estratégica de CELs} administrada por CNE para compensar déficits sistémicos justificados. \\ \hline
\textbf{Desconexión con Justicia Energética:} Emisión de CELs ciega al impacto social. & \textbf{LSE Art. 12 Fracc. VII:} Planeación con políticas de Justicia Energética y Sostenibilidad. & Otorgar CELs a proyectos con conflictos sociales activos, violando el principio de Sostenibilidad. & Vincular la emisión de CELs a la validación de la \textbf{Manifestación de Impacto Social} (RLSE Art. 42). \\ \hline
\bottomrule
\end{xltabular}
\end{tabladoradoLargo}


\subsection{C. Desarrollo Analítico: Hacia el Sistema de Monitoreo Soberano}
\nopagebreak

\subsubsection{1. Diagnóstico de la Situación Actual (Transición)}
\nopagebreak

El sistema heredado (S-CEL de la CRE) funcionaba como una cámara de compensación financiera. El nuevo sistema bajo la CNE debe funcionar como una herramienta de política pública para la \textbf{Descarbonización}.


\textbf{Datos Críticos para la CNE:}


\begin{itemize}
\item \textbf{Validación de Titularidad:} Migración de cuentas de la CRE a la CNE sin pérdida de datos históricos.
\item \textbf{Interconexión Legada:} Identificación precisa de capacidad excluida de contratos legados que ahora generan CELs bajo LSE.
\item \textbf{Sincronización CENACE-CNE:} Eliminación de la ventana de latencia de 10 días mediante interconexión directa de bases de datos.

\end{itemize}
\subsubsection{2. Estado Objetivo (Arquitectura LSE 2025)}
\nopagebreak

El modelo objetivo se alinea con el \textbf{Art. 147 de la LSE \textbf{, donde la CNE no solo \textquotedbl{}cuenta papeles\textquotedbl{}, sino que asegura la }Descarbonización Efectiva}:


\begin{itemize}
\item \textbf{Monitoreo Predictivo:} Uso de IA sobre los datos de \textit{Planeación Vinculante} (PLADESE) para prever la disponibilidad de CELs a 3 años (requisito RLSE Transitorio Décimo Sexto).
\item \textbf{Trazabilidad Total:} Desde el megawatt generado hasta el cumplimiento de la obligación del Centro de Carga.
\item \textbf{Auditoría de Eficiencia:} Integración de la metodología de Energía Libre de Combustible (ELC) directamente en el cálculo de emisión, verificando eficiencia térmica real.

\end{itemize}
\vspace{0.4cm}
\Needspace{25\baselineskip}
\subsubsection{3. Tabla Comparativa: Modelo Anterior vs. Modelo LSE 2025}
\nopagebreak

\vspace{0.3cm}
\Needspace{15\baselineskip}
\begin{tablaParadigmaV2}{Análisis Comparativo}{Modelo Anterior (Mercado)}{Modelo Objetivo (Soberanía/CNE)}
\tpRow{\textbf{Gestión de Datos}}{Reporte mensual estático (batch).}{\textbf{Gestión de Datos}}{Flujo continuo de datos CENACE → CNE.}
\tpRow{\textbf{Validación}}{Documental y de buena fe.}{\textbf{Validación}}{Técnica, basada en eficiencia real y cumplimiento social.}
\tpRow{\textbf{Flexibilidad}}{Basada en precio (\textgreater{}60 UDIs).}{\textbf{Flexibilidad}}{Basada en \textbf{Seguridad Energética} y disponibilidad física real.}
\tpRow{\textbf{Vigencia}}{Indefinida hasta consumo.}{\textbf{Vigencia}}{Condicionada a la actualización del Sistema Electrónico (RLSE Transitorio Sexto).}
\end{tablaParadigmaV2}


\subsection{D. Propuestas de Ajuste Normativo (Implementation Ready)}
\nopagebreak

\subsubsection{Instrumento A: Disposiciones Administrativas de Carácter General (DACG) para el Registro de CELs (Nuevo Art. 182 RLSE)}
\nopagebreak

\textbf{Propuesta de redacción para CNE:}


\vspace{0.2cm}
\Needspace{8\baselineskip}
\begin{glowBox}[gobmxDorado]{Nota Informativa}
\textquotedbl{}La Comisión Nacional de Energía, en coordinación con la Secretaría, establece que para el otorgamiento de CELs, los Generadores deberán acreditar no solo la inyección de energía limpia, sino el cumplimiento de los atributos de \textbf{Sostenibilidad y Justicia Energética} conforme al Art. 12 de la LSE. El sistema alertará automáticamente cuando el balance oferta-demanda proyectado a 6 meses sea inferior al umbral de seguridad definido por la Planeación Vinculante.\textquotedbl{}
\end{glowBox}

\subsubsection{Instrumento B: Mecanismo de Transición de Saldos (Transitorio Sexto RLSE)}
\nopagebreak

\textbf{Estrategia de Migración:}


\vspace{0.2cm}
\Needspace{8\baselineskip}
\begin{glowBox}[gobmxDorado]{Nota Informativa}
\textquotedbl{}Conforme al Transitorio Sexto del RLSE, los CELs emitidos anteriormente conservan validez. Sin embargo, su liquidación en el nuevo sistema estará condicionada a la \textbf{Validación de No-Duplicidad} y a la actualización del Registro de Certificados. Se establece una ventana de 180 días para la conciliación de saldos históricos entre la extinta CRE y la nueva CNE.\textquotedbl{}
\end{glowBox}


\subsection{E. Hoja de Ruta de Transitorios y Migración}
\nopagebreak

Para asegurar la continuidad operativa sin violar la nueva LSE:


\begin{enumerate}
\item \textbf{Fase 1 (Días 1-60):} CNE asume el control del S-CEL existente. Se suspende la emisión de nuevos criterios de flexibilidad financiera (A/013/2019 deja de aplicar por contravención a la LSE en materia de rectoría del Estado).
\item \textbf{Fase 2 (Días 61-120):} SENER publica los nuevos Requisitos de Energía Limpia 2025-2028 (Mandato RLSE Transitorio Décimo Sexto). CNE ajusta los algoritmos de validación.
\item \textbf{Fase 3 (Día 121+): \textbf{ Lanzamiento del }Módulo de Disponibilidad Real}, integrando proyecciones del PLADESE. Inicio de auditorías a centrales con \textquotedbl{}valores estimados\textquotedbl{} recurrentes.
\end{enumerate}
4.


\phantomsection
\addcontentsline{toc}{section}{Mecanismos de Flexibilidad y Diferimiento (Propuesta Regulatoria LSE 2025)}
\section*{Mecanismos de Flexibilidad y Diferimiento (Propuesta Regulatoria LSE 2025)}
\nopagebreak

\textbf{Autoridad Responsable:} Comisión Nacional de Energía (CNE) en coordinación con SENER

\textbf{Marco Jurídico Base:} Ley del Sector Eléctrico (2025), Artículos 14 (Confiabilidad), 20 (Fuerza Mayor) y 33 (Justicia Energética).


\subsection{A. Diagnóstico y Evolución Normativa (Tabla Comparativa)}
\nopagebreak

Ante la ausencia de un artículo explícito de \textquotedbl{}diferimiento\textquotedbl{} en la LSE 2025, se propone un mecanismo regulatorio basado en las facultades de emergencia y planeación de la CNE y SENER.


\vspace{0.3cm}
\Needspace{15\baselineskip}
\begin{tabladoradoLargo}
\renewcommand{\tabularxcolumn}[1]{p{#1}}
\begin{xltabular}{\textwidth}{|X|X|X|X|}
\caption{Matriz de Datos} \\
\toprule
\rowcolor{gobmxDorado} \encabezadotablaDorada{Variable} & \encabezadotablaDorada{Marco Anterior (LIE / CRE)} & \encabezadotablaDorada{Nuevo Marco (LSE 2025 / CNE)} & \encabezadotablaDorada{Justificación del Cambio} \\
\midrule
\endhead
\textbf{Mecanismo Base} & Lineamiento 25 (2014): Diferimiento del 25\% con costo financiero (+5\% anual). & \textbf{Mecanismo de Flexibilidad Soberana (Propuesta Reg.):} Diferimiento excepcional por causas operativas o de fuerza mayor (Art. 20 LSE). & Se elimina el diferimiento financiero especulativo. La flexibilidad ahora es una herramienta de seguridad operativa, no financiera. \\ \hline
\textbf{Gatillo de Activación} & \textbf{Acuerdo A/013/2019:} Automático si oferta CEL \textless{} 70\% demanda O Precio Implícito \textgreater{} 60 UDIS. & \textbf{Declaratoria de Estado Operativo (CENACE/CNE):} Activación por Alerta/Emergencia (Art. 14 LSE) o Fuerza Mayor certificada. & Se recupera la soberanía del control: el mercado no se \textquotedbl{}auto-regula\textquotedbl{} por precios, sino que la autoridad declara la excepción para proteger el suministro. \\ \hline
\textbf{Plataforma} & S-CEL (Aislado del control operativo). & \textbf{SNIEr 2.0 (Integrado):} Interoperabilidad total con CENACE. & Soberanía de datos: Validar que el CEL corresponda a energía inyectada y no a \textquotedbl{}papel\textquotedbl{} financiero. \\ \hline
\textbf{Sujetos} & Trato igualitario ciego (Suministradores y Usuarios Calificados). & \textbf{Diferenciación Estratégica (Art. 33 LSE):} Prioridad a Justicia Energética y Proyectos Estratégicos. & Reconocimiento del carácter de servicio público y estratégico de la electricidad (LSE Art. 4 y 33). \\ \hline
\bottomrule
\end{xltabular}
\end{tabladoradoLargo}

\subsection{B. Propuesta Técnica: Mecanismo de Flexibilidad Soberana}
\nopagebreak

Bajo la LSE 2025, el cumplimiento de CELs es un compromiso de descarbonización física. La flexibilidad no es un derecho automático del mercado, sino una medida de \textbf{salvaguarda operativa}.


\subsubsection{7.1 Criterios de Clasificación para Diferimiento}
\nopagebreak

La CNE clasificará a los Sujetos Obligados para la aplicación de este mecanismo, basándose en la \textbf{Justicia Energética} y la planeación estratégica (PLADESE):


\begin{enumerate}
\item \textbf{Entidades de Servicio Público (Suministrador de Servicios Básicos - CFE): \textbf{ Flexibilidad alineada a los tiempos de entrada en operación de }Proyectos Estratégicos} definidos por SENER (Art. 13 LSE).
\item \textbf{Usuarios Calificados (Privados):}
\end{enumerate}
\begin{itemize}
\item \textit{Industria Estratégica:} Esquemas flexibles para polos de desarrollo (Art. 8-II LSE), condicionados a inversión en infraestructura local.
\item \textit{Consumidores Generales:} Cumplimiento estricto. La \textquotedbl{}banca\textquotedbl{} de CELs solo aplica si existe imposibilidad material de adquisición validada por CNE.

\end{itemize}
\subsubsection{7.2 Reglas de Diferimiento y \textquotedbl{}Banca\textquotedbl{} de CELs (Protocolo Propuesto)}
\nopagebreak

Se sustituye el modelo de penalización financiera simple por un modelo de \textbf{Validación de Causa}:


\begin{itemize}
\item \textbf{Tope de Diferimiento:} Máximo 30\% de la obligación anual (Propuesta regulatoria para mantener disciplina).
\item \textbf{Condición de Activación (Art. 20 LSE):} Solo aplicable si el CENACE certifica:
\item Estados de Emergencia o Alerta en la Red Nacional de Transmisión.
\item Causa de Fuerza Mayor que impidió la inyección o retiro de energía limpia contratada.
\item \textbf{Costo de Regularización:}
\item \textit{Opción A (Sanción):} Si no hay causa justificada, aplica sanción directa (Multa) sin diferimiento.
\item \textit{Opción B (Regularización):} Con causa justificada, se permite diferir con un incremento del \textbf{10\%} de la obligación (desincentivo mayor al 5\% anterior) y compromiso de inversión en Generación Distribuida (Art. 48 LSE).

\end{itemize}
\subsection{C. Variables Críticas e Interoperabilidad (Requisitos de Datos)}
\nopagebreak

Para que la CNE autorice el diferimiento bajo el principio de Soberanía de Datos, se requiere:


\vspace{0.3cm}
\Needspace{15\baselineskip}
\begin{tablaParadigmaV2}{Análisis Comparativo}{Fuente de Validación}{Objetivo Operativo}
\tpRow{\textbf{Consumo Real (MWh)}}{Medición Fiscal (CENACE)}{\textbf{Consumo Real (MWh)}}{Evitar sub-reporte de obligaciones (Art. 25 LSE).}
\tpRow{\textbf{Declaratoria Operativa}}{CENACE (Protocolo de Emergencia)}{\textbf{Declaratoria Operativa}}{Verificar que existió una restricción física real en la red (Art. 14 LSE).}
\tpRow{\textbf{Causa de Fuerza Mayor}}{Bitácoras de Operación (CENACE)}{\textbf{Causa de Fuerza Mayor}}{Validar eventos fortuitos (Art. 20 LSE) vs negligencia administrativa.}
\tpRow{\textbf{Evaluación Impacto Social}}{Resolución SENER (LSE Art. 19-III)}{\textbf{Evaluación Impacto Social}}{\textbf{Nuevo Requisito:} No se otorga flexibilidad si el proyecto asociado tiene conflictos sociales no resueltos (Art. 33 LSE).}
\end{tablaParadigmaV2}

\subsection{D. Artículos Transitorios (Estrategia de Migración)}
\nopagebreak

Para asegurar la continuidad operativa y evitar litigios durante el cambio de LIE a LSE:


\textbf{PRIMERO. \textbf{ Se reconocen los saldos diferidos bajo el marco anterior (Lineamiento 25, 2014) correspondientes a los periodos 2023-2024. Estos deberán liquidarse en un plazo improrrogable de }18 meses} bajo las reglas originales (5\% interés). Al término, cualquier saldo pendiente se convertirá en crédito fiscal firme.


\textbf{SEGUNDO. \textbf{ Los procedimientos administrativos de sanción iniciados por la CRE antes de la entrada en vigor de la Ley de la CNE serán resueltos por la CNE aplicando el principio de \textit{in dubio pro reo} para cerrar expedientes rezagados, condicionando el cierre a la firma de un }Convenio de Regularización} con inversión en infraestructura de beneficio social (electrificación rural/urbana).


\textbf{TERCERO. \textbf{ Las solicitudes de diferimiento para el ciclo 2025-2026 se regirán exclusivamente por el nuevo }Protocolo de Emergencia y Confiabilidad} que emita la CNE, dejando sin efectos el Acuerdo A/013/2019.



\textbf{Estatus:} PROPUESTA DE AJUSTE REGULATORIO (Requiere Visto Bueno de SENER para emisión de DACGs).

\textbf{Fecha:} Enero 2026


\phantomsection
\addcontentsline{toc}{section}{Entidades Voluntarias y Cancelación Voluntaria de Certificados (Versión Final LSE 2025)}
\section*{Entidades Voluntarias y Cancelación Voluntaria de Certificados (Versión Final LSE 2025)}
\nopagebreak

\textbf{Marco Jurídico Base:} - Ley del Sector Eléctrico 2025 (LSE)


\begin{itemize}
\item Ley de la Comisión Nacional de Energía (Ley CNE)
\item Reglamento de la LSE 2025
\item Ley de Planeación y Transición Energética (LPTE)


\end{itemize}
\subsection{1. Auditoría de la \textquotedbl{}Foto Actual\textquotedbl{} (Marco Derogado RES/174/2016)}
\nopagebreak

\textbf{Diagnóstico:} El marco anterior (CRE) diseñó la participación voluntaria como un mecanismo financiero desconectado de la política climática real.


\begin{itemize}
\item \textbf{Falla de Definición (Disp. 7 RES/174/2016):} Se definía a la Entidad Voluntaria por la negativa (\textquotedbl{}quien no está obligado\textquotedbl{}), sin exigir un propósito de descarbonización.
\item \textbf{Cancelación \textquotedbl{}Ciega\textquotedbl{} (Disp. 53 RES/174/2016): \textbf{ Permitía la quema de certificados sin validar la }adicionalidad} (impacto real en reducción de emisiones), facilitando el \textit{greenwashing}.
\item \textbf{Opacidad de Datos (Disp. 60 RES/174/2016):} La publicación de datos agregados impedía auditar quién cancelaba y por qué, anulando el reconocimiento reputacional público.


\end{itemize}
\subsection{2. Reingeniería Legal (Nuevo Marco LSE 2025)}
\nopagebreak

Se reestructura el módulo para alinear la participación voluntaria con los objetivos de \textbf{Soberanía y Justicia Energética}.


\vspace{0.4cm}
\Needspace{25\baselineskip}
\subsubsection{A. Tabla Comparativa de Atribuciones}
\nopagebreak

\vspace{0.3cm}
\Needspace{15\baselineskip}
\begin{tablaParadigmaV2}{Análisis Comparativo}{Marco Anterior (CRE)}{Nuevo Marco (CNE / SENER - LSE 2025)}
\tpRow{\textbf{Autoridad}}{Comisión Reguladora de Energía (CRE)}{\textbf{Autoridad}}{\textbf{Comisión Nacional de Energía (CNE) \textbf{ (Operación) y }SENER} (Política)}
\tpRow{\textbf{Fundamento}}{LIE 2014 (Derogada)}{\textbf{Fundamento}}{\textbf{Art. 147, 149 LSE y Art. 8 Ley CNE}}
\tpRow{\textbf{Registro}}{S-CEL (Gestión administrativa)}{\textbf{Registro}}{\textbf{Registro de Certificados (Art. 149 LSE / Art. 182 RLSE):} Sistema de trazabilidad con interoperabilidad al SNIE.}
\tpRow{\textbf{Cancelación}}{Trámite automático administrativo.}{\textbf{Cancelación}}{\textbf{Validación de Impacto (Art. 181 RLSE):} Proceso que valida la contribución a metas de descarbonización.}
\end{tablaParadigmaV2}

\subsubsection{B. Matriz de Variables Críticas (Nuevos Requisitos de Datos)}
\nopagebreak

Para cumplir con el mandato de transparencia del \textbf{Art. 8 de la Ley CNE}, el nuevo sistema requerirá obligatoriamente los siguientes campos para cualquier Entidad Voluntaria:


\vspace{0.3cm}
\Needspace{15\baselineskip}
\begin{tablaParadigmaV2}{Análisis Comparativo}{Descripción Técnica}{Justificación Jurídica (LSE 2025)}
\tpRow{\textbf{Huella de Carbono (tCO2e)}}{Toneladas de CO2 evitadas por el paquete de CELs a cancelar.}{\textbf{Huella de Carbono (tCO2e)}}{Vinculación con control de emisiones (Art. 9 LPTE / Art. 150 LSE).}
\tpRow{\textbf{Beneficiario Final}}{RFC/Nombre de la entidad que reclama el beneficio ambiental (evita reventa).}{\textbf{Beneficiario Final}}{Principios de Integridad y Rendición de Cuentas (Art. 3 Ley CNE).}
\tpRow{\textbf{Motivo de Cancelación}}{Catálogo: \textit{NetZero, ESG, Contractual, Filantrópico}.}{\textbf{Motivo de Cancelación}}{Insumo para la Planeación Nacional (PLADESE - SENER).}
\end{tablaParadigmaV2}


\subsection{3. Propuesta Normativa (Texto para Nuevas DACG - CNE 2026)}
\nopagebreak

Se propone la inclusión del siguiente capítulo en las próximas Disposiciones Administrativas de Carácter General (DACG) que emitirá la CNE, fundamentado en su facultad de emitir regulación operativa (Art. 8 Fracción V, LCNE):


\vspace{0.2cm}
\Needspace{8\baselineskip}
\begin{glowBox}[gobmxDorado]{Nota Informativa}
\textbf{CAPÍTULO V. DE LA PARTICIPACIÓN VOLUNTARIA Y JUSTICIA ENERGÉTICA}
\end{glowBox}
\textgreater{}

\vspace{0.2cm}
\Needspace{8\baselineskip}
\begin{glowBox}[gobmxDorado]{Nota Informativa}
\textbf{Disposición 24. Registro Simplificado.}
\end{glowBox}
\vspace{0.2cm}
\Needspace{8\baselineskip}
\begin{glowBox}[gobmxDorado]{Nota Informativa}
Las Entidades Voluntarias podrán inscribirse en el Registro mediante un procedimiento digital expedito, declarando bajo protesta de decir verdad sus objetivos de descarbonización y aceptando la publicidad de sus acciones de cancelación.
\end{glowBox}
\textgreater{}

\vspace{0.2cm}
\Needspace{8\baselineskip}
\begin{glowBox}[gobmxDorado]{Nota Informativa}
\textbf{Disposición 25. Cancelación con Trazabilidad (Certificado de Retiro).}
\end{glowBox}
\vspace{0.2cm}
\Needspace{8\baselineskip}
\begin{glowBox}[gobmxDorado]{Nota Informativa}
I. La cancelación voluntaria de CELs es irrevocable y extingue el instrumento de manera definitiva del Registro (Art. 149 LSE).
\end{glowBox}
\vspace{0.2cm}
\Needspace{8\baselineskip}
\begin{glowBox}[gobmxDorado]{Nota Informativa}
II. Al ejecutarse la cancelación, la CNE emitirá un \textquotedbl{}Comprobante de Retiro\textquotedbl{} digital que validará la contribución del titular a las Metas de Energías Limpias (Art. 181 RLSE).
\end{glowBox}
\vspace{0.2cm}
\Needspace{8\baselineskip}
\begin{glowBox}[gobmxDorado]{Nota Informativa}
III. El sistema validará automáticamente que los CELs a cancelar no hayan sido utilizados previamente en otros mercados de carbono, garantizando la integridad de la contabilidad nacional.
\end{glowBox}
\textgreater{}

\vspace{0.2cm}
\Needspace{8\baselineskip}
\begin{glowBox}[gobmxDorado]{Nota Informativa}
\textbf{Disposición 26. Reporte de Transparencia.}
\end{glowBox}
\vspace{0.2cm}
\Needspace{8\baselineskip}
\begin{glowBox}[gobmxDorado]{Nota Informativa}
La CNE publicará mensualmente en su portal oficial el \textquotedbl{}Listado de Contribuyentes a la Transición\textquotedbl{}, detallando las Entidades Voluntarias y los volúmenes cancelados por tecnología y región, conforme al principio de máxima publicidad y rendición de cuentas (Art. 3 Ley CNE).
\end{glowBox}


\subsection{4. Artículos Transitorios (Estrategia de Migración)}
\nopagebreak

Para asegurar la continuidad operativa durante el cambio de LIE a LSE (conforme al Transitorio Sexto del Reglamento LSE):


\begin{enumerate}
\item \textbf{Reconocimiento de Saldos:} Los CELs adquiridos por voluntarios bajo el marco anterior (RES/174/2016) mantienen su validez jurídica.
\item \textbf{Ventana de Migración (72 horas):} Se establecerá un periodo de suspensión técnica para trasladar las cuentas de voluntarios del S-CEL legado al nuevo Sistema CNE.
\item \textbf{Actualización de Perfiles:} Al primer ingreso en la nueva plataforma, las Entidades Voluntarias deberán actualizar su perfil fiscal y aceptar los nuevos Términos y Condiciones de Trazabilidad.

\end{enumerate}
\portadaseccion{BLOQUE IV}{Mercado de Certificados de Energía Limpia y transacciones bilaterales}{}
\clearpage
\phantomsection
\addcontentsline{toc}{section}{Mercado de Certificados de Energía Limpia y transacciones bilaterales}
\section*{Mercado de Certificados de Energía Limpia y transacciones bilaterales}
\nopagebreak

\textbf{Alcance del punto:} Análisis de las áreas de oportunidad relacionadas con el mercado de CEL, mecanismos de transacción bilateral, transparencia de precios, gestión de riesgos crediticios y eficiencia operativa

\textbf{Instrumentos jurídicos revisados:} Reglamento de la Ley del Sector Eléctrico (LSE 2025), RES/174/2016 (Referencia Histórica), Bases del Mercado Eléctrico, Manual de Transacciones Bilaterales CENACE


\vspace{0.4cm}
\Needspace{25\baselineskip}
\subsection{A. Fuentes de información del S-CEL para el Punto 1}
\nopagebreak

\vspace{0.3cm}
\Needspace{15\baselineskip}
\begin{tabladoradoLargo}
\renewcommand{\tabularxcolumn}[1]{p{#1}}
\begin{xltabular}{\textwidth}{|X|X|X|X|}
\caption{Matriz de Datos} \\
\toprule
\rowcolor{gobmxDorado} \encabezadotablaDorada{Actor / Fuente} & \encabezadotablaDorada{Instrumento legal} & \encabezadotablaDorada{Artículo / Numeral} & \encabezadotablaDorada{Cita textual} \\
\midrule
\endhead
CENACE (Operación del Mercado) & Bases del Mercado Eléctrico & Base 3.1 & \textquotedbl{}Disposiciones generales\textquotedbl{} para registro y acreditación de Participantes del Mercado \\ \hline
Participantes del Mercado & Bases del Mercado Eléctrico & Base 1.3 & \textquotedbl{}Descripción del Mercado Eléctrico Mayorista\textquotedbl{} incluyendo mecanismos de transacción \\ \hline
Transacciones Bilaterales CEL & RES/174/2016 & Disposición 44 & \textquotedbl{}Cuando se trate de transacciones bilaterales, la parte que transfiere los CEL deberá comunicar de manera electrónica al S-CEL el número de CEL que serán transferidos y el nombre de usuario del destinatario\textquotedbl{} \\ \hline
CENACE (Mercado Centralizado) & DACG S-CEL (RES/174/2016) & Disp. 39 & \textquotedbl{}El Cenace operará un mercado de CEL conforme a lo señalado en las Bases del Mercado Eléctrico. De existir las condiciones para ello, el Cenace podrá proponer a la Comisión la operación con una frecuencia mayor a la establecida en las Bases del Mercado Eléctrico\textquotedbl{} \\ \hline
CENACE (Reporte de Operaciones) & DACG S-CEL (RES/174/2016) & Disp. 42 & \textquotedbl{}El día hábil siguiente de finalizada la operación del mercado de CEL, el Cenace entregará a la Comisión, entre otros, un reporte electrónico de las operaciones de dicho mercado, mismo que deberá incluir los siguientes campos: i) el nombre del Participante del Sistema; ii) el número de CEL que deberán ser deducidos... iv) Para fines estadísticos y de transparencia reportará también el precio de los CEL resultante al cierre de cada ejecución del mercado de CEL\textquotedbl{} \\ \hline
Participantes (Transacciones Bilaterales) & DACG S-CEL (RES/174/2016) & Disp. 44 & \textquotedbl{}Cuando se trate de transacciones bilaterales, la parte que transfiere los CEL deberá comunicar de manera electrónica al Sistema el número de CEL que serán transferidos y el nombre de usuario en el Sistema del destinatario, quien a su vez deberá manifestar por el mismo medio que acepta la transferencia\textquotedbl{} \\ \hline
CNE (Validación de Titularidad) & DACG S-CEL (RES/174/2016) & Disp. 41 & \textquotedbl{}Para la operación del mercado de CEL, tres días antes al inicio del periodo para aceptar ofertas, la Comisión entregará al Cenace un reporte electrónico actualizado del número de CEL que cada Participante del Mercado tiene en su posesión... El Cenace utilizará esta información para evitar que los Participantes del Mercado realicen ofertas de venta de CEL por cantidades mayores a las que efectivamente poseen\textquotedbl{} \\ \hline
Boletín Electrónico & DACG S-CEL (RES/174/2016) & Disp. 46 & \textquotedbl{}Los Participantes del Sistema que cuenten con CEL que no estén comprometidos podrán ofertarlos en el Boletín Electrónico que se ubicará en el portal electrónico del Sistema. Asimismo, los Participantes del Sistema interesados en adquirir CEL podrán publicar su interés en el Boletín Electrónico\textquotedbl{} \\ \hline
Contratos de Cobertura & DACG S-CEL (RES/174/2016) & Disp. 36 & \textquotedbl{}El Cenace notificará a la Comisión el fallo de adjudicación de cada uno de los Contratos de Cobertura Eléctrica en las subastas de largo plazo, indicando los términos de los correspondientes a CEL\textquotedbl{} \\ \hline
Modalidades de Transacción & DACG S-CEL (RES/174/2016) & Disp. 35 & \textquotedbl{}La compraventa de los CEL se podrán pactar en el mercado de CEL, las subastas de largo plazo organizadas por el Cenace o a través de las transacciones bilaterales\textquotedbl{} \\ \hline
\textbf{CNE (Nuevas Atribuciones)} & \textbf{Reglamento LSE 2025} & \textbf{Arts. 3, 4, 14, 19} & \textbf{Faculta a la CNE para emitir DACG, utilizar medios electrónicos y establecer criterios para Productos Asociados (CEL).} \\ \hline
\bottomrule
\end{xltabular}
\end{tabladoradoLargo}

\vspace{0.4cm}
\Needspace{25\baselineskip}
\subsection{B. Matriz de validación jurídica}
\nopagebreak

\vspace{0.3cm}
\Needspace{15\baselineskip}
\begin{tabladoradoLargo}
\renewcommand{\tabularxcolumn}[1]{p{#1}}
\begin{xltabular}{\textwidth}{|X|X|X|X|}
\caption{Matriz de Datos} \\
\toprule
\rowcolor{gobmxDorado} \encabezadotablaDorada{Hallazgo o limitación} & \encabezadotablaDorada{Fundamento jurídico (APA + cita textual)} & \encabezadotablaDorada{Riesgo (jurídico u operativo)} & \encabezadotablaDorada{Ajuste propuesto (alineado al marco vigente)} \\
\midrule
\endhead
Dependencia excesiva del mercado centralizado operado por CENACE & Comisión Reguladora de Energía. (2016). \textit{DACG del Sistema de CEL}. Disposición 39: \textquotedbl{}El Cenace operará un mercado de CEL conforme a lo señalado en las Bases del Mercado Eléctrico\textquotedbl{} & Operativo: Concentración de riesgo en una sola plataforma de transacción, vulnerabilidad ante fallas técnicas o operativas & Diversificar mecanismos de transacción con plataformas complementarias y sistemas de respaldo \\ \hline
Limitaciones del Boletín Electrónico como mecanismo de descubrimiento de precios & Comisión Reguladora de Energía. (2016). \textit{DACG del Sistema de CEL}. Disposición 46: \textquotedbl{}Los Participantes del Sistema que cuenten con CEL... podrán ofertarlos en el Boletín Electrónico... los interesados en adquirir CEL podrán publicar su interés\textquotedbl{} & Mercado: Ineficiencia en formación de precios, limitada liquidez y falta de matching automático & Mejorar funcionalidades del Boletín con matching automático, transparencia de precios y ejecución directa \\ \hline
Proceso de transacciones bilaterales con alta fricción operativa & Comisión Reguladora de Energía. (2016). \textit{DACG del Sistema de CEL}. Disposición 44: \textquotedbl{}la parte que transfiere los CEL deberá comunicar... el destinatario... deberá manifestar... que acepta la transferencia\textquotedbl{} & Eficiencia: Proceso manual que genera retrasos, errores y costos de transacción elevados & Automatizar proceso bilateral con \textbf{Smart Contracts} (respaldados en Art. 4 Reglamento LSE) y validación en tiempo real \\ \hline
Falta de transparencia en precios de transacciones bilaterales & Comisión Reguladora de Energía. (2016). \textit{DACG del Sistema de CEL}. Disposición 60: \textquotedbl{}los precios de los CEL en el mercado de CEL, transacciones bilaterales\textquotedbl{} sin especificar metodología de cálculo o reporte & Transparencia: Opacidad en formación de precios del mercado bilateral que afecta eficiencia del mercado & Establecer metodología específica y obligatoria para reporte de precios bilaterales \\ \hline
Ausencia de mecanismos de garantía para transacciones bilaterales & Comisión Reguladora de Energía. (2016). \textit{DACG del Sistema de CEL}. Disposición 44: \textquotedbl{}La Comisión sólo ejecutará la o las transferencias solicitadas y será ajena al pago de las contraprestaciones\textquotedbl{} & Riesgo crediticio: Falta de garantías que aseguren cumplimiento de pagos, exposición a incumplimientos & Implementar sistema de garantías y clearing para transacciones bilaterales \\ \hline
\bottomrule
\end{xltabular}
\end{tabladoradoLargo}

\subsection{C. Desarrollo analítico}
\nopagebreak

\subsubsection{Diagnóstico de la situación actual}
\nopagebreak

El mercado de CEL presenta una estructura híbrida con múltiples mecanismos de transacción que enfrentan limitaciones operativas significativas:


\textbf{Fortalezas identificadas:}


\begin{itemize}
\item Marco regulatorio establecido con múltiples canales de transacción
\item Mercado centralizado operado por CENACE con reportes estructurados
\item Mecanismo de validación de titularidad que previene overselling
\item Trazabilidad completa de transacciones desde origen hasta liquidación

\end{itemize}
\textbf{Limitaciones detectadas:}


\begin{itemize}
\item Fragmentación del mercado en múltiples plataformas con diferentes niveles de eficiencia
\item Dependencia excesiva del mercado centralizado sin alternativas robustas
\item Proceso manual de transacciones bilaterales con alta fricción operativa
\item Boletín Electrónico básico sin funcionalidades avanzadas de trading

\end{itemize}
\textbf{Problemáticas operativas identificadas:}


\begin{itemize}
\item Falta de transparencia en la formación de precios, especialmente en transacciones bilaterales
\item Ausencia de mecanismos de garantía que reduzcan el riesgo crediticio
\item Limitada liquidez debido a barreras operativas y falta de market makers
\item Dependencia de procesos manuales que generan retrasos y errores

\end{itemize}
\textbf{Datos críticos identificados:}


\begin{itemize}
\item Concentración de volúmenes en el mercado centralizado
\item Spreads bid-ask elevados por falta de liquidez
\item Tiempos de ejecución prolongados en transacciones bilaterales
\item Limitada participación de pequeños y medianos actores

\end{itemize}
\subsubsection{Estado objetivo}
\nopagebreak

El modelo objetivo debe crear un ecosistema de transacciones integrado que garantice:


\begin{enumerate}
\item \textbf{Eficiencia Operativa} mediante automatización de procesos y reducción de fricciones
\item \textbf{Transparencia Total} con visibilidad completa de precios y volúmenes
\item \textbf{Gestión de Riesgos} a través de sistemas de garantía y clearing
\item \textbf{Liquidez Mejorada} con integración de plataformas y market making
\item \textbf{Accesibilidad Universal} con interfaces modernas para todos los participantes

\end{enumerate}
\vspace{0.4cm}
\Needspace{25\baselineskip}
\subsubsection{Tabla comparativa: modelo actual vs modelo objetivo}
\nopagebreak

\vspace{0.3cm}
\Needspace{15\baselineskip}
\begin{tablaParadigmaV2}{Análisis Comparativo}{Modelo Actual}{Modelo Objetivo}
\tpRow{Estructura de Mercado}{Fragmentada: CENACE + bilaterales + boletín}{Estructura de Mercado}{Integrada con múltiples canales interoperables}
\tpRow{Proceso de Transacción}{Manual con verificación bidireccional}{Proceso de Transacción}{Automatizado con \textit{Smart Contracts} y validación en tiempo real}
\tpRow{Transparencia de Precios}{Limitada, especialmente en bilaterales}{Transparencia de Precios}{Total con metodología estándar para todos los mecanismos}
\tpRow{Gestión de Riesgos}{Ausente, CNE ajena a contraprestaciones}{Gestión de Riesgos}{Sistema integral de garantías y clearing}
\tpRow{Liquidez}{Limitada por barreras operativas}{Liquidez}{Mejorada con market makers y agregación de liquidez}
\tpRow{Acceso al Mercado}{Requiere contratos previos con CENACE}{Acceso al Mercado}{Acceso directo con diferentes niveles de participación}
\end{tablaParadigmaV2}

\subsubsection{Arquitectura del sistema}
\nopagebreak

El sistema objetivo debe integrar:


\begin{enumerate}
\item \textbf{Hub Central de Transacciones} con plataforma unificada que conecte todos los mecanismos
\item \textbf{Motor de Matching Inteligente} con sistema automatizado de emparejamiento
\item \textbf{Sistema de Clearing y Garantías} con cámara de compensación para gestión de riesgos
\item \textbf{Módulo de Transparencia} con dashboard en tiempo real
\item \textbf{API de Integración} con interfaces estándar para sistemas externos

\end{enumerate}
\subsubsection{Reingeniería de procesos}
\nopagebreak

\textbf{Proceso 1: Transacciones Integradas}


\begin{enumerate}
\item Registro de intención en plataforma unificada
\item Validación automática de titularidad y disponibilidad
\item Matching inteligente considerando precio, volumen y plazo
\item Ejecución automática con \textit{Smart Contracts}
\item Clearing y liquidación con gestión de garantías
\item Confirmación y registro con trazabilidad completa
\item Reporte de transparencia con publicación automática

\end{enumerate}
\textbf{Proceso 2: Market Making Dinámico}


\begin{enumerate}
\item Análisis continuo de liquidez por tecnología y plazo
\item Activación automática de market makers registrados
\item Cotización continua con spreads bid-ask en tiempo real
\item Incentivos dinámicos según condiciones de mercado
\item Monitoreo de performance y calidad de market making

\end{enumerate}
\textbf{Proceso 3: Gestión de Riesgos}


\begin{enumerate}
\item Evaluación automática de riesgo crediticio de participantes
\item Cálculo dinámico de garantías requeridas
\item Monitoreo en tiempo real de exposiciones
\item Ejecución automática de garantías en caso de incumplimiento
\item Liquidación neta multilateral para optimizar flujos

\end{enumerate}
\subsubsection{Beneficios esperados}
\nopagebreak

\begin{itemize}
\item \textbf{Reducción del 70\%} en costos de transacción
\item \textbf{Transparencia del 100\%} de precios y volúmenes en tiempo real
\item \textbf{Incremento del 150\%} en volúmenes transados
\item \textbf{Eliminación del 95\%} de incumplimientos por garantías
\item \textbf{Reducción del 40\%} en spreads bid-ask
\item \textbf{Incremento del 200\%} en participantes activos

\end{itemize}
\subsection{D. Propuestas de ajuste normativo}
\nopagebreak

\subsubsection{Instrumento A: Emisión de nuevas Disposiciones Administrativas de Carácter General (DACG) por la CNE}
\nopagebreak

\textit{(En sustitución del marco anterior de la CRE, con fundamento en el Art. 3 del Reglamento de la LSE 2025)}


\textbf{Disposición 39 Bis (Adición) - Hub Integrado de Transacciones:}

\textquotedbl{}La CNE operará un hub integrado de transacciones de CEL que incluirá: i) mercado centralizado operado por CENACE, ii) plataforma de transacciones bilaterales automatizadas, iii) sistema de market making, iv) cámara de compensación y garantías. Todos los mecanismos operarán con transparencia total y interoperabilidad completa.\textquotedbl{}

\textit{(Fundamento: \textbf{Art. 14 Reglamento LSE \textbf{ - Criterios para adquisición de Productos Asociados y }Art. 4} - Facultad para auxiliarse de herramientas tecnológicas)}


\textbf{Disposición 44 Bis (Adición) - Transacciones Automatizadas (Smart Contracts):}

\textquotedbl{}Las transacciones bilaterales podrán ejecutarse mediante \textbf{Smart Contracts} (programas tecnológicos de validación automatizada) que validen automáticamente: i) titularidad de CEL, ii) disponibilidad de garantías, iii) cumplimiento de términos contractuales. La ejecución será inmediata una vez cumplidas las condiciones programadas.\textquotedbl{}

\textit{(Fundamento: \textbf{Art. 4 Reglamento LSE} - \textquotedbl{}la Secretaría, el CENACE y la CNE, pueden auxiliarse de medios de comunicación electrónicos, de equipos, herramientas, sistemas y programas tecnológicos\textquotedbl{})}


\textbf{Disposición 46 Bis (Adición) - Boletín Electrónico Avanzado:}

\textquotedbl{}El Boletín Electrónico incluirá funcionalidades de: i) matching automático de ofertas y demandas, ii) cotización en tiempo real, iii) ejecución directa de transacciones, iv) integración con sistema de garantías. Los precios resultantes se publicarán en tiempo real.\textquotedbl{}


\textbf{Disposición 60 Bis (Adición) - Transparencia de Mercado:}

\textquotedbl{}La CNE publicará diariamente: i) precios promedio ponderados por mecanismo de transacción, ii) volúmenes transados por tecnología, iii) spreads bid-ask, iv) indicadores de liquidez, v) estadísticas de market making. La información será accesible vía API en tiempo real.\textquotedbl{}

\textit{(Fundamento: \textbf{Art. 6 Reglamento LSE} - Principios de Transparencia, Eficiencia y Competitividad)}


\subsubsection{Instrumento B: Manual de Operación del Hub de Transacciones}
\nopagebreak

\textbf{Sección propuesta 9.1: Criterios de Market Making}

\textquotedbl{}Los market makers deberán mantener: i) Spreads máximos del 5\% en condiciones normales, ii) Disponibilidad mínima del 80\% del tiempo de operación, iii) Volúmenes mínimos según categoría de participante, iv) Respuesta máxima de 5 segundos a solicitudes de cotización.\textquotedbl{}


\textbf{Sección propuesta 9.2: Sistema de Garantías}

\textquotedbl{}Las garantías se calcularán considerando: i) Riesgo crediticio del participante, ii) Volatilidad histórica de precios de CEL, iii) Volumen y plazo de la transacción, iv) Correlación con otros riesgos del participante. Las garantías se ajustarán diariamente.\textquotedbl{}


\subsection{E. Observaciones técnicas relevantes}
\nopagebreak

\subsubsection{Riesgos identificados}
\nopagebreak

\begin{itemize}
\item \textbf{Riesgo tecnológico:} Complejidad en la integración de múltiples sistemas y plataformas
\item \textbf{Riesgo operativo:} Posible resistencia de participantes acostumbrados a procesos manuales
\item \textbf{Riesgo de mercado:} Volatilidad inicial durante la transición al nuevo sistema

\end{itemize}
\subsubsection{Supuestos operativos}
\nopagebreak

\begin{itemize}
\item Disponibilidad de infraestructura tecnológica robusta para soportar trading en tiempo real
\item Capacidad de CENACE para integrar nuevas funcionalidades sin afectar operación actual
\item Participación activa de market makers para garantizar liquidez inicial

\end{itemize}
\subsubsection{Dependencias con otros bloques}
\nopagebreak

\begin{itemize}
\item \textbf{Bloque I:} Coordinación con sistemas de registro y administración de participantes
\item \textbf{Bloque II:} Integración con procesos de otorgamiento y trazabilidad de CEL
\item \textbf{Bloque V:} Articulación con mecanismos de formación de precios y señales ambientales


\end{itemize}
\textbf{Revisión técnica:} Área de Certificados de Energía Limpia

\textbf{Fecha de conclusión:} 12 de enero de 2026


\phantomsection
\addcontentsline{toc}{section}{Bolsa No Onerosa de Certificados de Energía Limpia}
\section*{Bolsa No Onerosa de Certificados de Energía Limpia}
\nopagebreak

\textbf{Alcance del punto:} Análisis de las áreas de oportunidad relacionadas con la Bolsa No Onerosa de CEL, distorsiones de mercado, criterios de asignación, transparencia y propuesta de eliminación del mecanismo

\textbf{Instrumentos jurídicos revisados:} Reglamento de la Ley de Planeación y Transición Energética (Art. 77 FOTEASE), Reglamento de la Ley del Sector Eléctrico 2025 (Arts. 6, 14), RES/174/2016 (Referencia Histórica), A/012/2020


\vspace{0.4cm}
\Needspace{25\baselineskip}
\subsection{A. Fuentes de información del S-CEL para el Punto 2}
\nopagebreak

\vspace{0.3cm}
\Needspace{15\baselineskip}
\begin{tabladoradoLargo}
\renewcommand{\tabularxcolumn}[1]{p{#1}}
\begin{xltabular}{\textwidth}{|X|X|X|X|}
\caption{Matriz de Datos} \\
\toprule
\rowcolor{gobmxDorado} \encabezadotablaDorada{Actor / Fuente} & \encabezadotablaDorada{Instrumento legal} & \encabezadotablaDorada{Artículo / Numeral} & \encabezadotablaDorada{Cita textual} \\
\midrule
\endhead
Cuenta CNE (Bolsa No Onerosa) & DACG S-CEL (RES/174/2016) & Disp. 34 & \textquotedbl{}Aquellos CEL que no fueron emitidos ni otorgados a pesar de que ello habría sido posible serán transferidos a la cuenta que para tal efecto administre la Comisión... la Comisión podrá asignarlos de manera no onerosa y proporcionalmente al consumo de todos los Participantes Obligados\textquotedbl{} \\ \hline
Generadores no inscritos & DACG S-CEL (RES/174/2016) & Disp. 34.1 & \textquotedbl{}El Generador Limpio o Suministrador que represente Generación Limpia Distribuida no se inscribió en el Sistema\textquotedbl{} \\ \hline
CENACE (ajustes y reliquidaciones) & DACG S-CEL (RES/174/2016) & Disp. 34.2 & \textquotedbl{}Hubo un ajuste por parte del Cenace en relación con la información que fue entregada a la Comisión y no fue registrada en el Sistema\textquotedbl{} \\ \hline
Generadores morosos & DACG S-CEL (RES/174/2016) & Disp. 34.3 & \textquotedbl{}El Generador Limpio o Suministrador que represente Generación Limpia Distribuida no cubrió el pago de derechos correspondiente a la cuota de dos años previos\textquotedbl{} \\ \hline
Cláusula abierta discrecional & DACG S-CEL (RES/174/2016) & Disp. 34.4 & \textquotedbl{}Aquellos otros que la Comisión considere justificados\textquotedbl{} \\ \hline
\textbf{CNE (Transparencia CEL)} & \textbf{Reglamento LPTE} & \textbf{Art. 76} & \textbf{\textquotedbl{}La Secretaría, con el apoyo de la CNE, debe integrar la información sobre los Certificados de Energías Limpias... precios promedio... suspensión y cancelación\textquotedbl{}} \\ \hline
\textbf{FOTEASE (Destino de Recursos)} & \textbf{Reglamento LPTE} & \textbf{Art. 77} & \textbf{\textquotedbl{}El Fondo para la Transición Energética... debe contribuir y ser coherente con las estrategias, líneas de acción y las metas establecidas\textquotedbl{}} \\ \hline
\bottomrule
\end{xltabular}
\end{tabladoradoLargo}

\vspace{0.4cm}
\Needspace{25\baselineskip}
\subsection{B. Matriz de validación jurídica}
\nopagebreak

\vspace{0.3cm}
\Needspace{15\baselineskip}
\begin{tabladoradoLargo}
\renewcommand{\tabularxcolumn}[1]{p{#1}}
\begin{xltabular}{\textwidth}{|X|X|X|X|}
\caption{Matriz de Datos} \\
\toprule
\rowcolor{gobmxDorado} \encabezadotablaDorada{Hallazgo o limitación} & \encabezadotablaDorada{Fundamento jurídico (APA + cita textual)} & \encabezadotablaDorada{Riesgo (jurídico u operativo)} & \encabezadotablaDorada{Ajuste propuesto (alineado al marco vigente)} \\
\midrule
\endhead
Distorsión estructural del mercado por asignación gratuita de CEL & Comisión Reguladora de Energía. (2016). \textit{DACG del Sistema de CEL}. Disposición 34: \textquotedbl{}la Comisión podrá asignarlos de manera no onerosa\textquotedbl{} & Mercado: Subsidio implícito que viola principios de competencia (Art. 6 RLSE) y desincentiva inversión & Eliminar gratuidad e implementar \textbf{Mecanismo Competitivo de Asignación} con precio de reserva \\ \hline
Criterios subjetivos y discrecionales para inclusión de CEL en la bolsa & Comisión Reguladora de Energía. (2016). \textit{DACG del Sistema de CEL}. Disposición 34.4: \textquotedbl{}Aquellos otros que la Comisión considere justificados\textquotedbl{} & Jurídico: Incertidumbre regulatoria incompatible con la LSE 2025 & Establecer supuestos taxativos y eliminar la cláusula discrecional \\ \hline
Falta de transparencia en inventarios y asignaciones & Acuerdo A/012/2020: \textquotedbl{}se realizará durante el proceso de otorgamiento... inmediato posterior\textquotedbl{} & Transparencia: Opacidad contraviene la obligación de información pública del \textbf{Art. 76 RLPTE} & Implementar publicación trimestral de inventarios y resultados en el S-CEL \\ \hline
Inequidad en asignación proporcional al consumo & DACG S-CEL Disp. 34: \textquotedbl{}proporcionalmente al consumo de todos los Participantes Obligados\textquotedbl{} & Equidad: Favorece grandes consumos sin criterios de eficiencia & Sustituir asignación por subasta; ingresos al \textbf{FOTEASE} para beneficio sistémico \\ \hline
Pérdida de recursos públicos por no monetización & DACG S-CEL Disp. 34: No establece contraprestación económica & Fiscal: Desaprovechamiento de activos del Estado que deberían financiar la transición (Art. 77 RLPTE) & Canalizar 100\% de ingresos de subastas al \textbf{FOTEASE} \\ \hline
\bottomrule
\end{xltabular}
\end{tabladoradoLargo}

\subsection{C. Desarrollo analítico}
\nopagebreak

\subsubsection{Diagnóstico de la situación actual}
\nopagebreak

La \textquotedbl{}Bolsa No Onerosa\textquotedbl{} representa una anomalía regulatoria heredada que opera como un subsidio cruzado opaco:


\begin{enumerate}
\item \textbf{Destrucción de Valor:} Al regalar certificados (más de 2 millones en 2018), se deprime el precio de mercado spot y se reduce el incentivo para nuevos proyectos limpios.
\item \textbf{Opacidad:} No existe un inventario público en tiempo real de los CELs que entran a esta bolsa, creando asimetría de información.
\item \textbf{Ineficiencia:} La asignación proporcional al consumo premia el volumen de energía sucia consumida, en lugar de la eficiencia.

\end{enumerate}
\subsubsection{Estado objetivo (LSE 2025)}
\nopagebreak

El nuevo marco exige eliminar la gratuidad y utilizar estos activos pare financiar la transición:


\begin{enumerate}
\item \textbf{Monetización:} Los CELs remanentes son activos del Estado y deben venderse mediante Mecaismos Competitivos.
\item \textbf{Destino Final: \textbf{ Los recursos no van a la tesorería general, sino al }FOTEASE} (Fondo para la Transición Energética y el Aprovechamiento Sustentable de la Energía), cerrando el ciclo virtuoso de financiamiento.
\item \textbf{Transparencia:} Cumplimiento estricto del Art. 76 del Reglamento LPTE sobre publicidad de información de CELs.

\end{enumerate}
\vspace{0.4cm}
\Needspace{25\baselineskip}
\subsubsection{Tabla comparativa: modelo actual vs modelo objetivo}
\nopagebreak

\vspace{0.3cm}
\Needspace{15\baselineskip}
\begin{tablaParadigmaV2}{Análisis Comparativo}{Modelo Actual (CRE)}{Modelo Objetivo (CNE/LSE 2025)}
\tpRow{Mecanismo de Salida}{Asignación Gratuita (\textquotedbl{}No Onerosa\textquotedbl{})}{Mecanismo de Salida}{\textbf{Mecanismo Competitivo (Subasta)}}
\tpRow{Beneficiario}{Participantes Obligados (Proporcional Consumo)}{Beneficiario}{\textbf{FOTEASE} (Financiamiento de Transición)}
\tpRow{Transparencia}{Opaca (Acuerdos retroactivos)}{Transparencia}{\textbf{Pública Trimestral} (Art. 76 RLPTE)}
\tpRow{Impacto Mercado}{Depresión de precios (Dump)}{Impacto Mercado}{Señal de precio piso (Reserva)}
\end{tablaParadigmaV2}

\subsubsection{Arquitectura del sistema}
\nopagebreak

\begin{enumerate}
\item \textbf{Detección:} El sistema identifica automáticamente CELs no reclamados por más de 6 meses (causas taxativas).
\item \textbf{Inventario:} Se transfieren a la \textquotedbl{}Cuenta de Remanentes CNE\textquotedbl{}.
\item \textbf{Subasta:} Trimestralmente se ejecuta el Mecanismo Competitivo con precio de reserva (promedio spot - spread).
\item \textbf{Fondeo:} Los ingresos se depositan en el Fideicomiso del FOTEASE.

\end{enumerate}
\subsection{D. Propuestas de ajuste normativo}
\nopagebreak

\subsubsection{Instrumento A: Emisión de nuevas DACG por la CNE (Sustitución de RES/174/2016)}
\nopagebreak

\textbf{Disposición 34 (Reforma Integral) - Gestión de CEL Remanentes:}

\textquotedbl{}Los CEL que no fueron emitidos ni otorgados por las causas establecidas en la normatividad serán transferidos a la Cuenta de Remanentes administrada por la CNE. Dichos CEL serán ofertados trimestralmente mediante un \textbf{Mecanismo Competitivo de Asignación}, garantizando los principios de transparencia y máxima concurrencia previstos en el artículo 6 del Reglamento de la Ley del Sector Eléctrico.\textquotedbl{}


\textbf{Disposición 34 Bis (Adición) - Destino de Recursos (FOTEASE):}

\textquotedbl{}Los ingresos obtenidos por la asignación de CEL remanentes se destinarán íntegramente al \textbf{Fondo para la Transición Energética y el Aprovechamiento Sustentable de la Energía (FOTEASE) \textbf{, de conformidad con lo establecido en el }artículo 77 del Reglamento de la Ley de Planeación y Transición Energética}, para el financiamiento de proyectos estratégicos de transición.\textquotedbl{}


\textbf{Disposición 34 Ter (Adición) - Prohibición de Gratuidad:}

\textquotedbl{}Queda prohibida cualquier forma de asignación no onerosa o gratuita de Certificados de Energía Limpia, a fin de preservar las señales económicas del mercado y evitar prácticas de subsidio cruzado incompatibles con la competencia del sector.\textquotedbl{}


\textbf{Disposición 34 Quáter (Adición) - Transparencia de Inventarios:}

\textquotedbl{}En cumplimiento del \textbf{artículo 76 del Reglamento de la Ley de Planeación y Transición Energética}, la CNE publicará trimestralmente en el S-CEL: i) el inventario detallado de CEL remanentes, ii) las causas de su ingreso a la cuenta, iii) el precio de reserva y resultado del Mecanismo Competitivo.\textquotedbl{}


\subsection{E. Observaciones técnicas relevantes}
\nopagebreak

\subsubsection{Riesgos identificados}
\nopagebreak

\begin{itemize}
\item \textbf{Resistencia de Participantes:} Los Obligados perderán el beneficio del \textquotedbl{}reparto gratuito\textquotedbl{}. Se debe comunicar como una medida de corrección de mercado necesaria.
\item \textbf{Liquidez de Subastas:} Riesgo de que las subastas queden desiertas si el precio de reserva es muy alto. El mecanismo debe incluir reglas de ajuste de precio.

\end{itemize}
\subsubsection{Dependencias}
\nopagebreak

\begin{itemize}
\item Requiere coordinación operativa con el Comité del FOTEASE para la recepción de recursos.
\item Depende de la implementación del Módulo de Transparencia del B-IV Punto 1.


\end{itemize}
\textbf{Revisión técnica:} Área de Certificados de Energía Limpia

\textbf{Fecha de conclusión:} 12 de enero de 2026


\portadaseccion{BLOQUE V}{Formación del precio del CEL y relación con el factor de emisiones}{}
\clearpage
\phantomsection
\addcontentsline{toc}{section}{Formación del precio del CEL y relación con el factor de emisiones}
\section*{Formación del precio del CEL y relación con el factor de emisiones}
\nopagebreak

\textbf{Alcance del punto:} Análisis de las áreas de oportunidad relacionadas con la formación de precios de los Certificados de Energía Limpia, transparencia, metodologías de descubrimiento de precios y vinculación ambiental.

\textbf{Instrumentos jurídicos revisados:} Reglamento de la Ley de Planeación y Transición Energética (Art. 76, 25, 36), Reglamento de la Ley del Sector Eléctrico 2025 (Art. 6, 4), RES/174/2016, RES/1838/2016.


\vspace{0.4cm}
\Needspace{25\baselineskip}
\subsection{A. Fuentes de información del S-CEL para el Punto 1}
\nopagebreak

\vspace{0.3cm}
\Needspace{15\baselineskip}
\begin{tabladoradoLargo}
\renewcommand{\tabularxcolumn}[1]{p{#1}}
\begin{xltabular}{\textwidth}{|X|X|X|X|}
\caption{Matriz de Datos} \\
\toprule
\rowcolor{gobmxDorado} \encabezadotablaDorada{Actor / Fuente} & \encabezadotablaDorada{Instrumento legal} & \encabezadotablaDorada{Artículo / Numeral} & \encabezadotablaDorada{Cita textual} \\
\midrule
\endhead
\textbf{CNE (Obligación de Información)} & \textbf{Reglamento LPTE} & \textbf{Art. 76} & \textbf{\textquotedbl{}La Secretaría, con el apoyo de la CNE, debe integrar la información sobre los Certificados de Energías Limpias... precios promedio de los certificados reportados\textquotedbl{}} \\ \hline
CENACE (Mercado Centralizado) & RES/174/2016 & Disp. 42 & \textquotedbl{}El día hábil siguiente... el Cenace entregará a la Comisión... un reporte electrónico de las operaciones\textquotedbl{} \\ \hline
Participantes (Transacciones Bilaterales) & RES/174/2016 & Disp. 60 & \textquotedbl{}los precios de los CEL en el mercado de CEL, transacciones bilaterales\textquotedbl{} \\ \hline
Boletín Electrónico & RES/174/2016 & Disp. 46 & \textquotedbl{}Los Participantes... podrán ofertarlos en el Boletín Electrónico indicando el precio\textquotedbl{} \\ \hline
\textbf{Metas de Emisiones} & \textbf{Reglamento LPTE} & \textbf{Art. 25} & \textbf{\textquotedbl{}Metas... referidas a lo siguiente: II. Reducción de... c) Emisiones de gases y compuestos de efecto invernadero\textquotedbl{}} \\ \hline
Criterios de Eficiencia & RES/1838/2016 & Tabla 1 & \textquotedbl{}Metodología aplicable para determinar el porcentaje de energía libre de combustible\textquotedbl{} \\ \hline
\bottomrule
\end{xltabular}
\end{tabladoradoLargo}

\vspace{0.4cm}
\Needspace{25\baselineskip}
\subsection{B. Matriz de validación jurídica}
\nopagebreak

\vspace{0.3cm}
\Needspace{15\baselineskip}
\begin{tabladoradoLargo}
\renewcommand{\tabularxcolumn}[1]{p{#1}}
\begin{xltabular}{\textwidth}{|X|X|X|X|}
\caption{Matriz de Datos} \\
\toprule
\rowcolor{gobmxDorado} \encabezadotablaDorada{Hallazgo o limitación} & \encabezadotablaDorada{Fundamento jurídico (APA + cita textual)} & \encabezadotablaDorada{Riesgo (jurídico u operativo)} & \encabezadotablaDorada{Ajuste propuesto (alineado al marco vigente)} \\
\midrule
\endhead
Ausencia de metodología de reporte para precios bilaterales (Opacidad) & RES/174/2016, Disp. 60: Menciona \textquotedbl{}precios\textquotedbl{} sin obligar metodología. & Transparencia: Incumplimiento del \textbf{Art. 76 R-LPTE} que exige integrar \textquotedbl{}precios promedio reportados\textquotedbl{}. & Establecer obligación de reporte de \textbf{Precio Real de Transacción} para cumplimiento del Art. 76 R-LPTE. \\ \hline
Desconexión valor CEL vs Valor Ambiental & Aviso FE 2025 (0.395 tCO2e/MWh) sin vínculo legal al precio. & Ambiental: El precio no refleja la contribución a las metas del \textbf{Art. 25 R-LPTE}. & Crear \textquotedbl{}Índice de Referencia Ambiental\textquotedbl{} derivado de las Metas de Transición (Art. 25 y 36 R-LPTE). \\ \hline
Ineficiencia del Boletín Electrónico & RES/174/2016, Disp. 46: Simple tablero de anuncios. & Eficiencia: Subutilización de tecnología permitida por el \textbf{Art. 4 RLSE}. & Implementar reglas de \textbf{Matching Automático} en el Boletín (Art. 4 RLSE - Herramientas Tecnológicas). \\ \hline
Falta de granularidad en la información pública & RES/174/2016, Disp. 60: Reportes agregados. & Mercado: Asimetría de información que favorece a grandes jugadores. & Publicación desglosada (anonimizada) para verdadera señal de precios (Art. 6 RLSE - Transparencia). \\ \hline
\bottomrule
\end{xltabular}
\end{tabladoradoLargo}

\subsection{C. Desarrollo analítico}
\nopagebreak

\subsubsection{Diagnóstico de la situación actual}
\nopagebreak

El mercado de CELs opera con una \textquotedbl{}ceguera parcial\textquotedbl{} en la formación de precios. Si bien el \textbf{70\% del volumen} se transa bilateralmente, no existe un mecanismo obligatorio y estandarizado para capturar el precio real de cierre de estos contratos. La autoridad solo ve \textquotedbl{}promedios\textquotedbl{} o estimaciones, lo que impide detectar abuso de mercado o volatilidad real.


\vspace{0.2cm}
\Needspace{8\baselineskip}
\begin{glowBox}[gobmxDorado]{Nota Informativa}

\end{glowBox}
\vspace{0.2cm}
\Needspace{8\baselineskip}
\begin{glowBox}[gobmxDorado]{Nota Informativa}
\textbf{Nota Técnica: Necesidad de Precios Reales}
\end{glowBox}
\vspace{0.2cm}
\Needspace{8\baselineskip}
\begin{glowBox}[gobmxDorado]{Nota Informativa}
Para un entendimiento preciso del comportamiento del mercado, \textbf{es indispensable contar con precios reales de transacción unitarios}, no solo índices agregados o estimaciones. Los promedios pueden enmascarar picos de volatilidad, tendencias de concentración y comportamientos anticompetitivos que solo son visibles analizando la dispersión real de los precios de cierre.
\end{glowBox}

\subsubsection{Estado objetivo}
\nopagebreak

El nuevo marco del Reglamento de la Ley de Planeación y Transición Energética (Art. 76) otorga a la CNE la facultad explícita de \textquotedbl{}integrar información... sobre precios promedio\textquotedbl{}. Para cumplirlo, el sistema debe evolucionar de un \textquotedbl{}buzón de reportes\textquotedbl{} a un \textbf{Monitor de Precios en Tiempo Real}.


La vinculación ambiental también cambia: el CEL deja de ser solo un requisito administrativo para convertirse en el instrumento financiero de las \textbf{Metas de Descarbonización (Art. 25 R-LPTE)}. El precio del CEL debería, teóricamente, converger hacia el Costo Marginal de Abatimiento de CO2.


\vspace{0.4cm}
\Needspace{25\baselineskip}
\subsubsection{Tabla comparativa: modelo actual vs modelo objetivo}
\nopagebreak

\vspace{0.3cm}
\Needspace{15\baselineskip}
\begin{tablaParadigmaV2}{Análisis Comparativo}{Modelo Actual (CRE)}{Modelo Objetivo (CNE/LSE 2025)}
\tpRow{Visibilidad Precios Bilaterales}{Nula/Voluntaria (Opacidad)}{Visibilidad Precios Bilaterales}{\textbf{Obligatoria y Real} (Art. 76 R-LPTE)}
\tpRow{Tecnología de Mercado}{Boletín pasivo (Clasificados)}{Tecnología de Mercado}{\textbf{Matching Activo} (Art. 4 RLSE)}
\tpRow{Señal de Valor}{Oferta/Demanda administrativa}{Señal de Valor}{\textbf{Referencia Ambiental} (Costo Social Carbono)}
\tpRow{Granularidad de Datos}{Promedios mensuales}{Granularidad de Datos}{Reporte transacción por transacción (Base de Datos)}
\end{tablaParadigmaV2}

\subsubsection{Arquitectura del sistema}
\nopagebreak

\begin{enumerate}
\item \textbf{Captura:} Cada contrato bilateral registrado en el S-CEL deberá incluir el campo obligatorio \textquotedbl{}Precio Unitario de Liquidación\textquotedbl{} (auditable).
\item \textbf{Validación:} Algoritmos automáticos detectan outliers (precios fuera de rango lógico) para revisión.
\item \textbf{Cálculo:} Se generan índices diarios: Volumen Ponderado (VWAP), Precio Cierre, Spread Bid-Ask.
\item \textbf{Referencia:} Se publica el \textquotedbl{}Precio Ambiental Teórico\textquotedbl{} basado en el Factor de Emisión SEN y el Precio Social del Carbono de la Estrategia de Transición.

\end{enumerate}
\subsection{D. Propuestas de ajuste normativo}
\nopagebreak

\subsubsection{Instrumento A: Emisión de nuevas DACG por la CNE (Sustitución RES/174/2016)}
\nopagebreak

\textbf{Disposición 60.A (Nueva) - Obligación de Reporte de Precios Reales:}

\textquotedbl{}En cumplimiento del \textbf{artículo 76 del Reglamento de la Ley de Planeación y Transición Energética}, los Participantes del Mercado deberán reportar al S-CEL el precio real unitario pactado en cada transacción bilateral. La CNE garantizará la confidencialidad de los datos individuales, publicando únicamente información agregada y anonimizada necesaria para las señales de mercado.\textquotedbl{}


\textbf{Disposición 46.A (Nueva) - Plataforma de Negociación (Smart Matching):}

\textquotedbl{}Con fundamento en el \textbf{artículo 4 del Reglamento de la LSE}, el Boletín Electrónico incorporará herramientas tecnológicas de emparejamiento automático (matching) de ofertas y demandas compatibles, facilitando la liquidez y el descubrimiento de precios eficiente.\textquotedbl{}


\textbf{Disposición 60.B (Nueva) - Índice de Referencia Ambiental:}

\textquotedbl{}Para coadyuvar al cumplimiento de las metas establecidas en el \textbf{artículo 25 del Reglamento de la Ley de Planeación y Transición Energética}, la CNE publicará anualmente un 'Precio de Referencia Ambiental del CEL', calculado metodológicamente a partir del Factor de Emisión del SEN y el Costo Social del Carbono definido en la Estrategia de Transición.\textquotedbl{}


\subsubsection{Instrumento B: Manual de Transparencia del Mercado}
\nopagebreak

\begin{itemize}
\item \textbf{Protocolo de Datos:} Definición de los campos \textquotedbl{}Precio\textquotedbl{}, \textquotedbl{}Divisa\textquotedbl{} y \textquotedbl{}Fecha de Pacto\textquotedbl{} como obligatorios e inalterables en el registro de bilaterales.
\item \textbf{Índices Públicos:} Metodología para el cálculo del VWAP (Volume Weighted Average Price) diario.

\end{itemize}
\subsection{E. Observaciones técnicas relevantes}
\nopagebreak

\subsubsection{Riesgos identificados}
\nopagebreak

\begin{itemize}
\item \textbf{Evasión de Reporte:} Participantes podrían reportar precios simbólicos (\$0 o \$0.01) para ocultar valor real. Se requieren reglas de validación (precios fuera de mercado activan auditoría).
\item \textbf{Complejidad de Datos:} El manejo de precios reales requiere protocolos de ciberseguridad robustos para proteger secretos comerciales (Art. 68 R-LPTE).


\end{itemize}
\textbf{Revisión técnica:} Área de Certificados de Energía Limpia

\textbf{Fecha de conclusión:} 12 de enero de 2026


\portadaseccion{BLOQUE VI}{Certificados de Energía Limpia en contratos de cobertura}{}
\clearpage
\phantomsection
\addcontentsline{toc}{section}{Certificados de Energía Limpia en contratos de cobertura}
\section*{Certificados de Energía Limpia en contratos de cobertura}
\nopagebreak

\textbf{Alcance del punto:} Análisis de la integración de CEL en Contratos de Cobertura bajo el esquema de Planeación Centralizada, asignación prioritaria a Centrales del Estado y cumplimiento de metas del PROSENER.

\textbf{Instrumentos jurídicos revisados:} Ley del Sector Eléctrico 2025 (Arts. 53, 54), Reglamento LSE 2025 (Art. 20), Programa Sectorial de Energía (PROSENER), Reglas del Mercado Eléctrico.


\vspace{0.4cm}
\Needspace{25\baselineskip}
\subsection{A. Fuentes de información del S-CEL para el Punto 1}
\nopagebreak

\vspace{0.3cm}
\Needspace{15\baselineskip}
\begin{tabladoradoLargo}
\renewcommand{\tabularxcolumn}[1]{p{#1}}
\begin{xltabular}{\textwidth}{|X|X|X|X|}
\caption{Matriz de Datos} \\
\toprule
\rowcolor{gobmxDorado} \encabezadotablaDorada{Actor / Fuente} & \encabezadotablaDorada{Instrumento legal} & \encabezadotablaDorada{Artículo / Numeral} & \encabezadotablaDorada{Cita textual} \\
\midrule
\endhead
\textbf{Suministro Básico (Asignación)} & \textbf{Ley del Sector Eléctrico 2025} & \textbf{Art. 53} & \textbf{\textquotedbl{}La CFE Suministrador de Servicios Básicos celebrará Contratos de Cobertura Eléctrica prioritariamente con las Centrales Eléctricas del Estado para garantizar la confiabilidad y soberanía...\textquotedbl{}} \\ \hline
\textbf{Planeación de Cobertura} & \textbf{Ley del Sector Eléctrico 2025} & \textbf{Art. 54} & \textbf{\textquotedbl{}La Secretaría establecerá los términos, plazos y porcentajes de cobertura que deberán observar los Suministradores... conforme a las metas del PROSENER\textquotedbl{}} \\ \hline
\textbf{Integración de CELs} & \textbf{Reglamento LSE 2025} & \textbf{Art. 20} & \textbf{\textquotedbl{}Los Contratos de Cobertura Eléctrica deberán especificar la cantidad de Productos Asociados (CEL) incluidos, garantizando el cumplimiento de las obligaciones de energías limpias\textquotedbl{}} \\ \hline
CENACE (Registro) & Reglas del Mercado & Regla 12 & \textquotedbl{}El CENACE llevará el registro de los Contratos de Cobertura y verificará su alineación con la Planeación Nacional\textquotedbl{} \\ \hline
Metas Nacionales & PROSENER 2025-2030 & Estrategia 2.1 & \textquotedbl{}Alcanzar el 54\% de generación limpia mediante la asignación eficiente de recursos en las Centrales del Estado\textquotedbl{} \\ \hline
\bottomrule
\end{xltabular}
\end{tabladoradoLargo}

\vspace{0.4cm}
\Needspace{25\baselineskip}
\subsection{B. Matriz de validación jurídica}
\nopagebreak

\vspace{0.3cm}
\Needspace{15\baselineskip}
\begin{tabladoradoLargo}
\renewcommand{\tabularxcolumn}[1]{p{#1}}
\begin{xltabular}{\textwidth}{|X|X|X|X|}
\caption{Matriz de Datos} \\
\toprule
\rowcolor{gobmxDorado} \encabezadotablaDorada{Hallazgo o limitación} & \encabezadotablaDorada{Fundamento jurídico (APA + cita textual)} & \encabezadotablaDorada{Riesgo (jurídico u operativo)} & \encabezadotablaDorada{Ajuste propuesto (alineado al marco vigente)} \\
\midrule
\endhead
Obsolescencia de \textquotedbl{}Contratos Legados\textquotedbl{} y Subastas & RES/008/2016 (Marco derogado). El concepto de \textquotedbl{}Legado\textquotedbl{} desaparece ante la reintegración de CFE. & Jurídico: Incompatibilidad con el \textbf{Art. 53 LSE 2025} que unifica el parque de generación estatal. & Sustituir por \textbf{\textquotedbl{}Contratos de Cobertura de Centrales del Estado\textquotedbl{}} con asignación directa basada en Planeación. \\ \hline
Porcentajes de cobertura desactualizados & Referencias a metas 2018 (RES/584). & Regulatorio: Vacío normativo sobre obligaciones actuales de contratación. & Vincular porcentajes dinámicamente a las \textbf{Metas Anuales del PROSENER} y notificaciones de SENER. \\ \hline
Falta de garantías en transferencias internas & Esquema anterior de mercado vs. transferencias entre empresas del Estado. & Operativo: Riesgo de \textquotedbl{}incumplimiento administrativo\textquotedbl{} entre filiales de CFE. & Implementar \textbf{Mecanismos de Cumplimiento Operativo} supervisados por CNE, en lugar de garantías financieras bancarias. \\ \hline
Descoordinación Planeación-Contratación & Contratos basados en pronósticos de mercado vs. Planeación Central. & Soberanía: Desviación de las metas de Seguridad Energética. & Alinear plazos y volúmenes estrictamente al \textbf{Programa de Desarrollo del Sector Eléctrico (PRODESEN)}. \\ \hline
\bottomrule
\end{xltabular}
\end{tabladoradoLargo}

\subsection{C. Desarrollo analítico}
\nopagebreak

\subsubsection{Diagnóstico de la situación actual}
\nopagebreak

El esquema anterior de \textquotedbl{}Contratos Legados\textquotedbl{} y \textquotedbl{}Subastas de Largo Plazo\textquotedbl{} fue diseñado para un mercado de múltiples compradores y vendedores privados. Bajo la \textbf{LSE 2025}, este modelo es incompatible con la re-centralización de la planeación:


\begin{enumerate}
\item \textbf{Fragmentación:} La distinción \textquotedbl{}Legado\textquotedbl{} vs \textquotedbl{}Nuevo\textquotedbl{} carece de sentido cuando el objetivo es optimizar todo el parque estatal.
\item \textbf{Incertidumbre:} La falta de subastas desde 2019 dejó a los suministradores sin mecanismo claro de cobertura a largo plazo.
\item \textbf{Soberanía:} El suministro básico depende de contratos que deben priorizar la seguridad nacional, no solo el precio marginal.

\end{enumerate}
\subsubsection{Estado objetivo (LSE 2025)}
\nopagebreak

El nuevo modelo se basa en la \textbf{Asignación Estratégica}:


\begin{enumerate}
\item \textbf{Prioridad Estatal:} El Suministro Básico se cubre prioritariamente con Centrales del Estado (Art. 53 LSE).
\item \textbf{Planeación Vinculante:} Los volúmenes de energía y CELs se definen en el PROSENER/PRODESEN, no por \textquotedbl{}oferta y demanda\textquotedbl{} espontánea.
\item \textbf{Integralidad:} El contrato de cobertura es un instrumento de política pública que asegura tarifas justas y cumplimiento de metas limpias.

\end{enumerate}
\vspace{0.4cm}
\Needspace{25\baselineskip}
\subsubsection{Tabla comparativa: modelo actual vs modelo objetivo}
\nopagebreak

\vspace{0.3cm}
\Needspace{15\baselineskip}
\begin{tablaParadigmaV2}{Análisis Comparativo}{Modelo Anterior (Mercado)}{Modelo Objetivo (Soberanía/Planeación)}
\tpRow{Mecanismo de Asignación}{Subastas de Largo Plazo (Precio)}{Mecanismo de Asignación}{\textbf{Asignación Directa / Procesos Competitivos Estratégicos}}
\tpRow{Contraparte Principal}{Generadores Privados / Legados}{Contraparte Principal}{\textbf{Centrales Eléctricas del Estado}}
\tpRow{Definición de Volúmenes}{Estimaciones de Mercado}{Definición de Volúmenes}{\textbf{Metas del PROSENER}}
\tpRow{Garantías}{Cartas de Crédito, Garantías Financieras}{Garantías}{\textbf{Acuerdos de Nivel de Servicio (SLA) Operativos}}
\tpRow{Objetivo}{Minimización de Costos Financieros}{Objetivo}{\textbf{Seguridad Energética y Tarifas Justas}}
\end{tablaParadigmaV2}

\subsubsection{Arquitectura del sistema}
\nopagebreak

\begin{enumerate}
\item \textbf{Planeación:} SENER define en PROSENER la demanda esperada y la cuota de CELs requerida.
\item \textbf{Asignación:} CNE asigna a las Centrales del Estado los volúmenes de cobertura para el Suministro Básico.
\item \textbf{Formalización:} Se firman \textquotedbl{}Contratos de Cobertura Estratégica\textquotedbl{} registrados en CENACE.
\item \textbf{Ejecución:} Transferencia automática de CELs mensual basada en generación real verificada.

\end{enumerate}
\subsection{D. Propuestas de ajuste normativo}
\nopagebreak

\subsubsection{Instrumento A: Disposiciones para la Asignación de Cobertura Eléctrica (Nueva DACG CNE)}
\nopagebreak

\textbf{Disposición 1 (Objeto):}

\textquotedbl{}Establecer los mecanismos para la asignación de Contratos de Cobertura Eléctrica entre el Suministrador de Servicios Básicos y las Centrales Eléctricas del Estado, en cumplimiento del \textbf{artículo 53 de la Ley del Sector Eléctrico}.\textquotedbl{}


\textbf{Disposición 5 (Determinación de Volúmenes):}

\textquotedbl{}Los volúmenes de energía y Certificados de Energía Limpia objeto de cobertura serán determinados anualmente por la Secretaría, en congruencia con las metas de generación limpia establecidas en el \textbf{PROSENER} y los estudios de confiabilidad del CENACE.\textquotedbl{}


\textbf{Disposición 8 (Mecanismo de Cumplimiento):}

\textquotedbl{}Las Centrales del Estado garantizarán el cumplimiento de entrega de CELs mediante Acuerdos de Nivel de Servicio Operativo. En caso de indisponibilidad, el CENACE gestionará la compensación mediante el portafolio consolidado de generación estatal.\textquotedbl{}


\subsubsection{Instrumento B: Actualización de Manuales del Mercado}
\nopagebreak

\textbf{Manual de Registro de Contratos:}

\textquotedbl{}Se crea la categoría 'Contrato de Cobertura Estratégica' para instrumentar las asignaciones del Art. 53 LSE, exentos de los requisitos de garantías financieras de mercado, sujetos a supervisión operativa de CNE.\textquotedbl{}


\subsection{E. Observaciones técnicas relevantes}
\nopagebreak

\subsubsection{Riesgos identificados}
\nopagebreak

\begin{itemize}
\item \textbf{Rigidez Operativa:} La asignación centralizada puede ser lenta para reaccionar a cambios bruscos en la demanda. Se requieren cláusulas de ajuste dinámico.
\item \textbf{Suficiencia Financiera:} Es crucial que las tarifas de estos contratos reconozcan los costos reales de generación limpia para garantizar el mantenimiento de las centrales estatales.


\end{itemize}
\textbf{Revisión técnica:} Área de Certificados de Energía Limpia

\textbf{Fecha de conclusión:} 12 de enero de 2026


\phantomsection
\addcontentsline{toc}{section}{Incorporación de Metas de CEL en la Planeación Centralizada (PROSENER)}
\section*{Incorporación de Metas de CEL en la Planeación Centralizada (PROSENER)}
\nopagebreak

\textbf{Alcance del punto:} Análisis de la integración de las metas de Certificados de Energía Limpia dentro del Sistema de Planeación Democrática y Soberana, específicamente en el PROSENER, sustituyendo los mecanismos de subasta por asignaciones estratégicas vinculantes.

\textbf{Instrumentos jurídicos revisados:} Reglamento de la Ley del Sector Eléctrico 2025 (Arts. 28, 30, 31), Reglamento de la Ley de Planeación y Transición Energética (Arts. 33, 35, 77), PROSENER.


\vspace{0.4cm}
\Needspace{25\baselineskip}
\subsection{A. Fuentes de información del S-CEL para el Punto 2}
\nopagebreak

\vspace{0.3cm}
\Needspace{15\baselineskip}
\begin{tabladoradoLargo}
\renewcommand{\tabularxcolumn}[1]{p{#1}}
\begin{xltabular}{\textwidth}{|X|X|X|X|}
\caption{Matriz de Datos} \\
\toprule
\rowcolor{gobmxDorado} \encabezadotablaDorada{Actor / Fuente} & \encabezadotablaDorada{Instrumento legal} & \encabezadotablaDorada{Artículo / Numeral} & \encabezadotablaDorada{Cita textual} \\
\midrule
\endhead
\textbf{Instrumento Rector} & \textbf{Reglamento LPTE} & \textbf{Art. 33} & \textbf{\textquotedbl{}El PROSENER es el instrumento rector de la planeación de corto plazo del Sector Energético...\textquotedbl{}} \\ \hline
\textbf{Contenido del PROSENER} & \textbf{Reglamento LPTE} & \textbf{Art. 35} & \textbf{\textquotedbl{}El PROSENER debe contener... metas de participación de energías limpias desglosadas por tecnología y región.\textquotedbl{}} \\ \hline
\textbf{Vinculación de Permisos} & \textbf{Reglamento LSE 2025} & \textbf{Art. 30} & \textbf{\textquotedbl{}La Secretaría y la CNE... deben considerar lo establecido en los instrumentos de planeación... y el cumplimiento de la planeación vinculante.\textquotedbl{}} \\ \hline
\textbf{Criterios de Evaluación} & \textbf{Reglamento LSE 2025} & \textbf{Art. 31} & \textbf{\textquotedbl{}Las disposiciones... deben establecer los criterios que la Secretaría y la CNE deben considerar al evaluar el cumplimiento de la planeación vinculante.\textquotedbl{}} \\ \hline
\textbf{Proyectos Estratégicos} & \textbf{Reglamento LSE 2025} & \textbf{Art. 28} & \textbf{\textquotedbl{}Permisos de generación... para Proyectos Estratégicos y aquellos necesarios para salvaguardar la integridad del SEN.\textquotedbl{}} \\ \hline
\textbf{Financiamiento} & \textbf{Reglamento LPTE} & \textbf{Art. 77} & \textbf{\textquotedbl{}El FOTEASE... debe contribuir y ser coherente con las estrategias... y las metas establecidas en los Instrumentos de Planeación.\textquotedbl{}} \\ \hline
\bottomrule
\end{xltabular}
\end{tabladoradoLargo}

\vspace{0.4cm}
\Needspace{25\baselineskip}
\subsection{B. Matriz de validación jurídica}
\nopagebreak

\vspace{0.3cm}
\Needspace{15\baselineskip}
\begin{tabladoradoLargo}
\renewcommand{\tabularxcolumn}[1]{p{#1}}
\begin{xltabular}{\textwidth}{|X|X|X|X|}
\caption{Matriz de Datos} \\
\toprule
\rowcolor{gobmxDorado} \encabezadotablaDorada{Hallazgo o limitación} & \encabezadotablaDorada{Fundamento jurídico (APA + cita textual)} & \encabezadotablaDorada{Riesgo (jurídico u operativo)} & \encabezadotablaDorada{Ajuste propuesto (alineado al marco vigente)} \\
\midrule
\endhead
Propuesta obsoleta de \textquotedbl{}Reactivar Subastas\textquotedbl{} & El término \textquotedbl{}Subasta\textquotedbl{} ha sido suprimido de la LSE para contratación básica. El paradigma actual es la \textbf{Planeación Vinculante} (Art. 30 R-LSE). & Jurídico: Inviabilidad legal de convocar subastas de mercado bajo un marco de asignación centralizada. & Sustituir \textquotedbl{}Subastas\textquotedbl{} por \textbf{\textquotedbl{}Mecanismos de Asignación Estratégica alineados al PROSENER\textquotedbl{}}. \\ \hline
Desconexión entre Metas CEL y Permisos & Art. 30 R-LSE obliga a alinear permisos con planeación, pero falta el mecanismo operativo. & Operativo: Otorgamiento de permisos desordenados que no abonan a las metas regionales de CEL. & Establecer que la CNE solo otorgará permisos que cubran las \textbf{Brechas de CEL identificadas en el PROSENER}. \\ \hline
Falta de señal de inversión a largo plazo & El mercado espera \textquotedbl{}Precios de Subasta\textquotedbl{}. El nuevo modelo ofrece \textquotedbl{}Certeza de Planeación\textquotedbl{}. & Económico: Retraimiento de inversión privada por falta de señales claras. & Publicar los \textbf{Mapas de Ruta Tecnológicos} (Art. 60 R-LPTE) como la nueva señal de demanda firme del Estado. \\ \hline
Financiamiento disperso & Propuesta anterior de \textquotedbl{}Fondos de Subasta\textquotedbl{} sin base legal. & Fiscal: Indeterminación de recursos. & Vincular el desarrollo de proyectos limpios al \textbf{FOTEASE} (Art. 77 R-LPTE) para apalancamiento financiero. \\ \hline
\bottomrule
\end{xltabular}
\end{tabladoradoLargo}

\subsection{C. Desarrollo analítico}
\nopagebreak

\subsubsection{Diagnóstico de la situación actual}
\nopagebreak

El análisis basado en la \textquotedbl{}ausencia de subastas\textquotedbl{} es un diagnóstico del pasado. El verdadero reto actual no es la falta de un mecanismo de compra (subasta), sino la falta de \textbf{articulación operativa} entre las Metas de Transición plasmadas en el PROSENER y la ejecución diaria de permisos y contratos por la CNE y CFE.


Bajo la LSE 2025, el Estado recupera la rectoría. El problema no es \textquotedbl{}¿quién vende más barato?\textquotedbl{}, sino \textquotedbl{}¿cómo aseguramos que se construya lo que el sistema necesita donde lo necesita?\textquotedbl{}.


\vspace{0.2cm}
\Needspace{8\baselineskip}
\begin{glowBox}[gobmxDorado]{Nota Informativa}

\end{glowBox}
\vspace{0.2cm}
\Needspace{8\baselineskip}
\begin{glowBox}[gobmxDorado]{Nota Informativa}
\textbf{Cambio de Paradigma}
\end{glowBox}
\vspace{0.2cm}
\Needspace{8\baselineskip}
\begin{glowBox}[gobmxDorado]{Nota Informativa}
De: \textquotedbl{}Competencia por el Mercado\textquotedbl{} (Subastas de Precio).
\end{glowBox}
\vspace{0.2cm}
\Needspace{8\baselineskip}
\begin{glowBox}[gobmxDorado]{Nota Informativa}
A: \textquotedbl{}Planeación Vinculante\textquotedbl{} (Asignación por Seguridad y Soberanía).
\end{glowBox}

\subsubsection{Estado objetivo}
\nopagebreak

El modelo objetivo es un ecosistema de \textbf{Ejecución de la Planeación}:


\begin{enumerate}
\item \textbf{PROSENER} define: \textquotedbl{}Se necesitan 2 GW eólicos en Tamaulipas para 2027 para cumplir metas de CEL\textquotedbl{}.
\item \textbf{CNE} ejecuta: \textquotedbl{}Solo otorgo permisos en Tamaulipas a proyectos eólicos alineados\textquotedbl{}.
\item \textbf{CFE} contrata: \textquotedbl{}Firma contratos de cobertura (Art. 53 LSE) con esos proyectos estratégicos\textquotedbl{}.
\item \textbf{FOTEASE} apoya: \textquotedbl{}Financia infraestructura de transmisión necesaria\textquotedbl{}.

\end{enumerate}
\vspace{0.4cm}
\Needspace{25\baselineskip}
\subsubsection{Tabla comparativa: modelo anterior vs modelo objetivo}
\nopagebreak

\vspace{0.3cm}
\Needspace{15\baselineskip}
\begin{tablaParadigmaV2}{Análisis Comparativo}{Modelo Anterior (Mercado)}{Modelo Objetivo (Planeación LSE 2025)}
\tpRow{Señal de Inversión}{Precios de Subasta (Volátiles)}{Señal de Inversión}{\textbf{Metas del PROSENER (Vinculantes)}}
\tpRow{Mecanismo de Entrada}{Libre Mercado (cualquiera oferta)}{Mecanismo de Entrada}{\textbf{Permisos Alineados} (Art. 30 R-LSE)}
\tpRow{Criterio de Selección}{Menor Costo Marginal}{Criterio de Selección}{\textbf{Seguridad Energética y Balance Regional}}
\tpRow{Rol del Privado}{Oferente especulativo}{Rol del Privado}{\textbf{Aliado Estratégico en Proyectos Definidos}}
\tpRow{Financiamiento}{Privado (Project Finance)}{Financiamiento}{\textbf{Mixto / FOTEASE}}
\end{tablaParadigmaV2}

\subsubsection{Arquitectura del sistema}
\nopagebreak

\begin{enumerate}
\item \textbf{Input:} SENER publica la actualización anual del PROSENER con el \textquotedbl{}Balance de Metas CEL\textquotedbl{} por región.
\item \textbf{Filtrado:} CNE actualiza los Criterios de Planeación Vinculante (Art. 31 R-LSE).
\item \textbf{Asignación:}
\end{enumerate}
\begin{itemize}
\item \textbf{Proyectos Públicos:} Asignación presupuestal y FOTEASE directa a CFE.
\item \textbf{Proyectos Privados:} \textquotedbl{}Convocatorias de Alineación\textquotedbl{} donde privados proponen cubrir brechas específicas.
\end{itemize}
\begin{enumerate}
\item \textbf{Cierre:} Firma de Contratos de Cobertura de Centrales (estatales o estratégicas).

\end{enumerate}
\subsection{D. Propuestas de ajuste normativo}
\nopagebreak

\subsubsection{Instrumento A: Disposiciones Administrativas de Planeación Vinculante (Art. 31 R-LSE)}
\nopagebreak

\textbf{Disposición 1 (Objeto):}

\textquotedbl{}Establecer los criterios técnicos que la CNE observará para otorgar permisos de generación, garantizando su estricta alineación con las metas de energías limpias del \textbf{PROSENER}.\textquotedbl{}


\textbf{Disposición 4 (Brechas de CEL):}

\textquotedbl{}La CNE publicará anualmente el 'Reporte de Brechas de CEL Regionales'. Solo se admitirán a trámite solicitudes de permiso que demuestren contribuir al cierre de dichas brechas.\textquotedbl{}


\subsubsection{Instrumento B: Manual de Uso del FOTEASE para Proyectos Estratégicos (Art. 77 R-LPTE)}
\nopagebreak

\textbf{Regla 3:}

\textquotedbl{}Prioridad de Financiamiento. El FOTEASE destinará recursos prioritarios a proyectos de generación limpia que la Secretaría califique como 'Estratégicos' conforme al Art. 28 del Reglamento de la LSE, asegurando el cumplimiento de las metas nacionales.\textquotedbl{}


\subsection{E. Observaciones técnicas relevantes}
\nopagebreak

\subsubsection{Riesgos identificados}
\nopagebreak

\begin{itemize}
\item \textbf{Capacidad de Ejecución Estatal:} Riesgo de que la CFE y FOTEASE no tengan la velocidad de despliegue que tenía el mercado privado. Se mitiga con \textquotedbl{}Convocatorias de Alianza\textquotedbl{} claras.
\item \textbf{Señal de Precios:} Al eliminar la subasta, se pierde la referencia de \textquotedbl{}precio eficiente\textquotedbl{}. Se requiere un \textit{benchmark} internacional de costos (LCOE) para evitar sobrecostos en los contratos asignados.


\end{itemize}
\textbf{Revisión técnica:} Área de Certificados de Energía Limpia

\textbf{Fecha de conclusión:} 12 de enero de 2026


\portadaseccion{BLOQUE VII}{Liquidación Anual de Obligaciones (Sustitución de DECLARACEL)}{}
\clearpage
\phantomsection
\addcontentsline{toc}{section}{Liquidación Anual de Obligaciones (Sustitución de DECLARACEL)}
\section*{Liquidación Anual de Obligaciones (Sustitución de DECLARACEL)}
\nopagebreak

\textbf{Alcance del punto:} Modernización del cumplimiento de obligaciones de CEL, transitando del modelo declarativo manual (DECLARACEL) a un sistema de liquidación administrativa automática basado en la interoperabilidad de datos CNE-CENACE.

\textbf{Instrumentos jurídicos revisados:} Reglamento de la Ley del Sector Eléctrico 2025 (Arts. 4, 182), Reglamento LPTE (Art. 76), RES/174/2016 (Marco Derogado).


\vspace{0.4cm}
\Needspace{25\baselineskip}
\subsection{A. Fuentes de información del S-CEL para el Punto 1}
\nopagebreak

\vspace{0.3cm}
\Needspace{15\baselineskip}
\begin{tabladoradoLargo}
\renewcommand{\tabularxcolumn}[1]{p{#1}}
\begin{xltabular}{\textwidth}{|X|X|X|X|}
\caption{Matriz de Datos} \\
\toprule
\rowcolor{gobmxDorado} \encabezadotablaDorada{Actor / Fuente} & \encabezadotablaDorada{Instrumento legal} & \encabezadotablaDorada{Artículo / Numeral} & \encabezadotablaDorada{Cita textual} \\
\midrule
\endhead
\textbf{CNE (Facultad de Liquidación)} & \textbf{Reglamento LSE 2025} & \textbf{Art. 182} & \textbf{\textquotedbl{}La CNE... debe emitir las disposiciones... que regulen el otorgamiento, liquidación, cancelación voluntaria y transacciones...\textquotedbl{}} \\ \hline
\textbf{Uso de Tecnología} & \textbf{Reglamento LSE 2025} & \textbf{Art. 4} & \textbf{\textquotedbl{}El SEN y el MEM deben incorporar las mejores prácticas... y hacer uso de herramientas tecnológicas avanzadas para la operación...\textquotedbl{}} \\ \hline
\textbf{Transparencia de Resultados} & \textbf{Reglamento LPTE} & \textbf{Art. 76} & \textbf{\textquotedbl{}La Secretaría... debe integrar información... sobre el otorgamiento, liquidación y cancelación [de CELs]...\textquotedbl{}} \\ \hline
\textit{Modulo Histórico (DECLARACEL)} & \textit{RES/174/2016} & \textit{Disp. 31/48} & \textit{\textquotedbl{}Plataforma informática... DECLARACEL que permitirá a los Participantes Obligados realizar las liquidaciones provisionales\textquotedbl{}} (Referencia de modelo anterior) \\ \hline
\bottomrule
\end{xltabular}
\end{tabladoradoLargo}

\vspace{0.4cm}
\Needspace{25\baselineskip}
\subsection{B. Matriz de validación jurídica}
\nopagebreak

\vspace{0.3cm}
\Needspace{15\baselineskip}
\begin{tabladoradoLargo}
\renewcommand{\tabularxcolumn}[1]{p{#1}}
\begin{xltabular}{\textwidth}{|X|X|X|X|}
\caption{Matriz de Datos} \\
\toprule
\rowcolor{gobmxDorado} \encabezadotablaDorada{Hallazgo o limitación} & \encabezadotablaDorada{Fundamento jurídico (APA + cita textual)} & \encabezadotablaDorada{Riesgo (jurídico u operativo)} & \encabezadotablaDorada{Ajuste propuesto (alineado al marco vigente)} \\
\midrule
\endhead
Obsolescencia del modelo \textquotedbl{}Autodeclarativo\textquotedbl{} (DECLARACEL) & RES/174/2016 basaba el cumplimiento en la carga manual de datos por el usuario. & Operativo: Alta incidencia de errores humanos y discrepancias con datos de medición de CENACE. & Sustituir la \textquotedbl{}Declaración\textquotedbl{} por \textbf{\textquotedbl{}Liquidación Administrativa Automática\textquotedbl{}} calculada por la autoridad (Art. 182 R-LSE). \\ \hline
Falta de certeza en tiempos de cierre & El proceso antiguo dependía de los tiempos de captura del usuario (Marzo). & Jurídico: Incertidumbre regulatoria sobre el momento exacto del  cumplimiento. & Establecer \textbf{Fecha de Corte Automático} (ej. 15 Febrero) para la ejecución del algoritmo de liquidación. \\ \hline
Ineficiencia en Reliquidaciones & Correcciones manuales \textquotedbl{}caso por caso\textquotedbl{} ante inconsistencias. & Administrativo: Carga burocrática excesiva para la CNE y participantes. & Implementar \textbf{\textquotedbl{}Ajustes de Reconciliación\textquotedbl{}} automáticos vinculados a las correcciones de medición del Código de Red. \\ \hline
Opacidad del proceso & El usuario solo veía el resultado final \textquotedbl{}aprobado/rechazado\textquotedbl{}. & Transparencia: Incumplimiento del principio de certidumbre. & Implementar el \textbf{\textquotedbl{}Módulo de Cumplimiento en Tiempo Real\textquotedbl{}} (Dashboard) visible todo el año. \\ \hline
\bottomrule
\end{xltabular}
\end{tabladoradoLargo}

\subsection{C. Desarrollo analítico}
\nopagebreak

\subsubsection{Diagnóstico de la situación actual (El legado de DECLARACEL)}
\nopagebreak

El sistema \textbf{DECLARACEL} fue diseñado bajo las Disposiciones Administrativas de 2016 (RES/174/2016) como un formulario digital donde los usuarios \textquotedbl{}confesaban\textquotedbl{} sus obligaciones. Este enfoque generó un ecosistema de:


\begin{enumerate}
\item \textbf{Duplicidad:} El CENACE ya tenía los datos de consumo; pedir al usuario capturarlos nuevamente era redundante.
\item \textbf{Error Humano:} Discrepancias constantes entre lo \textquotedbl{}declarado\textquotedbl{} y lo \textquotedbl{}medido\textquotedbl{} por CENACE.
\item \textbf{Reactividad:} Las fallas se detectaban hasta la auditoría anual, meses después de ocurrido el incumplimiento.

\end{enumerate}
\subsubsection{Estado objetivo (LSE 2025)}
\nopagebreak

El nuevo Reglamento LSE (Art. 182) faculta a la CNE para regular la \textit{\textquotedbl{}liquidación\textquotedbl{}}, entendida como un acto administrativo de cálculo y cierre. El objetivo es un \textbf{Sistema de Cero Captura}:


\begin{enumerate}
\item \textbf{Dato Único:} La CNE toma los consumos directamente de las bases de datos de liquidación del MEM (CENACE).
\item \textbf{Cálculo Oficial:} El sistema aplica los Requisitos de CEL vigentes y determina la obligación.
\item \textbf{Validación:} El usuario solo interviene si discrepa del cálculo oficial (principio de \textquotedbl{}afirmativa ficta\textquotedbl{} de datos).

\end{enumerate}
\vspace{0.4cm}
\Needspace{25\baselineskip}
\subsubsection{Tabla comparativa: modelo anterior vs modelo objetivo}
\nopagebreak

\vspace{0.3cm}
\Needspace{15\baselineskip}
\begin{tablaParadigmaV2}{Análisis Comparativo}{Modelo Anterior (DECLARACEL)}{Modelo Objetivo (Sistema de Liquidación)}
\tpRow{Acción Principal}{Captura Manual de Datos}{Acción Principal}{\textbf{Validación de Cálculo Automático}}
\tpRow{Fuente de Datos}{Declaración del Usuario}{Fuente de Datos}{\textbf{Interoperabilidad CENACE-CNE}}
\tpRow{Tiempo de Proceso}{Anual (Marzo)}{Tiempo de Proceso}{\textbf{Continuo (Cierre Mensual + Anual)}}
\tpRow{Resolución de Discrepancias}{Auditoría Ex-post}{Resolución de Discrepancias}{\textbf{Conciliación en Tiempo Real}}
\tpRow{Base Jurídica}{DACG CRE (Autodeclaración)}{Base Jurídica}{\textbf{R-LSE Art. 182 (Acto Administrativo)}}
\end{tablaParadigmaV2}

\subsubsection{Arquitectura del sistema}
\nopagebreak

\begin{enumerate}
\item \textbf{Ingesta:} El S-CEL recibe mensualmente de CENACE los archivos de \textquotedbl{}Liquidación Final de Consumos\textquotedbl{} de todos los Centros de Carga.
\item \textbf{Procesamiento:} El motor de cálculo aplica el \% de Requisito CEL y genera el \textquotedbl{}Estado de Cuenta de Obligaciones\textquotedbl{}.
\item \textbf{Cobertura:} El sistema verifica el saldo de CELs en la cuenta del Participante.
\item \textbf{Cierre Anual:} En la Fecha de Corte, el sistema \textquotedbl{}cobra\textquotedbl{} (cancela) automáticamente los CELs necesarios y emite el Certificado de Cumplimiento o la Nota de Incumplimiento.

\end{enumerate}
\subsection{D. Propuestas de ajuste normativo}
\nopagebreak

\subsubsection{Instrumento A: Disposiciones de Liquidación y Cumplimiento de CEL (Nueva DACG CNE)}
\nopagebreak

\textbf{Disposición 48.A (Liquidación Automática):}

\textquotedbl{}Se sustituye el mecanismo de declaración manual por el \textbf{Sistema de Liquidación Automática}. La CNE determinará las obligaciones de los Participantes basándose exclusivamente en la información de medición y facturación validada por el CENACE.\textquotedbl{}


\textbf{Disposición 53.A (Calendario de Liquidación):}

\textquotedbl{}El cierre anual de obligaciones se ejecutará automáticamente el \textbf{tercer viernes de Febrero} del año siguiente al periodo de cumplimiento. Los Participantes dispondrán del 1 al 31 de Enero para presentar solicitudes de reconciliación de datos.\textquotedbl{}


\textbf{Disposición Transitoria Tercera (Migración):}

\textquotedbl{}Los expedientes históricos contenidos en la plataforma \textbf{DECLARACEL} serán migrados al nuevo sistema como archivo muerto para fines de trazabilidad y auditoría. A partir de la entrada en vigor, se deshabilita la función de captura manual.\textquotedbl{}


\subsubsection{Instrumento B: Manual del Módulo de Cumplimiento}
\nopagebreak

\textbf{Sección 8.1 (Monitor de Cumplimiento):}

\textquotedbl{}El S-CEL dispondrá de un tablero de control donde el Participante podrá visualizar mes a mes su 'Obligación Acumulada' vs 'Saldo de CELs', con proyecciones de déficit o superávit al cierre de año.\textquotedbl{}


\subsection{E. Observaciones técnicas relevantes}
\nopagebreak

\subsubsection{Riesgos identificados}
\nopagebreak

\begin{itemize}
\item \textbf{Calidad del Dato CENACE:} Si el CENACE tiene retrasos o errores en sus liquidaciones finales, el cálculo de CELs arrastrará el error. Se requiere un protocolo de \textquotedbl{}Re-liquidación\textquotedbl{} (Art. 52 anterior modernizado).
\item \textbf{Resistencia al Cambio:} Usuarios acostumbrados a \textquotedbl{}ajustar\textquotedbl{} sus declaraciones manuales perderán esa flexibilidad indebida.


\end{itemize}
\textbf{Revisión técnica:} Área de Certificados de Energía Limpia

\textbf{Fecha de conclusión:} 12 de enero de 2026


\phantomsection
\addcontentsline{toc}{section}{Régimen de Sanciones Diferenciado y Justicia Energética}
\section*{Régimen de Sanciones Diferenciado y Justicia Energética}
\nopagebreak

\textbf{Alcance del punto:} Definición del nuevo esquema sancionador por incumplimiento de obligaciones de CEL, transitando de un modelo puramente recaudatorio a uno correctivo fundamentado en la Justicia Energética, bajo la autoridad de la CNE.

\textbf{Instrumentos jurídicos revisados:} Ley del Sector Eléctrico 2025 (Art. 165), Reglamento LSE (Art. 33 - Justicia Energética), RES/248/2016 (Marco Derogado).


\vspace{0.4cm}
\Needspace{25\baselineskip}
\subsection{A. Fuentes de información del S-CEL para el Punto 2}
\nopagebreak

\vspace{0.3cm}
\Needspace{15\baselineskip}
\begin{tabladoradoLargo}
\renewcommand{\tabularxcolumn}[1]{p{#1}}
\begin{xltabular}{\textwidth}{|X|X|X|X|}
\caption{Matriz de Datos} \\
\toprule
\rowcolor{gobmxDorado} \encabezadotablaDorada{Actor / Fuente} & \encabezadotablaDorada{Instrumento legal} & \encabezadotablaDorada{Artículo / Numeral} & \encabezadotablaDorada{Cita textual} \\
\midrule
\endhead
\textbf{Facultad Sancionadora} & \textbf{LSE 2025} & \textbf{Art. 165} & \textbf{\textquotedbl{}Las infracciones a esta Ley serán sancionadas por la Comisión Nacional de Energía (CNE)...\textquotedbl{}} \\ \hline
\textbf{Criterio de Multa} & \textbf{LSE 2025} & \textbf{Art. 165 Fracc. IV} & \textbf{\textquotedbl{}Multa equivalente al 0.1\% hasta el 10\% de los ingresos... o monto proporcional al daño causado al SEN.\textquotedbl{}} \\ \hline
\textbf{Enfoque Correctivo} & \textbf{Reglamento LSE} & \textbf{Art. 33} & \textbf{\textquotedbl{}La Justicia Energética debe ser considerada... asegurando que las medidas correctivas prioricen la restitución del equilibrio del sistema.\textquotedbl{}} \\ \hline
\textbf{Reincidencia} & \textbf{LSE 2025} & \textbf{Art. 169} & \textbf{\textquotedbl{}En caso de reincidencia, la multa podrá duplicarse.\textquotedbl{}} \\ \hline
\textit{Antecedente Derogado} & \textit{RES/248/2016} & \textit{Consid. 14} & \textit{\textquotedbl{}Multa de seis a cincuenta salarios mínimos por cada MWh de incumplimiento\textquotedbl{}} (Referencia obsoleta de la CRE). \\ \hline
\bottomrule
\end{xltabular}
\end{tabladoradoLargo}

\vspace{0.4cm}
\Needspace{25\baselineskip}
\subsection{B. Matriz de validación jurídica}
\nopagebreak

\vspace{0.3cm}
\Needspace{15\baselineskip}
\begin{tabladoradoLargo}
\renewcommand{\tabularxcolumn}[1]{p{#1}}
\begin{xltabular}{\textwidth}{|X|X|X|X|}
\caption{Matriz de Datos} \\
\toprule
\rowcolor{gobmxDorado} \encabezadotablaDorada{Hallazgo o limitación} & \encabezadotablaDorada{Fundamento jurídico (APA + cita textual)} & \encabezadotablaDorada{Riesgo (jurídico u operativo)} & \encabezadotablaDorada{Ajuste propuesto (alineado al marco vigente)} \\
\midrule
\endhead
Referencia a autoridad extinta (CRE) & El borrador citaba atribuciones de la CRE bajo la LIE anterior. & Jurídico: Nulidad de cualquier acto sancionador emitido por autoridad incompetente. & Actualizar atribución exclusiva a la \textbf{Comisión Nacional de Energía (CNE)}. \\ \hline
Modelo punitivo vs correctivo & El modelo anterior solo buscaba cobrar (recaudatorio). La LSE busca \textbf{Soberanía y Balance}. & Estratégico: Cobrar multas no genera energía limpia. El sistema fallaba en su objetivo final. & Implementar el \textbf{\textquotedbl{}Plan de Regularización Forzosa\textquotedbl{}} y la \textbf{\textquotedbl{}Inversión Sustituta\textquotedbl{}} como alternativas al pago en efectivo. \\ \hline
Discrecionalidad en el monto & El rango 0.1\% - 10\% es muy amplio sin criterios claros. & Justicia: Riesgo de sanciones desproporcionadas o impugnables. & Establecer una \textbf{Matriz de Graduación} basada en el \% de déficit de CELs y la capacidad del infractor. \\ \hline
Falta de incentivos & El sistema castigaba el error pero no premiaba la corrección temprana. & Operativo: Poca motivación para la autocorrección. & Introducir \textbf{Atenuantes por Regularización Espontánea} antes del inicio del procedimiento. \\ \hline
\bottomrule
\end{xltabular}
\end{tabladoradoLargo}

\subsection{C. Desarrollo analítico}
\nopagebreak

\subsubsection{Diagnóstico de la situación actual}
\nopagebreak

El régimen heredado (RES/248/2016) tenía un defecto estructural: \textbf{La multa era más barata que el cumplimiento}.

Muchos participantes preferían pagar la sanción mínima (o litigarla años en tribunales) en lugar de comprar CELs o invertir en tecnología. Esto convirtió al sistema de sanciones en un \textquotedbl{}costo operativo\textquotedbl{} más, trivializando la obligación de transición energética.


\subsubsection{Estado objetivo (LSE 2025 - Justicia Energética)}
\nopagebreak

El nuevo marco (Art. 33 R-LSE) introduce el concepto de \textbf{Justicia Energética Sancionadora}:


\begin{enumerate}
\item \textbf{Objetivo:} No es recaudar dinero, es asegurar que la energía limpia se genere.
\item \textbf{Mecanismo: \textbf{ Si un participante incumple, la prioridad no es multarlo, es obligarlo a financiar esa energía limpia faltante (}Inversión Sustituta}).
\item \textbf{Severidad:} Para quienes reincidan o actúen con dolo, la sanción escala hasta la posible revocación del Permiso de Suministro.

\end{enumerate}
\vspace{0.4cm}
\Needspace{25\baselineskip}
\subsubsection{Tabla comparativa: modelo anterior vs modelo objetivo}
\nopagebreak

\vspace{0.3cm}
\Needspace{15\baselineskip}
\begin{tablaParadigmaV2}{Análisis Comparativo}{Modelo Anterior (Punitivo)}{Modelo Objetivo (Correctivo/CNE)}
\tpRow{Autoridad}{CRE (Regulador Económico)}{Autoridad}{\textbf{CNE (Autoridad de Estado)}}
\tpRow{Fin de la Sanción}{Recaudación Fiscal}{Fin de la Sanción}{\textbf{Restitución del Balance Energético}}
\tpRow{Destino del Recurso}{Tesorería / Fondo Universal}{Destino del Recurso}{\textbf{FOTEASE (Etiquetado para Proyectos)}}
\tpRow{Alternativa a Multa}{Ninguna (Solo pago)}{Alternativa a Multa}{\textbf{Inversión Sustituta en Proyectos Estratégicos}}
\tpRow{Base de Cálculo}{Salarios Mínimos por MWh}{Base de Cálculo}{\textbf{\% de Ingresos + Daño al Sistema}}
\end{tablaParadigmaV2}

\subsubsection{Arquitectura del sistema sancionador}
\nopagebreak

\begin{enumerate}
\item \textbf{Detección:} El \textit{Sistema de Liquidación Automática} (ver Punto 1) emite la \textquotedbl{}Alerta de Déficit\textquotedbl{}.
\item \textbf{Pre-Sanción: \textbf{ Se otorga una ventana de 30 días para presentar un }Plan de Regularización}.
\item \textbf{Procedimiento:} Si no hay plan, la CNE inicia el Procedimiento Administrativo Sancionador (PAS).
\item \textbf{Resolución:}
\end{enumerate}
\begin{itemize}
\item \textbf{Opción A (Pago):} Multa económica al FOTEASE.
\item \textbf{Opción B (Inversión):} El infractor financia directamente un proyecto de generación limpia certificado por CNE (Inversión Sustituta).

\end{itemize}
\subsection{D. Propuestas de ajuste normativo}
\nopagebreak

\subsubsection{Instrumento A: Disposiciones de Disciplina de Mercado y Sanciones (CNE)}
\nopagebreak

\textbf{Disposición 60 (Matriz de Graduación):}

\textquotedbl{}Para determinar el monto de la multa dentro del rango legal (0.1\% al 10\%), la CNE aplicará la siguiente matriz:


\begin{itemize}
\item \textbf{Leve (0.1\% - 3\%):} Déficit \textless{} 5\% de la obligación, primer incumplimiento, sin dolo.
\item \textbf{Grave (3\% - 7\%):} Déficit \textgreater{} 5\%, reincidencia, o reporte de información falsa.
\item \textbf{Crítico (7\% - 10\%):} Incumplimiento sistemático, contumacia, o riesgo a la confiabilidad.\textquotedbl{}

\end{itemize}
\textbf{Disposición 65 (Inversión Sustituta):}

\textquotedbl{}Atenuante de Inversión. El infractor podrá solicitar la conmutación de hasta el 80\% de la multa económica a cambio de realizar una inversión equivalente en \textbf{Proyectos de Generación Limpia en Zonas Prioritarias} definidas en el PROSENER, bajo supervisión de la CNE.\textquotedbl{}


\subsubsection{Instrumento B: Procedimiento de Justicia Energética}
\nopagebreak

\textbf{Regla 12 (Debido Proceso):}

\textquotedbl{}Se garantiza el derecho de audiencia. Antes de la imposición de multa, el participante será citado a una \textbf{Junta de Conciliación Técnica} para explorar alternativas de cumplimiento forzoso que beneficien al SEN.\textquotedbl{}


\subsection{E. Observaciones técnicas relevantes}
\nopagebreak

\subsubsection{Riesgos identificados}
\nopagebreak

\begin{itemize}
\item \textbf{Riesgo de Litigio:} Los grandes consumidores podrían ampararse contra las multas altas (10\% de ingresos).
\item \textit{Mitigación:} La opción de \textquotedbl{}Inversión Sustituta\textquotedbl{} reduce el incentivo a litigar, pues el dinero se queda en activos productivos del sector (o del propio grupo) en lugar de irse a \textquotedbl{}fondo perdido\textquotedbl{}.
\item \textbf{Capacidad de Supervisión:} La CNE necesita un equipo robusto para validar que las \textquotedbl{}Inversiones Sustitutas\textquotedbl{} realmente se ejecuten.


\end{itemize}
\textbf{Revisión técnica:} Área de Certificados de Energía Limpia

\textbf{Fecha de conclusión:} 12 de enero de 2026


\phantomsection
\addcontentsline{toc}{section}{Transparencia, Información Pública y Reportes del Sistema de CEL}
\section*{Transparencia, Información Pública y Reportes del Sistema de CEL}
\nopagebreak

\textbf{Alcance del punto:} Definición del esquema de Transparencia Administrativa y Soberanía de Datos para el Sistema de Certificados de Energía Limpia (S-CEL), alineado con la rectoría del Estado.

\textbf{Instrumentos jurídicos revisados:} Reglamento de la Ley del Sector Eléctrico 2025 (R-LSE), Ley del Sector Eléctrico 2025 (LSE), Resolución RES/174/2016 (Marco Histórico).


\vspace{0.4cm}
\Needspace{25\baselineskip}
\subsection{A. Fuentes de información del S-CEL para el Punto 3}
\nopagebreak

\vspace{0.3cm}
\Needspace{15\baselineskip}
\begin{tabladoradoLargo}
\renewcommand{\tabularxcolumn}[1]{p{#1}}
\begin{xltabular}{\textwidth}{|X|X|X|X|}
\caption{Matriz de Datos} \\
\toprule
\rowcolor{gobmxDorado} \encabezadotablaDorada{Actor / Fuente} & \encabezadotablaDorada{Instrumento legal} & \encabezadotablaDorada{Artículo / Numeral} & \encabezadotablaDorada{Cita textual / Interpretación} \\
\midrule
\endhead
\textbf{Rectoría del Sistema} & \textbf{Reglamento LSE 2025} & \textbf{Artículo 182} & \textquotedbl{}La CNE, en coordinación con la Secretaría, debe emitir las disposiciones administrativas... que regulen el otorgamiento, liquidación... las cuales deben incluir el funcionamiento del sistema electrónico... Este sistema constituye el registro de certificados.\textquotedbl{} \\ \hline
\textbf{Principio de Transparencia} & \textbf{Reglamento LSE 2025} & \textbf{Artículo 6} & \textquotedbl{}La actuación administrativa... se debe desarrollar conforme a los principios de... transparencia, imparcialidad... y en su caso, considerar la Justicia Energética.\textquotedbl{} \\ \hline
\textbf{Publicidad de Información} & \textbf{Reglamento LSE 2025} & \textbf{Artículo 11} & Obligación de publicar informes pormenorizados sobre proyectos y estado del sistema (aplicación análoga para CEL). \\ \hline
\textbf{Antecedente (Limitado)} & \textbf{RES/174/2016} & \textbf{Disposición 60} & \textquotedbl{}Mensualmente, la Comisión publicará... información estadística que refleje de manera agregada, la cantidad de CEL vigentes... precios... transacciones.\textquotedbl{} \\ \hline
\textbf{Protección de Datos} & \textbf{LFTAIP / R-LSE} & \textbf{Art. 113 LFTAIP} & Protección de secretos comerciales e industriales, armonizado con la obligación de transparencia pública. \\ \hline
\bottomrule
\end{xltabular}
\end{tabladoradoLargo}

\vspace{0.4cm}
\Needspace{25\baselineskip}
\subsection{B. Matriz de validación jurídica}
\nopagebreak

\vspace{0.3cm}
\Needspace{15\baselineskip}
\begin{tabladoradoLargo}
\renewcommand{\tabularxcolumn}[1]{p{#1}}
\begin{xltabular}{\textwidth}{|X|X|X|X|}
\caption{Matriz de Datos} \\
\toprule
\rowcolor{gobmxDorado} \encabezadotablaDorada{Hallazgo o limitación (Marco Anterior - Mercantil)} & \encabezadotablaDorada{Fundamento jurídico (Corrección LSE 2025)} & \encabezadotablaDorada{Riesgo (Operativo / Soberanía)} & \encabezadotablaDorada{Ajuste propuesto (Nuevo Modelo CNE)} \\
\midrule
\endhead
\textbf{Opacidad por Agregación Excesiva:} La RES/174/2016 (Disp. 60) permitía una agregación tal que impedía validar el cumplimiento real de metas nacionales. & \textbf{Art. 182 R-LSE:} Faculta a CNE para regular el sistema sin restricciones arcaicas de agregación, priorizando la fiscalización pública. & \textbf{Riesgo de \textquotedbl{}Greenwashing\textquotedbl{}:} Imposibilidad de auditar si los CEL provienen de proyectos reales o especulativos. & \textbf{Publicidad Desagregada por Tecnologías Estratégicas:} Reportes públicos por zona y tecnología (sin revelar secretos comerciales unitarios). \\ \hline
\textbf{Enfoque en Especulación de Precios:} El marco anterior priorizaba \textquotedbl{}señales de precios\textquotedbl{} para traders financieros. & \textbf{Art. 6 R-LSE (Justicia Energética):} El objetivo es la \textquotedbl{}Justicia Energética\textquotedbl{}, no la rentabilidad financiera del certificado. & \textbf{Volatilidad Artificial:} Fomentar la especulación con el precio del CEL desvirtúa su fin ambiental. & \textbf{Dashboard de Cumplimiento Normativo:} Enfocar la transparencia en el avance de metas nacionales y cumplimiento de obligados, no en ticks de precios. \\ \hline
\textbf{Falta de Certeza en la Cancelación:} No existía un mecanismo público claro para verificar que un CEL \textquotedbl{}cancelado\textquotedbl{} realmente saliera del mercado. & \textbf{Art. 182 R-LSE:} \textquotedbl{}Liquidación, cancelación voluntaria...\textquotedbl{}. & \textbf{Doble Contabilidad:} Riesgo de reclamo de atributos limpios duplicados internacionalmente. & \textbf{Libro Mayor Público de Retiros:} Publicación de listas de series de CEL retiradas definitivamente (Trazabilidad Total). \\ \hline
\textbf{Dependencia de Plataformas Privadas:} Se delegaba la inteligencia de mercado a consultoras privadas por falta de datos oficiales claros. & \textbf{Soberanía de Datos (Política CNE):} El Estado debe ser la fuente única de la verdad. & \textbf{Asimetría de Información:} Actores grandes con mejores datos que pymes o la sociedad civil. & \textbf{Portal de Datos Abiertos CNE:} Acceso gratuito, directo y oficial a la estadística del S-CEL. \\ \hline
\bottomrule
\end{xltabular}
\end{tabladoradoLargo}

\subsection{C. Desarrollo analítico: De la Especulación a la Soberanía de Datos}
\nopagebreak

\subsubsection{1. Diagnóstico del Modelo Fallido}
\nopagebreak

El modelo de transparencia de la RES/174/2016 estaba diseñado para un \textquotedbl{}mercado financiero\textquotedbl{}. Se buscaba replicar esquemas bursátiles donde la variable crítica es el precio y la volatilidad. Esto resultó en:


\begin{itemize}
\item \textbf{Ruido en la señal:} Precios distorsionados por multas (sanciones) y falta de liquidez real.
\item \textbf{Cajas negras:} Agregación de datos que ocultaba el desempeño real de tecnologías específicas.
\item \textbf{Privatización de la información:} Bancos y consultoras vendían reportes que debían ser públicos.

\end{itemize}
\subsubsection{2. Nuevo Modelo: Transparencia Administrativa y Certeza}
\nopagebreak

Bajo la LSE 2025, el CEL es un instrumento de planeación y cumplimiento, no un activo financiero especulativo. La transparencia cambia de objetivo: de \textquotedbl{}descubrir precios\textquotedbl{} a \textbf{\textquotedbl{}Garantizar la Fe Pública del Atributo Limpio\textquotedbl{}}.


\textbf{Pilares del Nuevo S-CEL:}


\begin{enumerate}
\item \textbf{Fuente Única de Verdad (CNE):} Nadie debe pagar por saber cuánto se ha generado.
\item \textbf{Trazabilidad Social:} La sociedad debe saber qué empresas cumplen y cuáles simulan.
\item \textbf{Protección Estratégica:} Se protegen los costos unitarios (secretos industriales) pero se publican los volúmenes agregados (interés público).

\end{enumerate}
\vspace{0.4cm}
\Needspace{25\baselineskip}
\subsubsection{3. Tabla Comparativa: Transición de Modelo}
\nopagebreak

\vspace{0.3cm}
\Needspace{15\baselineskip}
\begin{tablaParadigmaV2}{Análisis Comparativo}{Modelo Neoliberal (RES/174/2016)}{Modelo Nacional (LSE 2025)}
\tpRow{\textbf{Objetivo de los Datos}}{Facilitar trading y especulación financiera.}{\textbf{Objetivo de los Datos}}{Fiscalización pública y planeación energética.}
\tpRow{\textbf{Métrica Principal}}{Precio Spot / Futuros.}{\textbf{Métrica Principal}}{\% de Cumplimiento de Metas Nacionales / Balance Energético.}
\tpRow{\textbf{Acceso a Datos}}{PDF estáticos mensuales (limitados).}{\textbf{Acceso a Datos}}{Portal de Datos Abiertos y Visor Geográfico Público.}
\tpRow{\textbf{Granularidad}}{\textquotedbl{}Caja Negra\textquotedbl{} (Nacional Agregado).}{\textbf{Granularidad}}{Desglose por Zona de Gerencia de Control y Tecnología.}
\tpRow{\textbf{Costo de Información}}{Gratuita (básica) / Consultoras (detallada).}{\textbf{Costo de Información}}{100\% Gratuita (Principio de Gobierno Abierto).}
\end{tablaParadigmaV2}

\subsection{D. Propuestas de ajuste normativo}
\nopagebreak

\subsubsection{1. Creación del Portal de Transparencia del S-CEL (CNE-DATA)}
\nopagebreak

Se propone instruir en las nuevas DACG la creación de un módulo público en el sitio web de la CNE:


\textbf{Contenido Obligatorio:}


\begin{itemize}
\item \textbf{Semáforo de Metas Nacionales:} \% de avance anual vs meta PROSENER.
\item \textbf{Inventario de Emisiones:} Total de CEL emitidos por mes, desglosados por tecnología (Solar, Eólica, Hidro, Nuclear, Bio).
\item \textbf{Registro de Retiros:} Listado de CEL liquidados (para cumplimiento) y Cancelados (voluntarios), garantizando que \textquotedbl{}1 MWh limpio = 1 CEL usado\textquotedbl{}.
\item \textbf{Indicadores de Precio Promedio Ponderado:} Calculado administrativamente para fines de planeación (no tiempo real).

\end{itemize}
\subsubsection{2. Reportes Ejecutivos para Planeación}
\nopagebreak

La CNE remitirá a la SENER reportes trimestrales para alimentar el PROSENER:


\begin{itemize}
\item Proyección de Oferta de CEL a 3 años.
\item Alertas de déficit regional.
\item Análisis de impacto de la \textquotedbl{}Inversión Sustituta\textquotedbl{} (ver Bloque VII Punto 2).

\end{itemize}
\subsubsection{3. Mecanismo de Protección de Datos}
\nopagebreak

Se establece una regla clara de anonimización:


\begin{itemize}
\item \textbf{Público:} Volúmenes totales por tecnología y zona. Lista de empresas cumplidas (Sí/No).
\item \textbf{Reservado:} Precios de contratos bilaterales específicos, datos técnicos de vulnerabilidad de red.

\end{itemize}
\subsection{E. Estrategia de Transición (Transitorios)}
\nopagebreak

\subsubsection{Migración de Históricos}
\nopagebreak

\textbf{Artículo Transitorio Primero.} A la entrada en vigor del nuevo S-CEL, la CNE publicará un \textquotedbl{}Libro Blanco 2018-2024\textquotedbl{} que consolide y depure la información histórica dispersa de los años anteriores, cerrando definitivamente los ciclos contables del modelo anterior.


\textbf{Artículo Transitorio Segundo.} Se suspende la obligación de reportar \textquotedbl{}precios spot\textquotedbl{} diarios en plataformas externas. Los participantes únicamente reportarán el valor de la contraprestación al momento del registro del contrato bilateral en el S-CEL, para fines estadísticos internos.


\subsubsection{Implementación Tecnológica}
\nopagebreak

\textbf{Fase 1 (Día 0):} Publicación de PDFs estructurados con el nuevo desglose.

\textbf{Fase 2 (Mes 6):} Lanzamiento del Visor de Datos Abiertos (CSV/API de consulta básica).



\textbf{Conclusión:} La transparencia en el nuevo modelo no busca \textquotedbl{}señales de mercado\textquotedbl{} para especuladores, sino \textquotedbl{}señales de certeza\textquotedbl{} para la inversión productiva y la confianza ciudadana en la transición energética.


\newpage
\contraportada

\end{document}
